% !TEX program = xelatex
\documentclass[a4paper,11pt,openany,fontset=none]{ctexrep}
\newcommand{\reporttitle}{A股盈利修复的再定价机会}
\newcommand{\reportsubtitle}{从2026H1预告扩散到可执行的行业与个股验证框架}
\newcommand{\reportkicker}{A股策略研究}
\newcommand{\reportscope}{中国A股 | 1,680家预告公司 | 39只优先池 | 4只可复核模型}
\newcommand{\reportdate}{2026年7月15日}
\newcommand{\reportdatacutoff}{官方预告与个股行情截至2026年7月15日；行业资金观察至2026年7月14日}
\newcommand{\reporttype}{A股策略与行业配置}
\newcommand{\reportauthor}{AStock策略研究}
\newcommand{\reporthouseview}{盈利预告改善并不等于全面风险偏好开启。当前更具可执行性的机会集中在分母可信、现金流可验证、价格尚未充分反映且能够获得当前外部证据支持的样本。}
\newcommand{\reportquality}{官方预告覆盖1,680家；高影响候选142只；优先池39只；Q1财务142/142；候选研报PDF 137/142；行业资金日/周表31/31。}
\newcommand{\reportdisclaimer}{本报告基于公开资料整理，不构成任何证券买卖建议。}
% Reusable IB-style preamble for XeLaTeX equity research reports
% Usage: % Reusable IB-style preamble for XeLaTeX equity research reports
% Usage: % Reusable IB-style preamble for XeLaTeX equity research reports
% Usage: \input{preamble} after \documentclass[a4paper,11pt,openany,fontset=none]{ctexrep}

% ============ 页面布局 ============
\usepackage[top=2cm, bottom=2cm, left=2cm, right=2cm]{geometry}
\setlength{\parskip}{0.3em}
\setlength{\parindent}{2em}
\renewcommand{\baselinestretch}{1.15}

% ============ 字体（macOS） ============
\setCJKmainfont[BoldFont={Heiti SC}]{STSong}
\setCJKsansfont{Heiti SC}
\setCJKmonofont{Heiti SC}
\setmainfont{Times New Roman}

% ============ 颜色（IB Navy/Grey palette） ============
\usepackage{xcolor}
\definecolor{navy}{HTML}{003366}
\definecolor{steelgrey}{HTML}{4A5568}
\definecolor{lightgrey}{HTML}{F7FAFC}
\definecolor{riskred}{HTML}{C53030}
\definecolor{riskamber}{HTML}{D69E2E}
\definecolor{riskgreen}{HTML}{38A169}
\definecolor{nvgreen}{HTML}{76B900}

% ============ 表格 ============
\usepackage{booktabs}
\usepackage{longtable}
\usepackage{tabularx}
\usepackage{multirow}
\usepackage{array}
\newcolumntype{Y}{>{\centering\arraybackslash}X}
\newcolumntype{L}[1]{>{\raggedright\arraybackslash}p{#1}}
\newcolumntype{C}[1]{>{\centering\arraybackslash}p{#1}}
\newcolumntype{R}[1]{>{\raggedleft\arraybackslash}p{#1}}

% ============ 图形 ============
\usepackage{tikz}
\usetikzlibrary{shapes,arrows.meta,positioning,calc,backgrounds,fit}
\usepackage{pgfplots}
\usepgfplotslibrary{polar}
\pgfplotsset{compat=1.18}
\usepackage[edges]{forest}

% ============ 页眉页脚（报告标题和日期需在main.tex中覆盖） ============
\usepackage{fancyhdr}
\pagestyle{fancy}
\fancyhf{}
\fancyhead[L]{\small\color{steelgrey} \reporttitle}
\fancyhead[R]{\small\color{steelgrey} \reportdate}
\fancyfoot[C]{\small\color{steelgrey} \thepage}
\fancyfoot[R]{\tiny\color{steelgrey} 本报告不构成任何投资建议}
\renewcommand{\headrulewidth}{0.4pt}
\renewcommand{\footrulewidth}{0.2pt}

% ============ 章节格式 ============
\usepackage{titlesec}
\titleformat{\chapter}[display]
  {\normalfont\huge\bfseries\color{navy}}
  {\chaptertitlename\ \thechapter}{20pt}{\Huge}
\titleformat{\section}
  {\normalfont\Large\bfseries\color{navy}}
  {\thesection}{1em}{}
\titleformat{\subsection}
  {\normalfont\large\bfseries\color{steelgrey}}
  {\thesubsection}{1em}{}

% ============ 超链接 ============
\usepackage{hyperref}
\hypersetup{
  colorlinks=true,
  linkcolor=navy,
  urlcolor=navy,
  citecolor=navy,
}

% ============ 其他工具 ============
\usepackage{enumitem}
\setlist{nosep, leftmargin=2em}
\usepackage{fontawesome5}
\usepackage{amssymb}
\usepackage{pifont}
\usepackage{tcolorbox}
\tcbuselibrary{skins,breakable}
\usepackage{caption}
\captionsetup{font=small, skip=4pt}
\usepackage{float}
\setlength{\textfloatsep}{8pt plus 2pt minus 2pt}
\setlength{\floatsep}{6pt plus 2pt minus 2pt}
\setlength{\intextsep}{6pt plus 2pt minus 2pt}
\setlength{\abovecaptionskip}{4pt}
\setlength{\belowcaptionskip}{2pt}

% ============ 自定义命令 ============
\newcommand{\rmark}{\textcolor{riskred}{\ding{108}}}
\newcommand{\amark}{\textcolor{riskamber}{\ding{108}}}
\newcommand{\gmark}{\textcolor{riskgreen}{\ding{108}}}
\newcommand{\starfive}{$\bigstar\bigstar\bigstar\bigstar\bigstar$}
\newcommand{\starfour}{$\bigstar\bigstar\bigstar\bigstar$}
\newcommand{\starthree}{$\bigstar\bigstar\bigstar$}

% 信息框
\newtcolorbox{keyinsight}[1][]{
  colback=lightgrey, colframe=navy,
  fonttitle=\bfseries, title={#1}, breakable
}
\newtcolorbox{riskbox}[1][]{
  colback=red!5, colframe=riskred,
  fonttitle=\bfseries, title={#1}, breakable
}
 after \documentclass[a4paper,11pt,openany,fontset=none]{ctexrep}

% ============ 页面布局 ============
\usepackage[top=2cm, bottom=2cm, left=2cm, right=2cm]{geometry}
\setlength{\parskip}{0.3em}
\setlength{\parindent}{2em}
\renewcommand{\baselinestretch}{1.15}

% ============ 字体（macOS） ============
\setCJKmainfont[BoldFont={Heiti SC}]{STSong}
\setCJKsansfont{Heiti SC}
\setCJKmonofont{Heiti SC}
\setmainfont{Times New Roman}

% ============ 颜色（IB Navy/Grey palette） ============
\usepackage{xcolor}
\definecolor{navy}{HTML}{003366}
\definecolor{steelgrey}{HTML}{4A5568}
\definecolor{lightgrey}{HTML}{F7FAFC}
\definecolor{riskred}{HTML}{C53030}
\definecolor{riskamber}{HTML}{D69E2E}
\definecolor{riskgreen}{HTML}{38A169}
\definecolor{nvgreen}{HTML}{76B900}

% ============ 表格 ============
\usepackage{booktabs}
\usepackage{longtable}
\usepackage{tabularx}
\usepackage{multirow}
\usepackage{array}
\newcolumntype{Y}{>{\centering\arraybackslash}X}
\newcolumntype{L}[1]{>{\raggedright\arraybackslash}p{#1}}
\newcolumntype{C}[1]{>{\centering\arraybackslash}p{#1}}
\newcolumntype{R}[1]{>{\raggedleft\arraybackslash}p{#1}}

% ============ 图形 ============
\usepackage{tikz}
\usetikzlibrary{shapes,arrows.meta,positioning,calc,backgrounds,fit}
\usepackage{pgfplots}
\usepgfplotslibrary{polar}
\pgfplotsset{compat=1.18}
\usepackage[edges]{forest}

% ============ 页眉页脚（报告标题和日期需在main.tex中覆盖） ============
\usepackage{fancyhdr}
\pagestyle{fancy}
\fancyhf{}
\fancyhead[L]{\small\color{steelgrey} \reporttitle}
\fancyhead[R]{\small\color{steelgrey} \reportdate}
\fancyfoot[C]{\small\color{steelgrey} \thepage}
\fancyfoot[R]{\tiny\color{steelgrey} 本报告不构成任何投资建议}
\renewcommand{\headrulewidth}{0.4pt}
\renewcommand{\footrulewidth}{0.2pt}

% ============ 章节格式 ============
\usepackage{titlesec}
\titleformat{\chapter}[display]
  {\normalfont\huge\bfseries\color{navy}}
  {\chaptertitlename\ \thechapter}{20pt}{\Huge}
\titleformat{\section}
  {\normalfont\Large\bfseries\color{navy}}
  {\thesection}{1em}{}
\titleformat{\subsection}
  {\normalfont\large\bfseries\color{steelgrey}}
  {\thesubsection}{1em}{}

% ============ 超链接 ============
\usepackage{hyperref}
\hypersetup{
  colorlinks=true,
  linkcolor=navy,
  urlcolor=navy,
  citecolor=navy,
}

% ============ 其他工具 ============
\usepackage{enumitem}
\setlist{nosep, leftmargin=2em}
\usepackage{fontawesome5}
\usepackage{amssymb}
\usepackage{pifont}
\usepackage{tcolorbox}
\tcbuselibrary{skins,breakable}
\usepackage{caption}
\captionsetup{font=small, skip=4pt}
\usepackage{float}
\setlength{\textfloatsep}{8pt plus 2pt minus 2pt}
\setlength{\floatsep}{6pt plus 2pt minus 2pt}
\setlength{\intextsep}{6pt plus 2pt minus 2pt}
\setlength{\abovecaptionskip}{4pt}
\setlength{\belowcaptionskip}{2pt}

% ============ 自定义命令 ============
\newcommand{\rmark}{\textcolor{riskred}{\ding{108}}}
\newcommand{\amark}{\textcolor{riskamber}{\ding{108}}}
\newcommand{\gmark}{\textcolor{riskgreen}{\ding{108}}}
\newcommand{\starfive}{$\bigstar\bigstar\bigstar\bigstar\bigstar$}
\newcommand{\starfour}{$\bigstar\bigstar\bigstar\bigstar$}
\newcommand{\starthree}{$\bigstar\bigstar\bigstar$}

% 信息框
\newtcolorbox{keyinsight}[1][]{
  colback=lightgrey, colframe=navy,
  fonttitle=\bfseries, title={#1}, breakable
}
\newtcolorbox{riskbox}[1][]{
  colback=red!5, colframe=riskred,
  fonttitle=\bfseries, title={#1}, breakable
}
 after \documentclass[a4paper,11pt,openany,fontset=none]{ctexrep}

% ============ 页面布局 ============
\usepackage[top=2cm, bottom=2cm, left=2cm, right=2cm]{geometry}
\setlength{\parskip}{0.3em}
\setlength{\parindent}{2em}
\renewcommand{\baselinestretch}{1.15}

% ============ 字体（macOS） ============
\setCJKmainfont[BoldFont={Heiti SC}]{STSong}
\setCJKsansfont{Heiti SC}
\setCJKmonofont{Heiti SC}
\setmainfont{Times New Roman}

% ============ 颜色（IB Navy/Grey palette） ============
\usepackage{xcolor}
\definecolor{navy}{HTML}{003366}
\definecolor{steelgrey}{HTML}{4A5568}
\definecolor{lightgrey}{HTML}{F7FAFC}
\definecolor{riskred}{HTML}{C53030}
\definecolor{riskamber}{HTML}{D69E2E}
\definecolor{riskgreen}{HTML}{38A169}
\definecolor{nvgreen}{HTML}{76B900}

% ============ 表格 ============
\usepackage{booktabs}
\usepackage{longtable}
\usepackage{tabularx}
\usepackage{multirow}
\usepackage{array}
\newcolumntype{Y}{>{\centering\arraybackslash}X}
\newcolumntype{L}[1]{>{\raggedright\arraybackslash}p{#1}}
\newcolumntype{C}[1]{>{\centering\arraybackslash}p{#1}}
\newcolumntype{R}[1]{>{\raggedleft\arraybackslash}p{#1}}

% ============ 图形 ============
\usepackage{tikz}
\usetikzlibrary{shapes,arrows.meta,positioning,calc,backgrounds,fit}
\usepackage{pgfplots}
\usepgfplotslibrary{polar}
\pgfplotsset{compat=1.18}
\usepackage[edges]{forest}

% ============ 页眉页脚（报告标题和日期需在main.tex中覆盖） ============
\usepackage{fancyhdr}
\pagestyle{fancy}
\fancyhf{}
\fancyhead[L]{\small\color{steelgrey} \reporttitle}
\fancyhead[R]{\small\color{steelgrey} \reportdate}
\fancyfoot[C]{\small\color{steelgrey} \thepage}
\fancyfoot[R]{\tiny\color{steelgrey} 本报告不构成任何投资建议}
\renewcommand{\headrulewidth}{0.4pt}
\renewcommand{\footrulewidth}{0.2pt}

% ============ 章节格式 ============
\usepackage{titlesec}
\titleformat{\chapter}[display]
  {\normalfont\huge\bfseries\color{navy}}
  {\chaptertitlename\ \thechapter}{20pt}{\Huge}
\titleformat{\section}
  {\normalfont\Large\bfseries\color{navy}}
  {\thesection}{1em}{}
\titleformat{\subsection}
  {\normalfont\large\bfseries\color{steelgrey}}
  {\thesubsection}{1em}{}

% ============ 超链接 ============
\usepackage{hyperref}
\hypersetup{
  colorlinks=true,
  linkcolor=navy,
  urlcolor=navy,
  citecolor=navy,
}

% ============ 其他工具 ============
\usepackage{enumitem}
\setlist{nosep, leftmargin=2em}
\usepackage{fontawesome5}
\usepackage{amssymb}
\usepackage{pifont}
\usepackage{tcolorbox}
\tcbuselibrary{skins,breakable}
\usepackage{caption}
\captionsetup{font=small, skip=4pt}
\usepackage{float}
\setlength{\textfloatsep}{8pt plus 2pt minus 2pt}
\setlength{\floatsep}{6pt plus 2pt minus 2pt}
\setlength{\intextsep}{6pt plus 2pt minus 2pt}
\setlength{\abovecaptionskip}{4pt}
\setlength{\belowcaptionskip}{2pt}

% ============ 自定义命令 ============
\newcommand{\rmark}{\textcolor{riskred}{\ding{108}}}
\newcommand{\amark}{\textcolor{riskamber}{\ding{108}}}
\newcommand{\gmark}{\textcolor{riskgreen}{\ding{108}}}
\newcommand{\starfive}{$\bigstar\bigstar\bigstar\bigstar\bigstar$}
\newcommand{\starfour}{$\bigstar\bigstar\bigstar\bigstar$}
\newcommand{\starthree}{$\bigstar\bigstar\bigstar$}

% 信息框
\newtcolorbox{keyinsight}[1][]{
  colback=lightgrey, colframe=navy,
  fonttitle=\bfseries, title={#1}, breakable
}
\newtcolorbox{riskbox}[1][]{
  colback=red!5, colframe=riskred,
  fonttitle=\bfseries, title={#1}, breakable
}

\sloppy
\setlength{\emergencystretch}{3em}
\hypersetup{pdfauthor={\reportauthor},pdftitle={\reporttitle}}
\usepackage{amsmath}
\begin{document}
\astockcover
\tableofcontents
\clearpage
\chapter{投资结论与行动框架}

\begin{houseviewbox}[核心判断｜盈利预告的广度不等于全面风险偏好]
我们不把1,680家公司的预告扩散解读为全面风险偏好开启。可执行机会集中在三类：第一，已经形成利润—现金流—外部锚闭环的可复核验证模型；第二，利润修复可信、但仍需下一季证据确认的事件验证标的；第三，价格已先行或分母质量不足、只能作为风险监控的标的。因此，本报告的重点不是寻找最高目标价，也不是直接给出配置建议，而是识别在当前价格下可被证据、外部锚和风险预算共同约束的验证队列。
\end{houseviewbox}

\section{投资行动摘要}

\noindent
\begin{exhibitbox}[图表 1｜当前可复核模型与行动含义]
\small
\begin{tabularx}{\textwidth}{L{1.1cm}L{1.7cm}L{1.5cm}R{1.0cm}R{1.0cm}R{1.1cm}X}
\toprule
\textbf{代码} & \textbf{标的} & \textbf{角色} & \textbf{现价} & \textbf{概率值} & \textbf{空间} & \textbf{行动与约束} \\
\midrule
000703 & 恒逸石化 & 可复核正向验证模型 & 14.48 & 23.42 & 61.7\% & 仅作验证模型；等待H2扣非、现金流和外部锚同步确认后再讨论风险暴露 \\
002379 & 宏桥控股 & 可复核正向验证模型 & 18.66 & 25.80 & 38.3\% & 仅作验证模型；等待H2扣非、现金流和外部锚同步确认后再讨论风险暴露 \\
300014 & 亿纬锂能 & 可复核正向验证模型 & 52.28 & 60.43 & 15.6\% & 仅作验证模型；等待H2扣非、现金流和外部锚同步确认后再讨论风险暴露 \\
000155 & 川能动力 & 下行纪律对照 & 11.09 & 10.04 & -9.5\% & 概率值低于现价；不新增风险暴露，用于约束追价和估值纪律 \\
\bottomrule
\end{tabularx}
\normalsize
\sourcenote{官方业绩预告、Q1财务、2026年7月15日收盘价及已归档原始研报；概率值为本机构情景值，非收益承诺。}
\end{exhibitbox}

\noindent
\begin{exhibitbox}[图表 2｜三层行动框架：从可复核价值到风险纪律]
\begin{minipage}[t]{0.31\textwidth}
\textbf{\color{riskgreen}可复核正向验证}
\par\smallskip
当前价格、情景分母、原始研报锚与风险触发器均可复核。
\par\smallskip
\textbf{代表：}恒逸石化、宏桥控股。
\par\smallskip
\textbf{纪律：}不直接配置；H2扣非或现金流偏离即降级。
\end{minipage}\hfill
\begin{minipage}[t]{0.31\textwidth}
\textbf{\color{riskamber}事件验证观察}
\par\smallskip
盈利修复有线索，但外部锚、现金流或持续性仍未闭环。
\par\smallskip
\textbf{代表：}中国船舶、中国铝业、三六零、中远海能。
\par\smallskip
\textbf{纪律：}不以模型空间作为行动理由。
\end{minipage}\hfill
\begin{minipage}[t]{0.31\textwidth}
\textbf{\color{riskred}估值纪律观察}
\par\smallskip
概率值低于现价、周期盈利可能回落或证据时效不足。
\par\smallskip
\textbf{代表：}川能动力及高估值样本。
\par\smallskip
\textbf{纪律：}不追价，作为行业和风险偏好的反向指标。
\end{minipage}
\sourcenote{行动分层由当前价格、估值分母、外部锚和证据质量共同决定，不等同于券商投资评级。}
\end{exhibitbox}

\section{执行纪律：先验证，再放大}

可复核正向模型只在利润、现金流与外部锚三项同时未偏离时维持；任一项失效即转入事件验证。事件验证观察不以模型空间作为建仓理由，须等待下一份财报、订单或价格信号把假设变为事实。估值纪律观察不追价，主要用于识别行业盈利、价格与风险偏好是否发生反向变化。

\section{补证后的候选池调整}

原中等证据标的共48只，补证后基础证据不再是主要阻断项；其中6只进入优先验证候选、11只进入扩展事件验证、3只进入安全边际观察，28只仅保留为风险监控或暂不纳入。没有任何原中等证据标的直接进入正式模型池，核心原因仍是当前可正权外部估值锚不足或价格纪律不满足。

\noindent
\begin{exhibitbox}[图表 3｜补证后候选准入摘要]
\scriptsize
\begin{tabularx}{\textwidth}{L{2.5cm}R{0.9cm}L{4.0cm}X}
\toprule
\textbf{层级} & \textbf{数量} & \textbf{代表标的} & \textbf{使用方式} \\
\midrule
优先验证候选 & 6 & 吉林敖东、辽宁成大、中信金属、三六零、宁波华翔、海德股份 & 进入优先验证队列；等待外部锚、H2扣非和现金流确认 \\
扩展事件验证 & 11 & 九安医疗、焦作万方、神火股份、香农芯创、海正药业、高德红外 & 列入扩展观察，不因空间大直接升级 \\
安全边际观察 & 3 & 中远海能、赣锋锂业、越秀资本 & 低优先级跟踪，等待安全边际扩大 \\
风险/不纳入 & 28 & 渤海租赁、翔鹭钨业、申通快递、长江证券、中泰证券、上海石化 & 概率值低于现价或现金流/外部锚不足，不新增为候选 \\
\bottomrule
\end{tabularx}
\normalsize
\sourcenote{完整逐票准入结论见附录C；本表为投委会摘要，不构成正式推荐。}
\end{exhibitbox}

\section{为什么没有按100\%以上空间直接入选}

全142只候选中，本轮模型空间超过100\%的标的共有6只，其中3只进入39只优先池，0只进入正式模型。专项核验已对6/6只完成原始API、原始PDF、目标价字段与现金流闭环；当前正权原始目标为0只。因此未升级不是资料空白，而是已核验证据不支持把内部高空间转换为正式模型。

\noindent
\begin{exhibitbox}[图表 3A｜100\%以上空间样本的拦截链条]
\scriptsize
\begin{tabularx}{\textwidth}{L{1.05cm}L{1.55cm}R{1.0cm}L{1.35cm}L{1.35cm}X}
\toprule
\textbf{代码} & \textbf{标的} & \textbf{空间} & \textbf{优先池} & \textbf{正式池} & \textbf{未直接入选原因} \\
\midrule
002432 & 九安医疗 & 323.4\% & 否 & 否 & 未入优先池；当前原始目标0；Q1现金流-0.92亿元；保留候选观察，维持未进入优先池及非正式模型 \\
000623 & 吉林敖东 & 317.0\% & 是 & 否 & 优先池内；当前原始目标0；Q1现金流0.11亿元；保留优先池验证，维持非正式模型 \\
600739 & 辽宁成大 & 300.0\% & 是 & 否 & 优先池内；当前原始目标0；Q1现金流-1.76亿元；保留优先池验证，维持非正式模型 \\
000685 & 中山公用 & 162.8\% & 否 & 否 & 未入优先池；当前原始目标0；Q1现金流-0.27亿元；保留候选观察，维持未进入优先池及非正式模型 \\
600150 & 中国船舶 & 142.9\% & 是 & 否 & 优先池内；当前原始目标0；Q1现金流81.52亿元；保留优先池验证，维持非正式模型 \\
301308 & 江波龙 & 119.2\% & 否 & 否 & 未入优先池；当前原始目标0；Q1现金流-28.75亿元；保留候选观察，维持未进入优先池及非正式模型 \\
\bottomrule
\end{tabularx}
\normalsize
\sourcenote{专项包已逐票核验东方财富原始API、已归档原始PDF、目标价字段与Q1现金流。媒体转载、聚合页和失败探针权重均为0。}
\end{exhibitbox}


\section{重点事件验证清单}

下表不是推荐名单，而是最需要在下一轮信息披露中验证的高敏感样本。其共同点是：模型空间可能存在，但当前尚不满足正式审计模型条件。

\noindent
\begin{exhibitbox}[图表 3B｜条件性机会：从模型空间到证据升级]
\scriptsize
\begin{tabularx}{\textwidth}{L{1.1cm}L{1.7cm}L{1.6cm}R{1.0cm}R{1.0cm}X}
\toprule
\textbf{代码} & \textbf{标的} & \textbf{估值框架} & \textbf{概率值} & \textbf{空间} & \textbf{必须验证的事实} \\
\midrule
600150 & 中国船舶 & 成长盈利情景 PE & 83.33 & 142.9\% & 订单交付、合同负债、产品结构和应收账款回款；反证：订单确认延后、交付节奏放缓、应收账款继续扩张或估值透支交付 \\
601600 & 中国铝业 & 周期调整扣非 EPS / PE & 14.07 & 56.7\% & 金属价格、产量、单位成本、库存和经营现金流；反证：价格回落、成本曲线抬升、产量不达预期或利润无法转化为现金 \\
601360 & 三六零 & 2026E 收入 PS & 13.10 & 45.6\% & 合同订单、项目交付、回款、研发投入和利润率；反证：订单确认延后、回款周期拉长、研发费用率上升或估值继续压缩 \\
600026 & 中远海能 & 周期调整扣非 EPS / PE & 17.45 & 10.2\% & 运价、运力利用率、燃料成本、资本开支和现金流；反证：运价回落、利用率下降、成本上升或周期盈利被错误年化 \\
002048 & 宁波华翔 & 正常化 EPS / PE & 28.50 & 33.0\% & 销量、单车盈利、产品结构、价格竞争和应收回款；反证：销量或单车盈利下滑、价格战加剧、产品结构恶化或现金流转弱 \\
002812 & 恩捷股份 & 产能周期正常化前瞻 EPS / PE & 85.29 & 55.2\% & 出货量、订单、产品价格、产能利用率和营运资金；反证：订单转化慢、价格竞争、利用率不足或库存/应收继续占用现金 \\
300390 & 天华新能 & 周期调整 EPS / PE & 77.30 & 28.8\% & 出货量、订单、产品价格、产能利用率和营运资金；反证：订单转化慢、价格竞争、利用率不足或库存/应收继续占用现金 \\
002240 & 盛新锂能 & 周期调整 EPS / PE & 37.89 & 21.8\% & 金属价格、产量、单位成本、库存和经营现金流；反证：价格回落、成本曲线抬升、产量不达预期或利润无法转化为现金 \\
\bottomrule
\end{tabularx}
\normalsize
\sourcenote{条件性模型仅用于安排验证优先级。若缺乏当前外部锚或现金流确认，则不升级为可复核正向验证。}
\end{exhibitbox}

本报告的反方情景同样明确：若H2扣非利润不延续、经营现金流继续背离、商品价格或行业价差回落，当前候选池中的高空间模型将被迅速下调，而不是通过提高倍数维持结论。

\chapter{市场状态与机会架构}

市场层面的结论是“选择性再定价”而非“全面抬估值”。行业资金状态只能回答资本是否开始关注，不能替代公司盈利质量、现金流或估值分母。我们将行业状态与公司级证据分开处理，以避免把短期资金流误解为可配置的基本面信号。

\noindent
\begin{exhibitbox}[图表 4｜31个一级行业的阶段划分与投资含义]
\small
\begin{tabularx}{\textwidth}{L{1.9cm}L{3.5cm}X L{4.2cm}}
\toprule
\textbf{阶段} & \textbf{观察到的市场特征} & \textbf{代表行业} & \textbf{投资含义} \\
\midrule
静默吸纳 & 日/周资金同向流入且价格、估值仍受约束 & 农林牧渔、房地产 & 优先寻找利润质量与现金流同步改善的公司 \\
启动确认 & 资金与价格已出现确认，静态低位优势下降 & 石油石化、煤炭 & 关注利润兑现，不追逐单日放量 \\
资金观察 & 周度资金存在但价格、估值或日度流向尚未闭环 & 计算机、银行、传媒、国防军工 & 只做事件验证，不根据主题直接配置 \\
单日反弹 & 单日流入未获得周度确认 & 美容护理、轻工制造、食品饮料、商贸零售、钢铁、纺织服饰、社会服务、环保、汽车、建筑装饰、医药生物、通信、建筑材料、基础化工、有色金属、电力设备、电子 & 不把单日反弹视为资本布局 \\
\bottomrule
\end{tabularx}
\normalsize
\sourcenote{申万一级行业官方指数历史与公开行业资金表。行业资金观察截至2026年7月14日，周度观察截至2026年7月10日。}
\end{exhibitbox}

\section{从行业阶段到公司动作}

静默吸纳并不自动意味着低风险；例如房地产的行业PB处于低分位，但公司层面仍须通过结算、税费、债务和现金流检验。启动确认也不自动意味着高估；石油石化和煤炭的价格确认必须与价差、库存、分红和现金流相匹配。因此，公司动作由三重门槛决定：行业阶段、盈利可持续性、当前价格相对情景价值。

\noindent
\begin{tabularx}{\textwidth}{L{3.2cm}X L{4.0cm}}
\toprule
\textbf{问题} & \textbf{判断规则} & \textbf{若不成立} \\
\midrule
行业是否值得跟踪？ & 资金阶段、估值分位和价格位置至少有两项支持。 & 只保留为市场观察，不进入公司验证。 \\
利润是否可进入估值？ & 扣非为正、现金流不明显背离、无低基数或结算型扭曲。 & 启动恢复模型或保留为边界标的。 \\
价格是否可承受？ & 概率值相对现价有安全边际，且外部锚/分母时效可复核。 & 降级为事件验证或估值纪律观察。 \\
\bottomrule
\end{tabularx}

\chapter{估值框架与情景分析}

估值不是把同一PE模板扩展到142只股票。我们先判断利润属于周期、成长、金融资产、常规经营还是恢复性分母，再选择对应的估值方法。所有外部目标价均按来源质量加权；没有可审计当前锚时，模型只能是条件性本机构区间。

\noindent
\begin{exhibitbox}[图表 5｜估值方法选择矩阵]
\small
\begin{tabularx}{\textwidth}{L{1.0cm}L{2.2cm}R{1.0cm}X L{4.0cm}}
\toprule
\textbf{代码} & \textbf{方法} & \textbf{覆盖} & \textbf{核心公式} & \textbf{为何使用} \\
\midrule
M1 & 周期调整扣非EPS/PE & 53 & \(EPS_s=EPS_{H1}^d(1+r_s);\ V_s=EPS_s\times PE_s\) & 商品、价差、运价和利用率主导的盈利不能被单一半年度利润代表。 \\
M2 & 成长盈利情景PE & 42 & \(EPS_s=EPS_{H1}^d(1+r_s);\ V_s=EPS_s\times PE_s\) & 仅对已进入盈利桥接的成长持续期赋予更高情景倍数。 \\
M3 & PB/ROE与盈利混合 & 18 & \(V_s=0.65(BPS\times PB_s)+0.35(EPS_s\times PE_s)\) & 金融资产需要同时锚定资产负债表与盈利能力。 \\
M4 & 常规扣非EPS/PE & 16 & \(EPS_s=EPS_{H1}^d(1+r_s);\ V_s=EPS_s\times PE_s\) & 正扣非分母且未触发低基数或项目结算阻断。 \\
R1 & 恢复模型 & 12 & \(V_s=EPS_{norm}\times PE_s\) 或 \(BPS\times PB_s\) 或 \(R_{2026E}/Shares\times PS_s\) & 低基数、结算型、近零EPS必须先更换分母。 \\
B1 & 上市边界 & 1 & \(P_{target}=\text{不适用}\) & 无二级市场价格时不制造伪精确目标。 \\
\bottomrule
\end{tabularx}
\normalsize
\sourcenote{方法分布来自142只候选的逐票估值审计包；每只的变量代入、依据和证据路径见附录B。}
\end{exhibitbox}

\section{情景权重与目标价治理}

\begin{align*}
V_{\mathrm{prob}} &= 0.30V_{\mathrm{bear}}+0.50V_{\mathrm{base}}+0.20V_{\mathrm{bull}}, \\
P_{\mathrm{target}} &= V_{\mathrm{prob}}(1-w_{\mathrm{external}})+P_{\mathrm{external}}w_{\mathrm{external}}, \\
\mathrm{Upside} &= \frac{P_{\mathrm{target}}}{P_{\mathrm{current}}}-1.
\end{align*}
50%基准权重反映“正常兑现”是最常见路径；30%熊市权重强制保留周期、现金流和分母失效的风险；20%牛市权重仅在盈利、现金流和外部证据同步改善时才有意义。权重是统一的本机构情景先验，不是回归模型估计出的客观概率。

\noindent
\begin{sourcequalitybox}[外部锚与本机构区间的边界]
原始券商PDF或可审计共识快照可作为外部锚；摘要、转载、陈旧目标价和用户估值文章仅提供背景，权重为零。本机构区间的作用是定义验证顺序和风险边界，不替代当前、可审计的市场一致目标价。
\end{sourcequalitybox}

模型升级已在第一章的行动分层中固化：正式模型要求当前价格、分母、外部锚和失效条件同步可核；条件性区间仅用于安排验证优先级；上市或证据边界标的在形成有效交易价格或补齐关键证据前不进入目标价模型。

\noindent
\begin{exhibitbox}[图表 6｜概率值不是承诺：三项不可放松的估值约束]
\small
\begin{tabularx}{\textwidth}{L{2.2cm}X L{4.6cm}}
\toprule
\textbf{约束} & \textbf{模型含义} & \textbf{行动后果} \\
\midrule
分母约束 & H2扣非、现金流和周期变量必须支撑熊/基准/牛的情景分母；半年度利润不等同于全年盈利。 & 分母失效先下调EPS或切换正常化口径，不以提高估值倍数补偿。 \\
外部锚约束 & 只有当前、可审计的原始研报或共识快照可获得正权重；摘要、转载和陈旧目标价仅保留为背景。 & 弱来源权重归零，条件性本机构区间不升级为可复核正向验证。 \\
价格约束 & 概率值相对现价的空间只是情景结果；价格位置、估值拥挤度和失效条件决定是否可执行。 & 概率值低于现价或证据时效不足时不追价，转入估值纪律观察。 \\
\bottomrule
\end{tabularx}
\normalsize
\sourcenote{概率值用于比较风险调整后的情景价值，不构成收益承诺；逐票输入、权重与公式代入见附录B。}
\end{exhibitbox}

\noindent
\begin{exhibitbox}[图表 7｜正式模型的概率值桥接：本机构情景与外部锚分层计入]
\small
\begin{tabularx}{\textwidth}{L{2.5cm}R{2.4cm}L{4.1cm}R{2.4cm}}
\toprule
\textbf{标的} & \textbf{本机构概率值} & \textbf{外部锚及权重} & \textbf{最终概率值} \\
\midrule
恒逸石化（000703） & 24.13元 & 17.05元 × 10\% & 23.42元 \\
宏桥控股（002379） & 25.45元 & 29.00元 × 10\% & 25.80元 \\
亿纬锂能（300014） & 58.59元 & 77.00元 × 10\% & 60.43元 \\
川能动力（000155） & 8.88元 & 20.43元 × 10\% & 10.04元 \\
\bottomrule
\end{tabularx}
\normalsize
\sourcenote{最终概率值=本机构概率值×（1－外部锚权重）＋外部目标价×外部锚权重。外部锚的正权重仅限于当前、可审计的原始研报字段。}
\end{exhibitbox}

\chapter{焦点标的与验证路径}

本章聚焦能够改变下一轮配置排序的公司，而不是重述142只候选。我们把标的分为可复核当前价模型、恢复性估值模型和高空间但证据不足的验证观察三类；任何一类都必须有可观察的升级条件和明确的反证。

\section{可复核当前价模型}
\subsection{恒逸石化（000703）｜可复核正向验证模型}
现价14.48元，本机构概率值23.42元，对应61.7\%空间。熊/基准/牛值为10.86/25.59/40.41元；外部锚17.05元，权重10\%。
我们将其定义为可复核正向验证模型而非一般观察名单，原因是当前价格、股本、情景分母和外部证据可以同时复核。下一轮的关键验证是产品价差、原料成本、装置运行率、库存和现金流；若出现价差收窄、原料成本上升、装置检修或盈利恢复无法持续，则该模型应被降级。

\subsection{宏桥控股（002379）｜可复核正向验证模型}
现价18.66元，本机构概率值25.80元，对应38.3\%空间。熊/基准/牛值为14.00/26.24/40.65元；外部锚29.00元，权重10\%。
我们将其定义为可复核正向验证模型而非一般观察名单，原因是当前价格、股本、情景分母和外部证据可以同时复核。下一轮的关键验证是金属价格、产量、单位成本、库存和经营现金流；若出现价格回落、成本曲线抬升、产量不达预期或利润无法转化为现金，则该模型应被降级。

\subsection{亿纬锂能（300014）｜可复核正向验证模型}
现价52.28元，本机构概率值60.43元，对应15.6\%空间。熊/基准/牛值为31.49/68.46/74.57元；外部锚77.00元，权重10\%。
我们将其定义为可复核正向验证模型而非一般观察名单，原因是当前价格、股本、情景分母和外部证据可以同时复核。下一轮的关键验证是出货量、订单、产品价格、产能利用率和营运资金；若出现订单转化慢、价格竞争、利用率不足或库存/应收继续占用现金，则该模型应被降级。

\subsection{川能动力（000155）｜下行纪律对照}
现价11.09元，本机构概率值10.04元，对应-9.5\%空间。熊/基准/牛值为5.43/9.03/13.67元；外部锚20.43元，权重10\%。
我们将其定义为下行纪律对照而非一般观察名单，原因是当前价格、股本、情景分母和外部证据可以同时复核。下一轮的关键验证是利用小时、上网电价、燃料/资源成本、资本开支和现金流；若出现利用率或电价下行、成本抬升、资本开支超预期或现金流不达标，则该模型应被降级。

\noindent
\begin{exhibitbox}[图表 8｜正式模型横向比较：价值空间必须接受同一组验证纪律]
\small
\begin{tabularx}{\textwidth}{L{2.1cm}L{2.4cm}X L{4.0cm}}
\toprule
\textbf{标的} & \textbf{现价/概率值} & \textbf{本轮定位} & \textbf{下一轮必须验证} \\
\midrule
恒逸石化（000703） & 14.48/23.42元 & 可复核正向验证模型；61.7\%空间 & 验证：产品价差、原料成本、装置运行率、库存和现金流；反证：价差收窄、原料成本上升、装置检修或盈利恢复无法持续 \\
宏桥控股（002379） & 18.66/25.80元 & 可复核正向验证模型；38.3\%空间 & 验证：金属价格、产量、单位成本、库存和经营现金流；反证：价格回落、成本曲线抬升、产量不达预期或利润无法转化为现金 \\
亿纬锂能（300014） & 52.28/60.43元 & 可复核正向验证模型；15.6\%空间 & 验证：出货量、订单、产品价格、产能利用率和营运资金；反证：订单转化慢、价格竞争、利用率不足或库存/应收继续占用现金 \\
川能动力（000155） & 11.09/10.04元 & 下行纪律对照；-9.5\%空间 & 验证：利用小时、上网电价、燃料/资源成本、资本开支和现金流；反证：利用率或电价下行、成本抬升、资本开支超预期或现金流不达标 \\
\bottomrule
\end{tabularx}
\normalsize
\sourcenote{正式模型均要求当前价格、分母、原始外部锚和失效条件同时可复核；若概率值低于现价，则仅承担下行纪律对照功能。}
\end{exhibitbox}

\section{恢复性估值：关注分母修复，而非追逐表面空间}

恢复模型的共同纪律是：当H1利润由低基数、扭亏、结算或周期集中驱动时，不允许直接以H1年化EPS乘常规行业PE。下表展示当前最具代表性的恢复样本及其下一步证据要求。

\noindent
\begin{exhibitbox}[图表 9｜恢复性估值：分母切换与验证条件]
\scriptsize
\begin{tabularx}{\textwidth}{L{1.1cm}L{1.7cm}L{2.4cm}R{1.0cm}R{1.0cm}X}
\toprule
\textbf{代码} & \textbf{标的} & \textbf{替代分母/方法} & \textbf{概率值} & \textbf{空间} & \textbf{升级条件与反证} \\
\midrule
601360 & 三六零 & 2026E 收入 PS；2026E收入95.91亿元与AI/安全变现 & 13.10 & 45.6\% & 验证：2026年收入接近95.91亿元，AI/安全形成可重复变现，EPS共识不低于0.06元且经营现金流转正；反证：AI产品变现不及预期、2026年EPS低于0.03元、经营现金流仍为负，或PS低于可比区间 \\
000042 & 中洲控股 & 项目结算正常化 EPS / PE；项目结算、税费与利息正常化EPS & 8.54 & 6.5\% & 验证：2026年EPS不低于1.10元，项目结算及税费/利息计提得到确认，债务与现金流压力改善且股价守住8元；反证：结算利润回转、土地增值税或利息计提上升、销售回款走弱，或净负债/流动性恶化 \\
002460 & 赣锋锂业 & 周期调整 EPS / PE；2026E周期调整EPS & 53.68 & 4.2\% & 验证：锂价、产量、资源自给率和H2经营现金流共同验证2026年EPS 4.45元；反证：锂价回落、资源项目延期、库存减值，或现金流弱于利润 \\
002812 & 恩捷股份 & 产能周期正常化前瞻 EPS / PE；隔膜价格/利用率恢复后的2027E EPS & 85.29 & 55.2\% & 验证：隔膜涨价、产能利用率、海外/3C订单和H2现金流共同验证2027年EPS 5.16元；反证：供需修复不及预期、价格回落、客户集中，或产能利用率不足 \\
300390 & 天华新能 & 周期调整 EPS / PE；锂价敏感的2026E周期调整EPS & 77.30 & 28.8\% & 验证：氢氧化锂价格、客户采购、产量和H2经营现金流共同验证2026年EPS 5.70元；反证：锂价回落、客户集中风险暴露，或经营现金流继续弱于利润 \\
002240 & 盛新锂能 & 周期调整 EPS / PE；锂价、自有矿与现金流桥接EPS & 37.89 & 21.8\% & 验证：锂价、出货、自有矿贡献和H2经营现金流共同验证2026年EPS 2.45元；反证：锂价回落、现金流持续为负，或资源项目贡献不及预期 \\
600736 & 苏州高新 & 正常化 EPS / PE（PB交叉检验）；园区/结算/投资收益正常化EPS & 7.19 & 4.7\% & 验证：产业园利润、地产结算、投资收益和经营现金流共同验证2026年EPS 0.14元；反证：地产结算延期、投资收益回落，或经营现金流继续为负 \\
000415 & 渤海租赁 & 正常化 EPS / PE（PB交叉检验）；租赁收入、资产质量与汇率正常化EPS & 4.40 & -1.1\% & 验证：租赁收入、资产质量、汇率和H2现金流共同验证2026年EPS 0.77元；反证：飞机减值、汇率损失、融资成本上升，或现金流风险上行 \\
\bottomrule
\end{tabularx}
\normalsize
\sourcenote{逐票正常化分母、可比集、倍数和公式代入见附录B；表中区间均为条件性本机构价值。}
\end{exhibitbox}

\section{高空间模型的正确使用方式}

高空间不是排序第一的理由。本轮已对全部6只100\%+空间样本完成逐票证据闭环：原始API、已归档PDF、目标价字段、Q1经营现金流和池内准入均已核验，闭环率为6/6。核验结果不是“仍待补资料”，而是当前正权原始目标为0只、正式升级为0只。

已找到的历史目标仅作为反向证据：九安医疗70.28元、吉林敖东约18.02元、辽宁成大22.79元、中山公用10.50--11.11元、中国船舶31.36元，均显著低于或仅接近内部高空间情景值，且因时效不足权重为0。江波龙在原始API及已归档PDF语料中未找到目标价；中国船舶50元媒体转载仅作零权重交叉验证。

\noindent
\begin{exhibitbox}[图表 9A｜高空间样本的已核验证据与准入结论]
\scriptsize
\begin{tabularx}{\textwidth}{L{0.95cm}L{1.45cm}R{0.95cm}L{5.0cm}X}
\toprule
\textbf{代码} & \textbf{标的} & \textbf{空间} & \textbf{已核验证据} & \textbf{准入结论/剩余事件} \\
\midrule
002432 & 九安医疗 & 323.4\% & API/PDF 2/2；当前原始目标0；2024-12-23 70.28元，权重0\%；Q1现金流-0.92亿元 & 保留候选观察，维持未进入优先池及非正式模型；未来：重新满足优先池的价格位置与行业阶段规则；H2扣非利润持续性与经营现金流转正 \\
000623 & 吉林敖东 & 317.0\% & API/PDF 2/2；当前原始目标0；2024-03-27 18.02元，权重0\%；Q1现金流0.11亿元 & 保留优先池验证，维持非正式模型；未来：2026-04-01之后原始券商目标价或合理价值区间；H2扣非利润持续性与投资收益质量 \\
600739 & 辽宁成大 & 300.0\% & API/PDF 4/4；当前原始目标0；2017-10-30 22.79元，权重0\%；Q1现金流-1.76亿元 & 保留优先池验证，维持非正式模型；未来：2026-04-01之后原始券商目标价或合理价值区间；H2扣非利润持续性与广发证券投资收益质量 \\
000685 & 中山公用 & 162.8\% & API/PDF 6/6；当前原始目标0；2025-03-04 10.50--11.11元，权重0\%；Q1现金流-0.27亿元 & 保留候选观察，维持未进入优先池及非正式模型；未来：重新满足优先池的价格位置与行业阶段规则；H2扣非利润与投资收益可持续性 \\
600150 & 中国船舶 & 142.9\% & API/PDF 56/8；当前原始目标0；2023-03-27 31.36元，权重0\%；Q1现金流81.52亿元 & 保留优先池验证，维持非正式模型；未来：当前原始券商目标价或合理价值区间；H2订单交付、船价、利润率与应收回款 \\
301308 & 江波龙 & 119.2\% & API/PDF 12/8；当前原始目标0；原始语料未找到可接受目标价；Q1现金流-28.75亿元 & 保留候选观察，维持未进入优先池及非正式模型；未来：重新满足优先池的价格位置与行业阶段规则；当前原始券商目标价或合理价值区间 \\
\bottomrule
\end{tabularx}
\normalsize
\sourcenote{专项证据包：\texttt{data/high\_upside\_evidence\_closure\_20260716.json}。当前原始目标缺失已按“核验未披露”关闭；H2利润、现金流、订单与价格周期保留为未来事件验证。}
\end{exhibitbox}

\chapter{风险、催化与监控}

本报告的核心风险不在于“预告会不会变差”，而在于市场将短期利润、周期价差或项目结算误读为可持续盈利。监控应围绕能改变分母、倍数和外部锚的变量，而不是围绕新闻标题。

\noindent
\begin{exhibitbox}[图表 10｜风险监控框架：从风险描述到可执行触发器]
\small
\begin{tabularx}{\textwidth}{L{2.2cm}L{4.0cm}L{4.4cm}X}
\toprule
\textbf{风险类别} & \textbf{领先指标} & \textbf{受影响模型} & \textbf{行动规则} \\
\midrule
盈利质量 & 扣非与归母差额扩大、非经常性收益上升、H2利润低于桥接。 & 常规EPS/PE、成长PE与恢复模型。 & 下调分母或切换至恢复模型；不以提高PE抵消盈利下修。 \\
现金流与营运资本 & 经营现金流持续低于利润、应收/库存占用上升。 & 周期、成长和项目型公司。 & 暂停升级；将牛市情景权重下调或取消。 \\
周期与价格 & 商品价格、炼化价差、运价、锂价、利用率反转。 & 周期调整PE与恢复模型。 & 先下调H2转化率，再下调PE区间；不直接沿用H1盈利。 \\
估值与拥挤度 & 当前价高于概率值、外部目标价陈旧、短期涨幅过大。 & 全部条件性模型。 & 转为观察；等待价格回撤或外部锚更新。 \\
证据质量 & 原始研报不可得、仅有摘要/转载、数据时点过旧。 & 外部锚和正式模型。 & 外部权重归零；仅保留本机构条件性区间。 \\
\bottomrule
\end{tabularx}
\normalsize
\end{exhibitbox}

\section{下一轮催化日历}

\noindent
\begin{tabularx}{\textwidth}{L{2.4cm}L{4.5cm}X}
\toprule
\textbf{时点} & \textbf{需要确认的变量} & \textbf{模型含义} \\
\midrule
下一次业绩披露 & H2扣非利润、毛利率、现金流和非经常性收益。 & 决定H2/H1转换率是否成立，影响全部EPS类模型。 \\
行业价格与订单更新 & 金属、能源、化工价差、运价、锂价、订单和交付。 & 决定周期模型的熊/基准/牛倍数与分母可持续性。 \\
当前研报或共识更新 & 2026E/2027E收入、EPS、目标价及评级变化。 & 决定外部锚是否可新增正权重，条件性模型是否可升级。 \\
资金与价格确认 & 行业日/周资金、价格位置、成交和回撤承接。 & 只用于确认市场是否接受基本面，不替代基本面证据。 \\
\bottomrule
\end{tabularx}

\noindent
\begin{sourcequalitybox}[数据边界]
候选Q1财务覆盖142/142，研报PDF覆盖137/142。行业日/周资金均为31/31完整表，但观察日分别截至7月14日和7月10日；连续流入公开表仅30/112，只作为辅助证据。
\end{sourcequalitybox}

\chapter{研究基础与数据边界}

本报告将事实、估计和观点分层呈现。事实层来自官方业绩预告、Q1财务包、行情和行业数据；估计层来自经审计的情景模型与有限的原始研报/公开共识；观点层只用于定义验证顺序，不替代事实或外部锚。

\noindent
\begin{tabularx}{\textwidth}{L{2.1cm}L{4.0cm}X}
\toprule
\textbf{证据层级} & \textbf{用途} & \textbf{边界} \\
\midrule
一级：官方披露与市场数据 & 半年报预告、Q1财务、当前价格、行业指数和资金表。 & 反映已发生或可观察数据，不证明未来盈利持续性。 \\
二级：原始券商PDF与可审计共识快照 & 盈利预测、目标价、估值方法和可比框架。 & 仅当前、可比较字段可获得正权重。 \\
三级：摘要、转载、陈旧目标与用户文章 & 识别市场叙事、历史区间与待验证假设。 & 权重为零，不能被呈现为市场一致预期。 \\
本机构情景 & 估值区间、概率权重、行动门槛与失效条件。 & 明确标注为本机构假设，需由下一轮证据验证。 \\
\bottomrule
\end{tabularx}

全量候选的逐票公式、输入变量、方法理由、倍数理由、概率权重和证据路径均在附录B保留。这使主报告保持可读，同时不牺牲模型的可复算性。

\noindent
\begin{disclosurebox}[使用边界]
“主力资金”沿用公开数据源按成交单大小分类的统计口径，不代表对资金最终受益人、机构身份或一致行动关系的确认。行业阶段使用截至2026年7月14日的完整31行业日报和截至2026年7月10日的完整周报；连续流入公开表仅30/112，已按部分证据处理。全部目标价均为概率情景研究框架，不构成收益承诺或证券买卖建议。
\end{disclosurebox}

\section{来源与复现说明}

主要研究输入包括：官方2026H1业绩预告全量表、候选公司Q1财务包、2026年7月15日价格与复权K线、已归档的原始券商PDF、可审计公开共识快照，以及截至2026年7月14日的行业资金观察。详细文件路径、字段版本与重算结果由附录B的逐票账本和结构化数据包保存。

\noindent
\begin{exhibitbox}[研究数据覆盖与复现状态]
\begin{tabularx}{\textwidth}{L{4.5cm}R{2.0cm}X}
\toprule
\textbf{项目} & \textbf{覆盖} & \textbf{复现说明} \\
\midrule
官方预告公司 & 1680 & 排除B股后按申万行业映射。 \\
高影响候选Q1财务 & 142/142 & 作为盈利质量、现金流和BPS输入。 \\
候选研报PDF & 137/142 & 原始PDF可用于外部锚或方法/可比背景。 \\
逐票估值审计 & 142/142 & 每只候选均有公式、理由、证据与审计状态。 \\
\bottomrule
\end{tabularx}
\sourcenote{完整结构化数据位于本研究案例的数据包中；读者可通过附录B的代码与方法索引追溯每个估值结论。}
\end{exhibitbox}

研究结论反映截至数据截止日的公开资料与本机构情景；附录的结构化审计包是可复算记录，而不是对任何个股的确定性价格预测。

\clearpage
\appendix
\chapter{全量候选清单与操作分层}

本附录保留142只候选的当前价格、盈利规模、概率值、空间、证据质量和行动标签。它是全量筛选账本，不是142只股票的统一推荐清单。

\footnotesize
\renewcommand{\arraystretch}{1.00}
\setlength{\tabcolsep}{2pt}
\begin{longtable}{L{1.0cm}L{1.55cm}L{1.4cm}R{1.0cm}R{1.0cm}R{1.0cm}R{1.05cm}L{3.6cm}L{0.9cm}}
\toprule
\textbf{代码} & \textbf{公司} & \textbf{行业} & \textbf{现价} & \textbf{H1扣非} & \textbf{概率值} & \textbf{空间} & \textbf{行动含义} & \textbf{证据} \\
\midrule
\endfirsthead
\toprule
\textbf{代码} & \textbf{公司} & \textbf{行业} & \textbf{现价} & \textbf{H1扣非} & \textbf{概率值} & \textbf{空间} & \textbf{行动含义} & \textbf{证据} \\
\midrule
\endhead
000042 & 中洲控股 & 房地产 & 8.02元 & 10.6 & 9元 & 6.5\% & 市场观察 & 基础高 \\
600150 & 中国船舶 & 国防军工 & 34.31元 & 99.0 & 83元 & 142.9\% & 高空间验证 & 基础高 \\
002379 & 宏桥控股 & 有色金属 & 18.66元 & 154.0 & 26元 & 38.3\% & 高空间验证 & 基础高 \\
601600 & 中国铝业 & 有色金属 & 8.98元 & 115.0 & 14元 & 56.7\% & 高空间验证 & 中低 \\
600346 & 恒力石化 & 石油石化 & 15.24元 & 53.3 & 14元 & -6.8\% & 估值观察 & 基础高 \\
000858 & 五粮液 & 食品饮料 & 76.40元 & 87.0 & 75元 & -1.5\% & 估值观察 & 基础高 \\
002460 & 赣锋锂业 & 有色金属 & 51.50元 & 36.0 & 54元 & 4.2\% & 市场观察 & 基础高 \\
002466 & 天齐锂业 & 有色金属 & 47.43元 & 35.0 & 44元 & -6.4\% & 估值观察 & 基础高 \\
601688 & 华泰证券 & 非银金融 & 20.63元 & 115.2 & 23元 & 9.4\% & 市场观察 & 基础高 \\
601336 & 新华保险 & 非银金融 & 65.55元 & 222.7 & 78元 & 19.7\% & 回撤验证 & 基础高 \\
600037 & 歌华有线 & 传媒 & 6.80元 & 2.1 & 6元 & -17.4\% & 高估值 & 基础高 \\
002738 & 中矿资源 & 有色金属 & 49.33元 & 11.0 & 39元 & -20.4\% & 高估值 & 基础高 \\
601069 & 西部黄金 & 有色金属 & 22.96元 & 5.3 & 13元 & -44.5\% & 高估值 & 基础高 \\
000567 & 海德股份 & 非银金融 & 6.17元 & 14.4 & 7元 & 20.8\% & 回撤验证 & 基础高 \\
002602 & 世纪华通 & 传媒 & 14.26元 & 43.1 & 17元 & 15.8\% & 回撤验证 & 基础高 \\
600026 & 中远海能 & 交通运输 & 15.83元 & 44.0 & 17元 & 10.2\% & 市场观察 & 基础高 \\
002048 & 宁波华翔 & 汽车 & 21.43元 & 5.3 & 28元 & 33.0\% & 高空间验证 & 基础高 \\
600688 & 上海石化 & 石油石化 & 2.84元 & 3.1 & 3元 & -8.1\% & 估值观察 & 基础高 \\
600095 & 湘财股份 & 非银金融 & 7.81元 & 5.2 & 4元 & -46.2\% & 高估值 & 基础高 \\
000426 & 兴业银锡 & 有色金属 & 30.56元 & 18.1 & 24元 & -22.9\% & 高估值 & 基础高 \\
600489 & 中金黄金 & 有色金属 & 20.12元 & 43.0 & 21元 & 5.2\% & 市场观察 & 基础高 \\
000737 & 北方铜业 & 有色金属 & 12.75元 & 13.5 & 15元 & 20.8\% & 回撤验证 & 中低 \\
002756 & 永兴材料 & 有色金属 & 44.90元 & 10.0 & 42元 & -6.2\% & 估值观察 & 基础高 \\
300014 & 亿纬锂能 & 电力设备 & 52.28元 & 25.2 & 60元 & 15.6\% & 回撤验证 & 基础高 \\
601061 & 中信金属 & 商贸零售 & 10.82元 & 26.4 & 17元 & 60.5\% & 高空间验证 & 基础高 \\
600739 & 辽宁成大 & 医药生物 & 10.95元 & 15.8 & 44元 & 300.0\% & 高空间验证 & 基础高 \\
601360 & 三六零 & 计算机 & 9.00元 & 1.6 & 13元 & 45.6\% & 高空间验证 & 基础高 \\
600918 & 中泰证券 & 非银金融 & 5.71元 & 17.0 & 5元 & -6.3\% & 估值观察 & 基础高 \\
000975 & 山金国际 & 有色金属 & 19.34元 & 23.4 & 20元 & 5.0\% & 市场观察 & 基础高 \\
000623 & 吉林敖东 & 医药生物 & 18.66元 & 22.2 & 78元 & 317.0\% & 高空间验证 & 基础高 \\
600292 & 电投水电 & 公用事业 & 12.50元 & 11.8 & 7元 & -47.5\% & 高估值 & 基础高 \\
000987 & 越秀资本 & 非银金融 & 7.77元 & 23.3 & 8元 & 1.4\% & 市场观察 & 基础高 \\
000155 & 川能动力 & 公用事业 & 11.09元 & 6.8 & 10元 & -9.5\% & 估值观察 & 基础高 \\
601377 & 兴业证券 & 非银金融 & 6.21元 & 20.9 & 6元 & -0.8\% & 估值观察 & 基础高 \\
600120 & 浙江东方 & 非银金融 & 4.91元 & 8.4 & 5元 & 2.2\% & 市场观察 & 中低 \\
000703 & 恒逸石化 & 石油石化 & 14.48元 & 57.2 & 23元 & 61.7\% & 高空间验证 & 基础高 \\
000301 & 东方盛虹 & 石油石化 & 11.61元 & 44.2 & 11元 & -3.7\% & 估值观察 & 基础高 \\
001339 & 智微智能 & 计算机 & 89.21元 & 3.7 & 65元 & -26.7\% & 高估值 & 基础高 \\
002558 & 巨人网络 & 传媒 & 29.56元 & 23.0 & 34元 & 15.0\% & 市场观察 & 基础高 \\
301308 & 江波龙 & 电子 & 488.00元 & 97.5 & 1069元 & 119.2\% & 高空间验证 & 基础高 \\
601628 & 中国人寿 & 非银金融 & 40.42元 & 1332.8 & 49元 & 20.2\% & 回撤验证 & 基础高 \\
002493 & 荣盛石化 & 石油石化 & 11.10元 & 53.0 & 9元 & -20.0\% & 高估值 & 基础高 \\
001309 & 德明利 & 电子 & 662.40元 & 60.5 & 482元 & -27.3\% & 高估值 & 基础高 \\
601233 & 桐昆股份 & 石油石化 & 20.72元 & 40.8 & 29元 & 38.0\% & 高空间验证 & 基础高 \\
603986 & 兆易创新 & 电子 & 571.79元 & 48.5 & 351元 & -38.6\% & 高估值 & 基础高 \\
601225 & 陕西煤业 & 煤炭 & 24.61元 & 97.5 & 18元 & -28.4\% & 高估值 & 基础高 \\
601138 & 工业富联 & 电子 & 63.62元 & 232.0 & 69元 & 8.1\% & 市场观察 & 基础高 \\
000807 & 云铝股份 & 有色金属 & 26.10元 & 75.5 & 45元 & 72.5\% & 高空间验证 & 基础高 \\
603993 & 洛阳钼业 & 有色金属 & 17.91元 & 155.0 & 16元 & -8.2\% & 估值观察 & 基础高 \\
601872 & 招商轮船 & 交通运输 & 15.26元 & 68.9 & 22元 & 41.4\% & 高空间验证 & 基础高 \\
600999 & 招商证券 & 非银金融 & 20.90元 & 105.0 & 19元 & -11.4\% & 估值观察 & 基础高 \\
601899 & 紫金矿业 & 有色金属 & 28.83元 & 379.0 & 35元 & 21.9\% & 回撤验证 & 基础高 \\
600309 & 万华化学 & 基础化工 & 70.09元 & 99.0 & 80元 & 14.3\% & 市场观察 & 基础高 \\
600989 & 宝丰能源 & 基础化工 & 22.07元 & 91.5 & 32元 & 43.5\% & 高空间验证 & 基础高 \\
300475 & 香农芯创 & 电子 & 221.20元 & 36.4 & 373元 & 68.4\% & 高空间验证 & 基础高 \\
000776 & 广发证券 & 非银金融 & 22.78元 & 116.9 & 23元 & 0.3\% & 市场观察 & 基础高 \\
600030 & 中信证券 & 非银金融 & 28.60元 & 236.1 & 25元 & -12.7\% & 估值观察 & 基础高 \\
601088 & 中国神华 & 煤炭 & 43.69元 & 271.0 & 22元 & -50.4\% & 高估值 & 基础高 \\
600362 & 江西铜业 & 有色金属 & 42.28元 & 81.0 & 51元 & 20.1\% & 回撤验证 & 基础高 \\
002709 & 天赐材料 & 电力设备 & 37.55元 & 28.0 & 60元 & 61.0\% & 高空间验证 & 基础高 \\
002414 & 高德红外 & 国防军工 & 13.32元 & 13.2 & 21元 & 56.4\% & 高空间验证 & 基础高 \\
300390 & 天华新能 & 电力设备 & 60.01元 & 22.6 & 77元 & 28.8\% & 回撤验证 & 基础高 \\
002648 & 卫星化学 & 基础化工 & 23.43元 & 61.9 & 42元 & 80.9\% & 高空间验证 & 基础高 \\
000792 & 盐湖股份 & 基础化工 & 25.27元 & 61.5 & 29元 & 16.5\% & 回撤验证 & 基础高 \\
601869 & 长飞光纤 & 通信 & 425.45元 & 23.0 & 148元 & -65.3\% & 高估值 & 基础高 \\
002432 & 九安医疗 & 医药生物 & 66.99元 & 31.5 & 284元 & 323.4\% & 高空间验证 & 基础高 \\
000657 & 中钨高新 & 有色金属 & 69.80元 & 20.6 & 23元 & -67.0\% & 高估值 & 基础高 \\
601995 & 中金公司 & 非银金融 & 38.34元 & 78.1 & 26元 & -31.7\% & 高估值 & 基础高 \\
002497 & 雅化集团 & 基础化工 & 18.07元 & 12.2 & 32元 & 76.0\% & 高空间验证 & 基础高 \\
000933 & 神火股份 & 有色金属 & 25.51元 & 48.1 & 44元 & 73.3\% & 高空间验证 & 基础高 \\
300223 & 北京君正 & 电子 & 190.18元 & 11.5 & 121元 & -36.4\% & 高估值 & 基础高 \\
002192 & 融捷股份 & 有色金属 & 66.42元 & 10.0 & 84元 & 26.2\% & 回撤验证 & 基础高 \\
601066 & 中信建投 & 非银金融 & 27.82元 & 77.9 & 15元 & -46.5\% & 高估值 & 基础高 \\
002812 & 恩捷股份 & 电力设备 & 54.96元 & 8.5 & 85元 & 55.2\% & 高空间验证 & 基础高 \\
002636 & 金安国纪 & 电子 & 94.79元 & 7.8 & 54元 & -42.5\% & 高估值 & 基础高 \\
002378 & 章源钨业 & 有色金属 & 25.84元 & 6.8 & 12元 & -52.4\% & 高估值 & 基础高 \\
601211 & 国泰海通 & 非银金融 & 18.58元 & 195.0 & 20元 & 10.1\% & 市场观察 & 基础高 \\
600884 & 杉杉股份 & 电力设备 & 12.14元 & 7.8 & 15元 & 21.5\% & 回撤验证 & 基础高 \\
002842 & 翔鹭钨业 & 有色金属 & 35.96元 & 5.3 & 35元 & -1.3\% & 估值观察 & 基础高 \\
002407 & 多氟多 & 基础化工 & 33.88元 & 4.9 & 10元 & -69.3\% & 高估值 & 基础高 \\
002653 & 海思科 & 医药生物 & 74.13元 & 7.3 & 32元 & -57.0\% & 高估值 & 基础高 \\
603629 & 利通电子 & 电子 & 112.09元 & 6.1 & 85元 & -24.4\% & 高估值 & 基础高 \\
600961 & 株冶集团 & 有色金属 & 24.76元 & 19.1 & 38元 & 51.8\% & 高空间验证 & 基础高 \\
002128 & 电投能源 & 煤炭 & 26.51元 & 40.5 & 24元 & -9.5\% & 估值观察 & 基础高 \\
601168 & 西部矿业 & 有色金属 & 33.18元 & 41.9 & 38元 & 15.1\% & 回撤验证 & 基础高 \\
600188 & 兖矿能源 & 煤炭 & 20.23元 & 45.0 & 10元 & -50.8\% & 高估值 & 基础高 \\
603162 & 海通发展 & 交通运输 & 11.24元 & 5.2 & 9元 & -24.1\% & 高估值 & 基础高 \\
002240 & 盛新锂能 & 有色金属 & 31.11元 & 14.0 & 38元 & 21.8\% & 回撤验证 & 基础高 \\
002440 & 闰土股份 & 基础化工 & 11.29元 & 4.4 & 10元 & -12.1\% & 估值观察 & 基础高 \\
002532 & 天山铝业 & 有色金属 & 12.75元 & 40.9 & 20元 & 60.5\% & 高空间验证 & 基础高 \\
301377 & 鼎泰高科 & 机械设备 & 467.88元 & 6.6 & 61元 & -87.1\% & 高估值 & 基础高 \\
688257 & 新锐股份 & 机械设备 & 80.06元 & 5.9 & 64元 & -20.7\% & 高估值 & 基础高 \\
000408 & 藏格矿业 & 基础化工 & 71.50元 & 37.1 & 71元 & -0.1\% & 估值观察 & 基础高 \\
600601 & 方正科技 & 电子 & 11.71元 & 5.5 & 7元 & -43.6\% & 高估值 & 基础高 \\
002475 & 立讯精密 & 电子 & 60.47元 & 64.9 & 43元 & -29.4\% & 高估值 & 基础高 \\
600595 & 中孚实业 & 有色金属 & 6.46元 & 19.2 & 11元 & 66.6\% & 高空间验证 & 基础高 \\
600487 & 亨通光电 & 通信 & 71.05元 & 33.3 & 70元 & -1.2\% & 估值观察 & 基础高 \\
600183 & 生益科技 & 电子 & 152.43元 & 28.2 & 63元 & -58.4\% & 高估值 & 基础高 \\
000630 & 铜陵有色 & 有色金属 & 6.06元 & 28.5 & 6元 & -7.4\% & 估值观察 & 基础高 \\
600549 & 厦门钨业 & 有色金属 & 58.08元 & 21.8 & 32元 & -44.7\% & 高估值 & 基础高 \\
000100 & TCL科技 & 电子 & 5.02元 & 30.3 & 8元 & 61.2\% & 高空间验证 & 基础高 \\
001287 & 中电港 & 电子 & 26.65元 & 5.2 & 41元 & 52.0\% & 高空间验证 & 基础高 \\
600111 & 北方稀土 & 有色金属 & 39.56元 & 20.3 & 13元 & -68.0\% & 高估值 & 基础高 \\
000725 & 京东方A & 电子 & 6.38元 & 31.9 & 6元 & -13.6\% & 估值观察 & 基础高 \\
600176 & 中国巨石 & 建筑材料 & 53.00元 & 29.8 & 28元 & -46.8\% & 高估值 & 基础高 \\
002266 & 浙富控股 & 环保 & 5.09元 & 12.5 & 8元 & 52.6\% & 高空间验证 & 基础高 \\
000612 & 焦作万方 & 有色金属 & 10.84元 & 12.1 & 21元 & 90.6\% & 高空间验证 & 基础高 \\
605117 & 德业股份 & 电力设备 & 85.40元 & 25.2 & 94元 & 9.9\% & 市场观察 & 基础高 \\
000783 & 长江证券 & 非银金融 & 9.06元 & 31.8 & 9元 & -3.8\% & 估值观察 & 基础高 \\
600233 & 圆通速递 & 交通运输 & 17.50元 & 31.9 & 21元 & 18.1\% & 回撤验证 & 基础高 \\
600392 & 盛和资源 & 有色金属 & 23.86元 & 8.6 & 11元 & -55.6\% & 高估值 & 基础高 \\
603225 & 新凤鸣 & 基础化工 & 17.97元 & 15.2 & 23元 & 26.5\% & 回撤验证 & 基础高 \\
002008 & 大族激光 & 机械设备 & 114.36元 & 14.0 & 52元 & -54.2\% & 高估值 & 基础高 \\
000685 & 中山公用 & 环保 & 11.19元 & 14.3 & 29元 & 162.8\% & 高空间验证 & 基础高 \\
600909 & 华安证券 & 非银金融 & 9.50元 & 21.6 & 7元 & -29.5\% & 高估值 & 基础高 \\
000833 & 粤桂股份 & 综合 & 20.07元 & 6.5 & 27元 & 34.7\% & 高空间验证 & 基础高 \\
300037 & 新宙邦 & 电力设备 & 67.20元 & 9.8 & 69元 & 2.5\% & 市场观察 & 基础高 \\
000060 & 中金岭南 & 有色金属 & 6.26元 & 11.1 & 6元 & -8.9\% & 估值观察 & 基础高 \\
600522 & 中天科技 & 通信 & 40.70元 & 22.3 & 33元 & -18.3\% & 高估值 & 基础高 \\
002468 & 申通快递 & 交通运输 & 15.22元 & 10.1 & 15元 & -2.2\% & 估值观察 & 基础高 \\
600267 & 海正药业 & 医药生物 & 11.98元 & 5.1 & 20元 & 65.7\% & 高空间验证 & 基础高 \\
600906 & 财达证券 & 非银金融 & 7.25元 & 7.6 & 4元 & -40.6\% & 高估值 & 基础高 \\
603268 & 松发股份 & 国防军工 & 160.46元 & 35.0 & 243元 & 51.7\% & 高空间验证 & 基础高 \\
002458 & 益生股份 & 农林牧渔 & 9.02元 & 2.8 & 8元 & -15.0\% & 估值观察 & 基础高 \\
002384 & 东山精密 & 电子 & 262.49元 & 24.5 & 82元 & -68.7\% & 高估值 & 基础高 \\
000977 & 浪潮信息 & 计算机 & 84.70元 & 23.1 & 96元 & 13.8\% & 市场观察 & 基础高 \\
603127 & 昭衍新药 & 医药生物 & 50.62元 & 7.0 & 46元 & -9.6\% & 估值观察 & 基础高 \\
301200 & 大族数控 & 机械设备 & 312.00元 & 9.5 & 80元 & -74.4\% & 高估值 & 基础高 \\
000938 & 紫光股份 & 计算机 & 35.53元 & 18.1 & 39元 & 8.5\% & 市场观察 & 基础高 \\
002463 & 沪电股份 & 电子 & 137.23元 & 28.1 & 75元 & -45.3\% & 高估值 & 基础高 \\
688183 & 生益电子 & 电子 & 128.87元 & 11.1 & 67元 & -47.8\% & 高估值 & 基础高 \\
300604 & 长川科技 & 电子 & 304.44元 & 9.1 & 109元 & -64.3\% & 高估值 & 基础高 \\
002916 & 深南电路 & 电子 & 377.23元 & 21.8 & 178元 & -52.9\% & 高估值 & 基础高 \\
688825 & 长鑫科技 & 电子 & 未上市；发行价8.66元，暂无二级市场当前价 & 未披露 & 不适用；待上市交易 & 不适用 & 不可定价 & 低 \\
001248 & 华润新能源 & 公用事业 & 13.41元 & 35.4 & 7元 & -50.7\% & 质量排除 & 基础高 \\
601633 & 长城汽车 & 汽车 & 15.58元 & 16.2 & 11元 & -27.8\% & 质量排除 & 基础高 \\
600736 & 苏州高新 & 房地产 & 6.87元 & 3.0 & 7元 & 4.7\% & 质量排除 & 基础高 \\
000415 & 渤海租赁 & 非银金融 & 4.45元 & 17.0 & 4元 & -1.1\% & 质量排除 & 基础高 \\
002221 & 东华能源 & 石油石化 & 5.44元 & 0.9 & 1元 & -81.4\% & 质量排除 & 基础高 \\
002611 & 东方精工 & 机械设备 & 15.02元 & 0.9 & 3元 & -81.2\% & 质量排除 & 基础高 \\
002174 & 游族网络 & 传媒 & 12.54元 & 0.2 & 1元 & -95.4\% & 质量排除 & 基础高 \\
603296 & 华勤技术 & 电子 & 78.20元 & 16.8 & 78元 & -0.1\% & 质量排除 & 基础高 \\
\bottomrule
\end{longtable}
\renewcommand{\arraystretch}{1.12}
\normalsize

\chapter{全量估值推演与证据账本}

本附录保留每只候选的估值公式代入、情景结果、方法选择逻辑和证据角色。方法代码对应第三章的估值选择矩阵；所有数值按数据截止日重算，未四舍五入输入保存在结构化审计包中。
\section{M1｜周期盈利：以周期调整扣非EPS/PE估值}
适用于商品、价差、运价和利用率主导的行业。重点不是把H1利润线性外推，而是用熊/基准/牛H2转化率测试周期持续性，并以较低熊市倍数吸收周期反转。

\footnotesize
\setlength{\tabcolsep}{2.5pt}
\begin{longtable}{L{1.55cm}L{2.45cm}L{5.0cm}L{2.75cm}L{4.1cm}}
\toprule
\textbf{代码/标的} & \textbf{适用理由} & \textbf{分母与公式代入} & \textbf{情景与概率值} & \textbf{证据与验证} \\
\midrule
\endfirsthead
\toprule
\textbf{代码/标的} & \textbf{适用理由} & \textbf{分母与公式代入} & \textbf{情景与概率值} & \textbf{证据与验证} \\
\midrule
\endhead
002379 宏桥控股 & 周期变量主导，禁止直接年化H1利润 & H1扣非EPS 1.1818元；H2/H1=0.55/0.85/1.15；\(\mathrm{EPS}_s=1.1818\times(1+\{0.55,0.85,1.15\});\ V_s=\mathrm{EPS}_s\times\{8.0,12.0,16.0\}\Rightarrow\{14.00,26.24,40.65\}\) & 熊/基准/牛：14.00/\allowbreak26.24/\allowbreak40.65元；概率值：25.80元；空间：38.3\%。 & 质量基础高；事实输入：官方预告、Q1财务与行情；外部背景：招银国际（2026-05-13）。监控：验证：金属价格、产量、单位成本、库存和经营现金流；反证：价格回落、成本曲线抬升、产量不达预期或利润无法转化为现金 \\
601600 中国铝业 & 周期变量主导，禁止直接年化H1利润 & H1扣非EPS 0.6704元；H2/H1=0.55/0.85/1.15；\(\mathrm{EPS}_s=0.6704\times(1+\{0.55,0.85,1.15\});\ V_s=\mathrm{EPS}_s\times\{8.0,12.0,16.0\}\Rightarrow\{6.74,14.88,23.06\}\) & 熊/基准/牛：6.74/\allowbreak14.88/\allowbreak23.06元；概率值：14.07元；空间：56.7\%。 & 质量中低；事实输入：官方预告、Q1财务与行情；外部背景：无当前券商目标锚。监控：验证：金属价格、产量、单位成本、库存和经营现金流；反证：价格回落、成本曲线抬升、产量不达预期或利润无法转化为现金 \\
600346 恒力石化 & 周期变量主导，禁止直接年化H1利润 & H1扣非EPS 0.7572元；H2/H1=0.55/0.90/1.25；\(\mathrm{EPS}_s=0.7572\times(1+\{0.55,0.90,1.25\});\ V_s=\mathrm{EPS}_s\times\{6.0,9.0,12.0\}\Rightarrow\{7.04,16.02,20.44\}\) & 熊/基准/牛：7.04/\allowbreak16.02/\allowbreak20.44元；概率值：14.21元；空间：-6.8\%。 & 质量基础高；事实输入：官方预告、Q1财务与行情；外部背景：中邮证券（2026-07-10）。监控：验证：产品价差、原料成本、装置运行率、库存和现金流；反证：价差收窄、原料成本上升、装置检修或盈利恢复无法持续 \\
002466 天齐锂业 & 周期变量主导，禁止直接年化H1利润 & H1扣非EPS 2.0459元；H2/H1=0.55/0.85/1.15；\(\mathrm{EPS}_s=2.0459\times(1+\{0.55,0.85,1.15\});\ V_s=\mathrm{EPS}_s\times\{8.0,12.0,16.0\}\Rightarrow\{25.37,45.42,70.38\}\) & 熊/基准/牛：25.37/\allowbreak45.42/\allowbreak70.38元；概率值：44.40元；空间：-6.4\%。 & 质量基础高；事实输入：官方预告、Q1财务与行情；外部背景：东莞证券（2026-05-22）。监控：验证：金属价格、产量、单位成本、库存和经营现金流；反证：价格回落、成本曲线抬升、产量不达预期或利润无法转化为现金 \\
002738 中矿资源 & 周期变量主导，禁止直接年化H1利润 & H1扣非EPS 1.5246元；H2/H1=0.55/0.85/1.15；\(\mathrm{EPS}_s=1.5246\times(1+\{0.55,0.85,1.15\});\ V_s=\mathrm{EPS}_s\times\{8.0,12.0,16.0\}\Rightarrow\{18.91,46.20,52.45\}\) & 熊/基准/牛：18.91/\allowbreak46.20/\allowbreak52.45元；概率值：39.26元；空间：-20.4\%。 & 质量基础高；事实输入：官方预告、Q1财务与行情；外部背景：东吴证券（2026-07-14）。监控：验证：金属价格、产量、单位成本、库存和经营现金流；反证：价格回落、成本曲线抬升、产量不达预期或利润无法转化为现金 \\
601069 西部黄金 & 周期变量主导，禁止直接年化H1利润 & H1扣非EPS 0.5873元；H2/H1=0.55/0.85/1.15；\(\mathrm{EPS}_s=0.5873\times(1+\{0.55,0.85,1.15\});\ V_s=\mathrm{EPS}_s\times\{8.0,12.0,16.0\}\Rightarrow\{7.28,13.04,20.20\}\) & 熊/基准/牛：7.28/\allowbreak13.04/\allowbreak20.20元；概率值：12.74元；空间：-44.5\%。 & 质量基础高；事实输入：官方预告、Q1财务与行情；外部背景：民生证券（2025-10-31）。监控：验证：金属价格、产量、单位成本、库存和经营现金流；反证：价格回落、成本曲线抬升、产量不达预期或利润无法转化为现金 \\
600026 中远海能 & 周期变量主导，禁止直接年化H1利润 & H1扣非EPS 0.8042元；H2/H1=0.55/0.85/1.15；\(\mathrm{EPS}_s=0.8042\times(1+\{0.55,0.85,1.15\});\ V_s=\mathrm{EPS}_s\times\{8.0,12.0,16.0\}\Rightarrow\{9.97,17.85,27.66\}\) & 熊/基准/牛：9.97/\allowbreak17.85/\allowbreak27.66元；概率值：17.45元；空间：10.2\%。 & 质量基础高；事实输入：官方预告、Q1财务与行情；外部背景：国金证券（2025-10-31）。监控：验证：运价、运力利用率、燃料成本、资本开支和现金流；反证：运价回落、利用率下降、成本上升或周期盈利被错误年化 \\
000426 兴业银锡 & 周期变量主导，禁止直接年化H1利润 & H1扣非EPS 1.0165元；H2/H1=0.55/0.85/1.15；\(\mathrm{EPS}_s=1.0165\times(1+\{0.55,0.85,1.15\});\ V_s=\mathrm{EPS}_s\times\{8.0,12.0,16.0\}\Rightarrow\{12.61,25.56,34.97\}\) & 熊/基准/牛：12.61/\allowbreak25.56/\allowbreak34.97元；概率值：23.56元；空间：-22.9\%。 & 质量基础高；事实输入：官方预告、Q1财务与行情；外部背景：国信证券（2026-04-24）。监控：验证：金属价格、产量、单位成本、库存和经营现金流；反证：价格回落、成本曲线抬升、产量不达预期或利润无法转化为现金 \\
600489 中金黄金 & 周期变量主导，禁止直接年化H1利润 & H1扣非EPS 0.8871元；H2/H1=0.55/0.85/1.15；\(\mathrm{EPS}_s=0.8871\times(1+\{0.55,0.85,1.15\});\ V_s=\mathrm{EPS}_s\times\{8.0,12.0,16.0\}\Rightarrow\{11.00,23.52,30.52\}\) & 熊/基准/牛：11.00/\allowbreak23.52/\allowbreak30.52元；概率值：21.16元；空间：5.2\%。 & 质量基础高；事实输入：官方预告、Q1财务与行情；外部背景：中邮证券（2026-05-08）。监控：验证：金属价格、产量、单位成本、库存和经营现金流；反证：价格回落、成本曲线抬升、产量不达预期或利润无法转化为现金 \\
000737 北方铜业 & 周期变量主导，禁止直接年化H1利润 & H1扣非EPS 0.7099元；H2/H1=0.55/0.85/1.15；\(\mathrm{EPS}_s=0.7099\times(1+\{0.55,0.85,1.15\});\ V_s=\mathrm{EPS}_s\times\{8.0,12.0,16.0\}\Rightarrow\{8.80,15.76,24.42\}\) & 熊/基准/牛：8.80/\allowbreak15.76/\allowbreak24.42元；概率值：15.40元；空间：20.8\%。 & 质量中低；事实输入：官方预告、Q1财务与行情；外部背景：无当前券商目标锚。监控：验证：金属价格、产量、单位成本、库存和经营现金流；反证：价格回落、成本曲线抬升、产量不达预期或利润无法转化为现金 \\
002756 永兴材料 & 周期变量主导，禁止直接年化H1利润 & H1扣非EPS 1.8549元；H2/H1=0.55/0.85/1.15；\(\mathrm{EPS}_s=1.8549\times(1+\{0.55,0.85,1.15\});\ V_s=\mathrm{EPS}_s\times\{8.0,12.0,16.0\}\Rightarrow\{23.00,44.88,63.81\}\) & 熊/基准/牛：23.00/\allowbreak44.88/\allowbreak63.81元；概率值：42.10元；空间：-6.2\%。 & 质量基础高；事实输入：官方预告、Q1财务与行情；外部背景：东吴证券（2026-07-14）。监控：验证：金属价格、产量、单位成本、库存和经营现金流；反证：价格回落、成本曲线抬升、产量不达预期或利润无法转化为现金 \\
000975 山金国际 & 周期变量主导，禁止直接年化H1利润 & H1扣非EPS 0.8458元；H2/H1=0.55/0.85/1.15；\(\mathrm{EPS}_s=0.8458\times(1+\{0.55,0.85,1.15\});\ V_s=\mathrm{EPS}_s\times\{8.0,12.0,16.0\}\Rightarrow\{10.49,22.68,29.10\}\) & 熊/基准/牛：10.49/\allowbreak22.68/\allowbreak29.10元；概率值：20.31元；空间：5.0\%。 & 质量基础高；事实输入：官方预告、Q1财务与行情；外部背景：太平洋（2026-05-13）。监控：验证：金属价格、产量、单位成本、库存和经营现金流；反证：价格回落、成本曲线抬升、产量不达预期或利润无法转化为现金 \\
000703 恒逸石化 & 周期变量主导，禁止直接年化H1利润 & H1扣非EPS 1.4968元；H2/H1=0.55/0.90/1.25；\(\mathrm{EPS}_s=1.4968\times(1+\{0.55,0.90,1.25\});\ V_s=\mathrm{EPS}_s\times\{6.0,9.0,12.0\}\Rightarrow\{10.86,25.59,40.41\}\) & 熊/基准/牛：10.86/\allowbreak25.59/\allowbreak40.41元；概率值：23.42元；空间：61.7\%。 & 质量基础高；事实输入：官方预告、Q1财务与行情；外部背景：国信证券（2026-07-05）。监控：验证：产品价差、原料成本、装置运行率、库存和现金流；反证：价差收窄、原料成本上升、装置检修或盈利恢复无法持续 \\
000301 东方盛虹 & 周期变量主导，禁止直接年化H1利润 & H1扣非EPS 0.6680元；H2/H1=0.55/0.90/1.25；\(\mathrm{EPS}_s=0.6680\times(1+\{0.55,0.90,1.25\});\ V_s=\mathrm{EPS}_s\times\{6.0,9.0,12.0\}\Rightarrow\{6.21,11.42,18.03\}\) & 熊/基准/牛：6.21/\allowbreak11.42/\allowbreak18.03元；概率值：11.18元；空间：-3.7\%。 & 质量基础高；事实输入：官方预告、Q1财务与行情；外部背景：中邮证券（2026-05-06）。监控：验证：产品价差、原料成本、装置运行率、库存和现金流；反证：价差收窄、原料成本上升、装置检修或盈利恢复无法持续 \\
002493 荣盛石化 & 周期变量主导，禁止直接年化H1利润 & H1扣非EPS 0.5306元；H2/H1=0.55/0.90/1.25；\(\mathrm{EPS}_s=0.5306\times(1+\{0.55,0.90,1.25\});\ V_s=\mathrm{EPS}_s\times\{6.0,9.0,12.0\}\Rightarrow\{4.93,9.07,14.33\}\) & 熊/基准/牛：4.93/\allowbreak9.07/\allowbreak14.33元；概率值：8.88元；空间：-20.0\%。 & 质量基础高；事实输入：官方预告、Q1财务与行情；外部背景：国信证券（2026-07-15）。监控：验证：产品价差、原料成本、装置运行率、库存和现金流；反证：价差收窄、原料成本上升、装置检修或盈利恢复无法持续 \\
601233 桐昆股份 & 周期变量主导，禁止直接年化H1利润 & H1扣非EPS 1.7150元；H2/H1=0.55/0.90/1.25；\(\mathrm{EPS}_s=1.7150\times(1+\{0.55,0.90,1.25\});\ V_s=\mathrm{EPS}_s\times\{6.0,9.0,12.0\}\Rightarrow\{15.54,29.33,46.31\}\) & 熊/基准/牛：15.54/\allowbreak29.33/\allowbreak46.31元；概率值：28.59元；空间：38.0\%。 & 质量基础高；事实输入：官方预告、Q1财务与行情；外部背景：东海证券（2026-05-08）。监控：验证：产品价差、原料成本、装置运行率、库存和现金流；反证：价差收窄、原料成本上升、装置检修或盈利恢复无法持续 \\
601225 陕西煤业 & 周期变量主导，禁止直接年化H1利润 & H1扣非EPS 1.0053元；H2/H1=0.55/0.90/1.25；\(\mathrm{EPS}_s=1.0053\times(1+\{0.55,0.90,1.25\});\ V_s=\mathrm{EPS}_s\times\{7.0,9.0,11.0\}\Rightarrow\{10.91,18.72,24.88\}\) & 熊/基准/牛：10.91/\allowbreak18.72/\allowbreak24.88元；概率值：17.61元；空间：-28.4\%。 & 质量基础高；事实输入：官方预告、Q1财务与行情；外部背景：国信证券（2026-04-28）。监控：验证：收入/利润兑现、现金流、行业价格和估值分母；反证：利润兑现不及预期、现金流背离或当前估值缺少安全边际 \\
000807 云铝股份 & 周期变量主导，禁止直接年化H1利润 & H1扣非EPS 2.1771元；H2/H1=0.55/0.85/1.15；\(\mathrm{EPS}_s=2.1771\times(1+\{0.55,0.85,1.15\});\ V_s=\mathrm{EPS}_s\times\{8.0,12.0,16.0\}\Rightarrow\{19.58,48.33,74.89\}\) & 熊/基准/牛：19.58/\allowbreak48.33/\allowbreak74.89元；概率值：45.02元；空间：72.5\%。 & 质量基础高；事实输入：官方预告、Q1财务与行情；外部背景：东吴证券（2026-04-27）。监控：验证：金属价格、产量、单位成本、库存和经营现金流；反证：价格回落、成本曲线抬升、产量不达预期或利润无法转化为现金 \\
603993 洛阳钼业 & 周期变量主导，禁止直接年化H1利润 & H1扣非EPS 0.7245元；H2/H1=0.55/0.85/1.15；\(\mathrm{EPS}_s=0.7245\times(1+\{0.55,0.85,1.15\});\ V_s=\mathrm{EPS}_s\times\{8.0,12.0,16.0\}\Rightarrow\{8.98,17.52,24.92\}\) & 熊/基准/牛：8.98/\allowbreak17.52/\allowbreak24.92元；概率值：16.44元；空间：-8.2\%。 & 质量基础高；事实输入：官方预告、Q1财务与行情；外部背景：太平洋（2026-05-14）。监控：验证：金属价格、产量、单位成本、库存和经营现金流；反证：价格回落、成本曲线抬升、产量不达预期或利润无法转化为现金 \\
601872 招商轮船 & 周期变量主导，禁止直接年化H1利润 & H1扣非EPS 0.8533元；H2/H1=0.55/0.85/1.15；\(\mathrm{EPS}_s=0.8533\times(1+\{0.55,0.85,1.15\});\ V_s=\mathrm{EPS}_s\times\{8.0,12.0,16.0\}\Rightarrow\{10.58,25.08,29.35\}\) & 熊/基准/牛：10.58/\allowbreak25.08/\allowbreak29.35元；概率值：21.58元；空间：41.4\%。 & 质量基础高；事实输入：官方预告、Q1财务与行情；外部背景：华源证券（2026-07-06）。监控：验证：运价、运力利用率、燃料成本、资本开支和现金流；反证：运价回落、利用率下降、成本上升或周期盈利被错误年化 \\
601899 紫金矿业 & 周期变量主导，禁止直接年化H1利润 & H1扣非EPS 1.4253元；H2/H1=0.55/0.85/1.15；\(\mathrm{EPS}_s=1.4253\times(1+\{0.55,0.85,1.15\});\ V_s=\mathrm{EPS}_s\times\{8.0,12.0,16.0\}\Rightarrow\{17.67,40.08,49.03\}\) & 熊/基准/牛：17.67/\allowbreak40.08/\allowbreak49.03元；概率值：35.15元；空间：21.9\%。 & 质量基础高；事实输入：官方预告、Q1财务与行情；外部背景：中邮证券（2026-04-28）。监控：验证：金属价格、产量、单位成本、库存和经营现金流；反证：价格回落、成本曲线抬升、产量不达预期或利润无法转化为现金 \\
600309 万华化学 & 周期变量主导，禁止直接年化H1利润 & H1扣非EPS 3.1625元；H2/H1=0.55/0.85/1.15；\(\mathrm{EPS}_s=3.1625\times(1+\{0.55,0.85,1.15\});\ V_s=\mathrm{EPS}_s\times\{10.0,14.0,18.0\}\Rightarrow\{49.02,81.91,122.39\}\) & 熊/基准/牛：49.02/\allowbreak81.91/\allowbreak122.39元；概率值：80.14元；空间：14.3\%。 & 质量基础高；事实输入：官方预告、Q1财务与行情；外部背景：国信证券（2026-07-07）。监控：验证：收入/利润兑现、现金流、行业价格和估值分母；反证：利润兑现不及预期、现金流背离或当前估值缺少安全边际 \\
600989 宝丰能源 & 周期变量主导，禁止直接年化H1利润 & H1扣非EPS 1.2477元；H2/H1=0.55/0.85/1.15；\(\mathrm{EPS}_s=1.2477\times(1+\{0.55,0.85,1.15\});\ V_s=\mathrm{EPS}_s\times\{10.0,14.0,18.0\}\Rightarrow\{16.55,32.32,48.29\}\) & 熊/基准/牛：16.55/\allowbreak32.32/\allowbreak48.29元；概率值：31.66元；空间：43.5\%。 & 质量基础高；事实输入：官方预告、Q1财务与行情；外部背景：环球富盛理财（2026-05-05）。监控：验证：收入/利润兑现、现金流、行业价格和估值分母；反证：利润兑现不及预期、现金流背离或当前估值缺少安全边际 \\
601088 中国神华 & 周期变量主导，禁止直接年化H1利润 & H1扣非EPS 1.2495元；H2/H1=0.55/0.90/1.25；\(\mathrm{EPS}_s=1.2495\times(1+\{0.55,0.90,1.25\});\ V_s=\mathrm{EPS}_s\times\{7.0,9.0,11.0\}\Rightarrow\{13.56,22.86,30.92\}\) & 熊/基准/牛：13.56/\allowbreak22.86/\allowbreak30.92元；概率值：21.68元；空间：-50.4\%。 & 质量基础高；事实输入：官方预告、Q1财务与行情；外部背景：华源证券（2026-04-11）。监控：验证：收入/利润兑现、现金流、行业价格和估值分母；反证：利润兑现不及预期、现金流背离或当前估值缺少安全边际 \\
600362 江西铜业 & 周期变量主导，禁止直接年化H1利润 & H1扣非EPS 2.3406元；H2/H1=0.55/0.85/1.15；\(\mathrm{EPS}_s=2.3406\times(1+\{0.55,0.85,1.15\});\ V_s=\mathrm{EPS}_s\times\{8.0,12.0,16.0\}\Rightarrow\{29.02,51.96,80.52\}\) & 熊/基准/牛：29.02/\allowbreak51.96/\allowbreak80.52元；概率值：50.79元；空间：20.1\%。 & 质量基础高；事实输入：官方预告、Q1财务与行情；外部背景：国信证券（2025-09-02）。监控：验证：金属价格、产量、单位成本、库存和经营现金流；反证：价格回落、成本曲线抬升、产量不达预期或利润无法转化为现金 \\
002648 卫星化学 & 周期变量主导，禁止直接年化H1利润 & H1扣非EPS 1.8367元；H2/H1=0.55/0.85/1.15；\(\mathrm{EPS}_s=1.8367\times(1+\{0.55,0.85,1.15\});\ V_s=\mathrm{EPS}_s\times\{10.0,14.0,18.0\}\Rightarrow\{17.57,47.57,71.08\}\) & 熊/基准/牛：17.57/\allowbreak47.57/\allowbreak71.08元；概率值：42.38元；空间：80.9\%。 & 质量基础高；事实输入：官方预告、Q1财务与行情；外部背景：西南证券（2026-04-24）。监控：验证：收入/利润兑现、现金流、行业价格和估值分母；反证：利润兑现不及预期、现金流背离或当前估值缺少安全边际 \\
000792 盐湖股份 & 周期变量主导，禁止直接年化H1利润 & H1扣非EPS 1.1622元；H2/H1=0.55/0.85/1.15；\(\mathrm{EPS}_s=1.1622\times(1+\{0.55,0.85,1.15\});\ V_s=\mathrm{EPS}_s\times\{10.0,14.0,18.0\}\Rightarrow\{18.01,30.10,44.98\}\) & 熊/基准/牛：18.01/\allowbreak30.10/\allowbreak44.98元；概率值：29.45元；空间：16.5\%。 & 质量基础高；事实输入：官方预告、Q1财务与行情；外部背景：华鑫证券（2026-04-09）。监控：验证：收入/利润兑现、现金流、行业价格和估值分母；反证：利润兑现不及预期、现金流背离或当前估值缺少安全边际 \\
000657 中钨高新 & 周期变量主导，禁止直接年化H1利润 & H1扣非EPS 0.9019元；H2/H1=0.55/0.85/1.15；\(\mathrm{EPS}_s=0.9019\times(1+\{0.55,0.85,1.15\});\ V_s=\mathrm{EPS}_s\times\{8.0,12.0,16.0\}\Rightarrow\{11.18,27.00,31.02\}\) & 熊/基准/牛：11.18/\allowbreak27.00/\allowbreak31.02元；概率值：23.06元；空间：-67.0\%。 & 质量基础高；事实输入：官方预告、Q1财务与行情；外部背景：东吴证券（2026-07-14）。监控：验证：金属价格、产量、单位成本、库存和经营现金流；反证：价格回落、成本曲线抬升、产量不达预期或利润无法转化为现金 \\
002497 雅化集团 & 周期变量主导，禁止直接年化H1利润 & H1扣非EPS 1.0585元；H2/H1=0.55/0.85/1.15；\(\mathrm{EPS}_s=1.0585\times(1+\{0.55,0.85,1.15\});\ V_s=\mathrm{EPS}_s\times\{10.0,14.0,18.0\}\Rightarrow\{13.55,36.82,40.96\}\) & 熊/基准/牛：13.55/\allowbreak36.82/\allowbreak40.96元；概率值：31.80元；空间：76.0\%。 & 质量基础高；事实输入：官方预告、Q1财务与行情；外部背景：东吴证券（2026-07-08）。监控：验证：收入/利润兑现、现金流、行业价格和估值分母；反证：利润兑现不及预期、现金流背离或当前估值缺少安全边际 \\
000933 神火股份 & 周期变量主导，禁止直接年化H1利润 & H1扣非EPS 2.1369元；H2/H1=0.55/0.85/1.15；\(\mathrm{EPS}_s=2.1369\times(1+\{0.55,0.85,1.15\});\ V_s=\mathrm{EPS}_s\times\{8.0,12.0,16.0\}\Rightarrow\{19.13,47.54,73.51\}\) & 熊/基准/牛：19.13/\allowbreak47.54/\allowbreak73.51元；概率值：44.21元；空间：73.3\%。 & 质量基础高；事实输入：官方预告、Q1财务与行情；外部背景：国金证券（2026-04-14）。监控：验证：金属价格、产量、单位成本、库存和经营现金流；反证：价格回落、成本曲线抬升、产量不达预期或利润无法转化为现金 \\
002192 融捷股份 & 周期变量主导，禁止直接年化H1利润 & H1扣非EPS 3.8629元；H2/H1=0.55/0.85/1.15；\(\mathrm{EPS}_s=3.8629\times(1+\{0.55,0.85,1.15\});\ V_s=\mathrm{EPS}_s\times\{8.0,12.0,16.0\}\Rightarrow\{47.90,85.76,132.88\}\) & 熊/基准/牛：47.90/\allowbreak85.76/\allowbreak132.88元；概率值：83.83元；空间：26.2\%。 & 质量基础高；事实输入：官方预告、Q1财务与行情；外部背景：东吴证券（2022-10-26）。监控：验证：金属价格、产量、单位成本、库存和经营现金流；反证：价格回落、成本曲线抬升、产量不达预期或利润无法转化为现金 \\
002378 章源钨业 & 周期变量主导，禁止直接年化H1利润 & H1扣非EPS 0.5660元；H2/H1=0.55/0.85/1.15；\(\mathrm{EPS}_s=0.5660\times(1+\{0.55,0.85,1.15\});\ V_s=\mathrm{EPS}_s\times\{8.0,12.0,16.0\}\Rightarrow\{7.02,12.57,19.47\}\) & 熊/基准/牛：7.02/\allowbreak12.57/\allowbreak19.47元；概率值：12.29元；空间：-52.4\%。 & 质量基础高；事实输入：官方预告、Q1财务与行情；外部背景：民生证券（2025-08-27）。监控：验证：金属价格、产量、单位成本、库存和经营现金流；反证：价格回落、成本曲线抬升、产量不达预期或利润无法转化为现金 \\
002842 翔鹭钨业 & 周期变量主导，禁止直接年化H1利润 & H1扣非EPS 1.6352元；H2/H1=0.55/0.85/1.15；\(\mathrm{EPS}_s=1.6352\times(1+\{0.55,0.85,1.15\});\ V_s=\mathrm{EPS}_s\times\{8.0,12.0,16.0\}\Rightarrow\{20.28,36.30,56.25\}\) & 熊/基准/牛：20.28/\allowbreak36.30/\allowbreak56.25元；概率值：35.48元；空间：-1.3\%。 & 质量基础高；事实输入：官方预告、Q1财务与行情；外部背景：中邮证券（2018-12-04）。监控：验证：金属价格、产量、单位成本、库存和经营现金流；反证：价格回落、成本曲线抬升、产量不达预期或利润无法转化为现金 \\
002407 多氟多 & 周期变量主导，禁止直接年化H1利润 & H1扣非EPS 0.4106元；H2/H1=0.55/0.85/1.15；\(\mathrm{EPS}_s=0.4106\times(1+\{0.55,0.85,1.15\});\ V_s=\mathrm{EPS}_s\times\{10.0,14.0,18.0\}\Rightarrow\{6.37,10.64,15.89\}\) & 熊/基准/牛：6.37/\allowbreak10.64/\allowbreak15.89元；概率值：10.41元；空间：-69.3\%。 & 质量基础高；事实输入：官方预告、Q1财务与行情；外部背景：民生证券（2025-04-24）。监控：验证：收入/利润兑现、现金流、行业价格和估值分母；反证：利润兑现不及预期、现金流背离或当前估值缺少安全边际 \\
600961 株冶集团 & 周期变量主导，禁止直接年化H1利润 & H1扣非EPS 1.7803元；H2/H1=0.55/0.85/1.15；\(\mathrm{EPS}_s=1.7803\times(1+\{0.55,0.85,1.15\});\ V_s=\mathrm{EPS}_s\times\{8.0,12.0,16.0\}\Rightarrow\{18.57,39.52,61.24\}\) & 熊/基准/牛：18.57/\allowbreak39.52/\allowbreak61.24元；概率值：37.58元；空间：51.8\%。 & 质量基础高；事实输入：官方预告、Q1财务与行情；外部背景：国信证券（2026-04-21）。监控：验证：金属价格、产量、单位成本、库存和经营现金流；反证：价格回落、成本曲线抬升、产量不达预期或利润无法转化为现金 \\
002128 电投能源 & 周期变量主导，禁止直接年化H1利润 & H1扣非EPS 1.3701元；H2/H1=0.55/0.90/1.25；\(\mathrm{EPS}_s=1.3701\times(1+\{0.55,0.90,1.25\});\ V_s=\mathrm{EPS}_s\times\{7.0,9.0,11.0\}\Rightarrow\{14.87,25.47,33.91\}\) & 熊/基准/牛：14.87/\allowbreak25.47/\allowbreak33.91元；概率值：23.98元；空间：-9.5\%。 & 质量基础高；事实输入：官方预告、Q1财务与行情；外部背景：开源证券（2026-04-17）。监控：验证：收入/利润兑现、现金流、行业价格和估值分母；反证：利润兑现不及预期、现金流背离或当前估值缺少安全边际 \\
601168 西部矿业 & 周期变量主导，禁止直接年化H1利润 & H1扣非EPS 1.7583元；H2/H1=0.55/0.85/1.15；\(\mathrm{EPS}_s=1.7583\times(1+\{0.55,0.85,1.15\});\ V_s=\mathrm{EPS}_s\times\{8.0,12.0,16.0\}\Rightarrow\{21.80,39.03,60.49\}\) & 熊/基准/牛：21.80/\allowbreak39.03/\allowbreak60.49元；概率值：38.19元；空间：15.1\%。 & 质量基础高；事实输入：官方预告、Q1财务与行情；外部背景：国信证券（2026-07-06）。监控：验证：金属价格、产量、单位成本、库存和经营现金流；反证：价格回落、成本曲线抬升、产量不达预期或利润无法转化为现金 \\
600188 兖矿能源 & 周期变量主导，禁止直接年化H1利润 & H1扣非EPS 0.4483元；H2/H1=0.55/0.90/1.25；\(\mathrm{EPS}_s=0.4483\times(1+\{0.55,0.90,1.25\});\ V_s=\mathrm{EPS}_s\times\{7.0,9.0,11.0\}\Rightarrow\{4.86,11.10,14.76\}\) & 熊/基准/牛：4.86/\allowbreak11.10/\allowbreak14.76元；概率值：9.96元；空间：-50.8\%。 & 质量基础高；事实输入：官方预告、Q1财务与行情；外部背景：开源证券（2026-05-05）。监控：验证：收入/利润兑现、现金流、行业价格和估值分母；反证：利润兑现不及预期、现金流背离或当前估值缺少安全边际 \\
603162 海通发展 & 周期变量主导，禁止直接年化H1利润 & H1扣非EPS 0.3802元；H2/H1=0.55/0.85/1.15；\(\mathrm{EPS}_s=0.3802\times(1+\{0.55,0.85,1.15\});\ V_s=\mathrm{EPS}_s\times\{8.0,12.0,16.0\}\Rightarrow\{4.71,9.00,13.08\}\) & 熊/基准/牛：4.71/\allowbreak9.00/\allowbreak13.08元；概率值：8.53元；空间：-24.1\%。 & 质量基础高；事实输入：官方预告、Q1财务与行情；外部背景：西南证券（2026-07-09）。监控：验证：运价、运力利用率、燃料成本、资本开支和现金流；反证：运价回落、利用率下降、成本上升或周期盈利被错误年化 \\
002440 闰土股份 & 周期变量主导，禁止直接年化H1利润 & H1扣非EPS 0.3915元；H2/H1=0.55/0.85/1.15；\(\mathrm{EPS}_s=0.3915\times(1+\{0.55,0.85,1.15\});\ V_s=\mathrm{EPS}_s\times\{10.0,14.0,18.0\}\Rightarrow\{6.07,10.14,15.15\}\) & 熊/基准/牛：6.07/\allowbreak10.14/\allowbreak15.15元；概率值：9.92元；空间：-12.1\%。 & 质量基础高；事实输入：官方预告、Q1财务与行情；外部背景：华鑫证券（2025-05-02）。监控：验证：收入/利润兑现、现金流、行业价格和估值分母；反证：利润兑现不及预期、现金流背离或当前估值缺少安全边际 \\
002532 天山铝业 & 周期变量主导，禁止直接年化H1利润 & H1扣非EPS 0.8836元；H2/H1=0.55/0.85/1.15；\(\mathrm{EPS}_s=0.8836\times(1+\{0.55,0.85,1.15\});\ V_s=\mathrm{EPS}_s\times\{8.0,12.0,16.0\}\Rightarrow\{9.56,23.03,30.40\}\) & 熊/基准/牛：9.56/\allowbreak23.03/\allowbreak30.40元；概率值：20.46元；空间：60.5\%。 & 质量基础高；事实输入：官方预告、Q1财务与行情；外部背景：国金证券（2026-04-29）。监控：验证：金属价格、产量、单位成本、库存和经营现金流；反证：价格回落、成本曲线抬升、产量不达预期或利润无法转化为现金 \\
000408 藏格矿业 & 周期变量主导，禁止直接年化H1利润 & H1扣非EPS 2.3666元；H2/H1=0.55/0.85/1.15；\(\mathrm{EPS}_s=2.3666\times(1+\{0.55,0.85,1.15\});\ V_s=\mathrm{EPS}_s\times\{10.0,14.0,18.0\}\Rightarrow\{36.68,71.68,91.59\}\) & 熊/基准/牛：36.68/\allowbreak71.68/\allowbreak91.59元；概率值：71.44元；空间：-0.1\%。 & 质量基础高；事实输入：官方预告、Q1财务与行情；外部背景：东吴证券（2026-04-18）。监控：验证：收入/利润兑现、现金流、行业价格和估值分母；反证：利润兑现不及预期、现金流背离或当前估值缺少安全边际 \\
600595 中孚实业 & 周期变量主导，禁止直接年化H1利润 & H1扣非EPS 0.4803元；H2/H1=0.55/0.85/1.15；\(\mathrm{EPS}_s=0.4803\times(1+\{0.55,0.85,1.15\});\ V_s=\mathrm{EPS}_s\times\{8.0,12.0,16.0\}\Rightarrow\{4.84,12.00,16.52\}\) & 熊/基准/牛：4.84/\allowbreak12.00/\allowbreak16.52元；概率值：10.76元；空间：66.6\%。 & 质量基础高；事实输入：官方预告、Q1财务与行情；外部背景：国信证券（2026-04-28）。监控：验证：金属价格、产量、单位成本、库存和经营现金流；反证：价格回落、成本曲线抬升、产量不达预期或利润无法转化为现金 \\
000630 铜陵有色 & 周期变量主导，禁止直接年化H1利润 & H1扣非EPS 0.2125元；H2/H1=0.55/0.85/1.15；\(\mathrm{EPS}_s=0.2125\times(1+\{0.55,0.85,1.15\});\ V_s=\mathrm{EPS}_s\times\{8.0,12.0,16.0\}\Rightarrow\{2.64,6.72,7.31\}\) & 熊/基准/牛：2.64/\allowbreak6.72/\allowbreak7.31元；概率值：5.61元；空间：-7.4\%。 & 质量基础高；事实输入：官方预告、Q1财务与行情；外部背景：国信证券（2026-04-22）。监控：验证：金属价格、产量、单位成本、库存和经营现金流；反证：价格回落、成本曲线抬升、产量不达预期或利润无法转化为现金 \\
600549 厦门钨业 & 周期变量主导，禁止直接年化H1利润 & H1扣非EPS 1.3707元；H2/H1=0.55/0.85/1.15；\(\mathrm{EPS}_s=1.3707\times(1+\{0.55,0.85,1.15\});\ V_s=\mathrm{EPS}_s\times\{8.0,12.0,16.0\}\Rightarrow\{17.00,35.16,47.15\}\) & 熊/基准/牛：17.00/\allowbreak35.16/\allowbreak47.15元；概率值：32.11元；空间：-44.7\%。 & 质量基础高；事实输入：官方预告、Q1财务与行情；外部背景：太平洋（2026-06-21）。监控：验证：金属价格、产量、单位成本、库存和经营现金流；反证：价格回落、成本曲线抬升、产量不达预期或利润无法转化为现金 \\
600111 北方稀土 & 周期变量主导，禁止直接年化H1利润 & H1扣非EPS 0.5615元；H2/H1=0.55/0.85/1.15；\(\mathrm{EPS}_s=0.5615\times(1+\{0.55,0.85,1.15\});\ V_s=\mathrm{EPS}_s\times\{8.0,12.0,16.0\}\Rightarrow\{6.96,13.44,19.32\}\) & 熊/基准/牛：6.96/\allowbreak13.44/\allowbreak19.32元；概率值：12.67元；空间：-68.0\%。 & 质量基础高；事实输入：官方预告、Q1财务与行情；外部背景：东莞证券（2026-04-23）。监控：验证：金属价格、产量、单位成本、库存和经营现金流；反证：价格回落、成本曲线抬升、产量不达预期或利润无法转化为现金 \\
000612 焦作万方 & 周期变量主导，禁止直接年化H1利润 & H1扣非EPS 1.0133元；H2/H1=0.55/0.85/1.15；\(\mathrm{EPS}_s=1.0133\times(1+\{0.55,0.85,1.15\});\ V_s=\mathrm{EPS}_s\times\{8.0,12.0,16.0\}\Rightarrow\{8.13,22.50,34.86\}\) & 熊/基准/牛：8.13/\allowbreak22.50/\allowbreak34.86元；概率值：20.66元；空间：90.6\%。 & 质量基础高；事实输入：官方预告、Q1财务与行情；外部背景：开源证券（2025-09-26）。监控：验证：金属价格、产量、单位成本、库存和经营现金流；反证：价格回落、成本曲线抬升、产量不达预期或利润无法转化为现金 \\
600233 圆通速递 & 周期变量主导，禁止直接年化H1利润 & H1扣非EPS 0.9320元；H2/H1=0.55/0.85/1.15；\(\mathrm{EPS}_s=0.9320\times(1+\{0.55,0.85,1.15\});\ V_s=\mathrm{EPS}_s\times\{8.0,12.0,16.0\}\Rightarrow\{11.56,21.96,32.06\}\) & 熊/基准/牛：11.56/\allowbreak21.96/\allowbreak32.06元；概率值：20.67元；空间：18.1\%。 & 质量基础高；事实输入：官方预告、Q1财务与行情；外部背景：群益证券（2026-07-01）。监控：验证：运价、运力利用率、燃料成本、资本开支和现金流；反证：运价回落、利用率下降、成本上升或周期盈利被错误年化 \\
600392 盛和资源 & 周期变量主导，禁止直接年化H1利润 & H1扣非EPS 0.4878元；H2/H1=0.55/0.85/1.15；\(\mathrm{EPS}_s=0.4878\times(1+\{0.55,0.85,1.15\});\ V_s=\mathrm{EPS}_s\times\{8.0,12.0,16.0\}\Rightarrow\{6.05,10.83,16.78\}\) & 熊/基准/牛：6.05/\allowbreak10.83/\allowbreak16.78元；概率值：10.59元；空间：-55.6\%。 & 质量基础高；事实输入：官方预告、Q1财务与行情；外部背景：国元证券（2025-11-12）。监控：验证：金属价格、产量、单位成本、库存和经营现金流；反证：价格回落、成本曲线抬升、产量不达预期或利润无法转化为现金 \\
603225 新凤鸣 & 周期变量主导，禁止直接年化H1利润 & H1扣非EPS 0.9031元；H2/H1=0.55/0.85/1.15；\(\mathrm{EPS}_s=0.9031\times(1+\{0.55,0.85,1.15\});\ V_s=\mathrm{EPS}_s\times\{10.0,14.0,18.0\}\Rightarrow\{13.48,23.39,34.95\}\) & 熊/基准/牛：13.48/\allowbreak23.39/\allowbreak34.95元；概率值：22.73元；空间：26.5\%。 & 质量基础高；事实输入：官方预告、Q1财务与行情；外部背景：中邮证券（2026-07-10）。监控：验证：收入/利润兑现、现金流、行业价格和估值分母；反证：利润兑现不及预期、现金流背离或当前估值缺少安全边际 \\
000060 中金岭南 & 周期变量主导，禁止直接年化H1利润 & H1扣非EPS 0.2488元；H2/H1=0.55/0.85/1.15；\(\mathrm{EPS}_s=0.2488\times(1+\{0.55,0.85,1.15\});\ V_s=\mathrm{EPS}_s\times\{8.0,12.0,16.0\}\Rightarrow\{3.09,6.12,8.56\}\) & 熊/基准/牛：3.09/\allowbreak6.12/\allowbreak8.56元；概率值：5.70元；空间：-8.9\%。 & 质量基础高；事实输入：官方预告、Q1财务与行情；外部背景：中邮证券（2026-05-13）。监控：验证：金属价格、产量、单位成本、库存和经营现金流；反证：价格回落、成本曲线抬升、产量不达预期或利润无法转化为现金 \\
002468 申通快递 & 周期变量主导，禁止直接年化H1利润 & H1扣非EPS 0.6565元；H2/H1=0.55/0.85/1.15；\(\mathrm{EPS}_s=0.6565\times(1+\{0.55,0.85,1.15\});\ V_s=\mathrm{EPS}_s\times\{8.0,12.0,16.0\}\Rightarrow\{8.14,15.84,22.58\}\) & 熊/基准/牛：8.14/\allowbreak15.84/\allowbreak22.58元；概率值：14.88元；空间：-2.2\%。 & 质量基础高；事实输入：官方预告、Q1财务与行情；外部背景：群益证券（2026-04-15）。监控：验证：运价、运力利用率、燃料成本、资本开支和现金流；反证：运价回落、利用率下降、成本上升或周期盈利被错误年化 \\
002221 东华能源 & 周期变量主导，禁止直接年化H1利润 & H1扣非EPS 0.0603元；H2/H1=0.55/0.90/1.25；\(\mathrm{EPS}_s=0.0603\times(1+\{0.55,0.90,1.25\});\ V_s=\mathrm{EPS}_s\times\{6.0,9.0,12.0\}\Rightarrow\{0.56,1.03,1.63\}\) & 熊/基准/牛：0.56/\allowbreak1.03/\allowbreak1.63元；概率值：1.01元；空间：-81.4\%。 & 质量基础高；事实输入：官方预告、Q1财务与行情；外部背景：东吴证券（2024-12-30）。监控：验证：产品价差、原料成本、装置运行率、库存和现金流；反证：价差收窄、原料成本上升、装置检修或盈利恢复无法持续 \\
\bottomrule
\end{longtable}
\normalsize

\section{M2｜成长盈利：以情景EPS/PE估值}
适用于存在盈利持续期、产品升级或订单释放的成长行业。模型只对已进入盈利桥接的扣非利润赋予成长估值，不把泛化主题叙事直接转化为EPS。

\footnotesize
\setlength{\tabcolsep}{2.5pt}
\begin{longtable}{L{1.55cm}L{2.45cm}L{5.0cm}L{2.75cm}L{4.1cm}}
\toprule
\textbf{代码/标的} & \textbf{适用理由} & \textbf{分母与公式代入} & \textbf{情景与概率值} & \textbf{证据与验证} \\
\midrule
\endfirsthead
\toprule
\textbf{代码/标的} & \textbf{适用理由} & \textbf{分母与公式代入} & \textbf{情景与概率值} & \textbf{证据与验证} \\
\midrule
\endhead
600150 中国船舶 & 成长持续期需由盈利兑现与现金流验证 & H1扣非EPS 1.3155元；H2/H1=0.70/1.05/1.40；\(\mathrm{EPS}_s=1.3155\times(1+\{0.70,1.05,1.40\});\ V_s=\mathrm{EPS}_s\times\{25.0,35.0,45.0\}\Rightarrow\{25.73,94.39,142.07\}\) & 熊/基准/牛：25.73/\allowbreak94.39/\allowbreak142.07元；概率值：83.33元；空间：142.9\%。 & 质量基础高；事实输入：官方预告、Q1财务与行情；外部背景：华源证券（2026-07-14）。监控：验证：订单交付、合同负债、产品结构和应收账款回款；反证：订单确认延后、交付节奏放缓、应收账款继续扩张或估值透支交付 \\
002602 世纪华通 & 成长持续期需由盈利兑现与现金流验证 & H1扣非EPS 0.5854元；H2/H1=0.70/1.05/1.40；\(\mathrm{EPS}_s=0.5854\times(1+\{0.70,1.05,1.40\});\ V_s=\mathrm{EPS}_s\times\{10.0,14.0,18.0\}\Rightarrow\{9.95,16.94,25.29\}\) & 熊/基准/牛：9.95/\allowbreak16.94/\allowbreak25.29元；概率值：16.51元；空间：15.8\%。 & 质量基础高；事实输入：官方预告、Q1财务与行情；外部背景：开源证券（2026-06-30）。监控：验证：产品上线、流水/付费、内容成本和经营现金流；反证：产品或内容兑现延后、用户/付费不及预期或现金流不能跟随利润 \\
300014 亿纬锂能 & 成长持续期需由盈利兑现与现金流验证 & H1扣非EPS 1.1579元；H2/H1=0.70/1.00/1.30；\(\mathrm{EPS}_s=1.1579\times(1+\{0.70,1.00,1.30\});\ V_s=\mathrm{EPS}_s\times\{16.0,22.0,28.0\}\Rightarrow\{31.49,68.46,74.57\}\) & 熊/基准/牛：31.49/\allowbreak68.46/\allowbreak74.57元；概率值：60.43元；空间：15.6\%。 & 质量基础高；事实输入：官方预告、Q1财务与行情；外部背景：群益证券（2026-06-26）。监控：验证：出货量、订单、产品价格、产能利用率和营运资金；反证：订单转化慢、价格竞争、利用率不足或库存/应收继续占用现金 \\
600739 辽宁成大 & 成长持续期需由盈利兑现与现金流验证 & H1扣非EPS 1.0407元；H2/H1=0.70/1.00/1.30；\(\mathrm{EPS}_s=1.0407\times(1+\{0.70,1.00,1.30\});\ V_s=\mathrm{EPS}_s\times\{18.0,25.0,32.0\}\Rightarrow\{8.21,52.03,76.60\}\) & 熊/基准/牛：8.21/\allowbreak52.03/\allowbreak76.60元；概率值：43.80元；空间：300.0\%。 & 质量基础高；事实输入：官方预告、Q1财务与行情；外部背景：太平洋（2017-10-30）。监控：验证：产品放量、集采/价格、研发进度、许可收入和现金流；反证：产品放量不及预期、价格压力、研发里程碑延后或一次性收益回落 \\
000623 吉林敖东 & 成长持续期需由盈利兑现与现金流验证 & H1扣非EPS 1.8534元；H2/H1=0.70/1.00/1.30；\(\mathrm{EPS}_s=1.8534\times(1+\{0.70,1.00,1.30\});\ V_s=\mathrm{EPS}_s\times\{18.0,25.0,32.0\}\Rightarrow\{14.00,92.67,136.41\}\) & 熊/基准/牛：14.00/\allowbreak92.67/\allowbreak136.41元；概率值：77.82元；空间：317.0\%。 & 质量基础高；事实输入：官方预告、Q1财务与行情；外部背景：国元证券（2024-09-03）。监控：验证：产品放量、集采/价格、研发进度、许可收入和现金流；反证：产品放量不及预期、价格压力、研发里程碑延后或一次性收益回落 \\
001339 智微智能 & 成长持续期需由盈利兑现与现金流验证 & H1扣非EPS 1.1199元；H2/H1=0.70/1.05/1.40；\(\mathrm{EPS}_s=1.1199\times(1+\{0.70,1.05,1.40\});\ V_s=\mathrm{EPS}_s\times\{18.0,28.0,38.0\}\Rightarrow\{34.27,69.44,102.13\}\) & 熊/基准/牛：34.27/\allowbreak69.44/\allowbreak102.13元；概率值：65.43元；空间：-26.7\%。 & 质量基础高；事实输入：官方预告、Q1财务与行情；外部背景：中邮证券（2026-06-18）。监控：验证：合同订单、项目交付、回款、研发投入和利润率；反证：订单确认延后、回款周期拉长、研发费用率上升或估值继续压缩 \\
002558 巨人网络 & 成长持续期需由盈利兑现与现金流验证 & H1扣非EPS 1.2102元；H2/H1=0.70/1.05/1.40；\(\mathrm{EPS}_s=1.2102\times(1+\{0.70,1.05,1.40\});\ V_s=\mathrm{EPS}_s\times\{10.0,14.0,18.0\}\Rightarrow\{20.57,34.73,52.28\}\) & 熊/基准/牛：20.57/\allowbreak34.73/\allowbreak52.28元；概率值：33.99元；空间：15.0\%。 & 质量基础高；事实输入：官方预告、Q1财务与行情；外部背景：西南证券（2026-07-15）。监控：验证：产品上线、流水/付费、内容成本和经营现金流；反证：产品或内容兑现延后、用户/付费不及预期或现金流不能跟随利润 \\
301308 江波龙 & 成长持续期需由盈利兑现与现金流验证 & H1扣非EPS 23.0463元；H2/H1=0.70/1.00/1.30；\(\mathrm{EPS}_s=23.0463\times(1+\{0.70,1.00,1.30\});\ V_s=\mathrm{EPS}_s\times\{18.0,26.0,34.0\}\Rightarrow\{366.00,1198.41,1802.22\}\) & 熊/基准/牛：366.00/\allowbreak1198.41/\allowbreak1802.22元；概率值：1069.45元；空间：119.2\%。 & 质量基础高；事实输入：官方预告、Q1财务与行情；外部背景：国信证券（2026-05-07）。监控：验证：收入/利润兑现、现金流、行业价格和估值分母；反证：利润兑现不及预期、现金流背离或当前估值缺少安全边际 \\
603986 兆易创新 & 成长持续期需由盈利兑现与现金流验证 & H1扣非EPS 6.9113元；H2/H1=0.70/1.00/1.30；\(\mathrm{EPS}_s=6.9113\times(1+\{0.70,1.00,1.30\});\ V_s=\mathrm{EPS}_s\times\{18.0,26.0,34.0\}\Rightarrow\{211.49,359.39,540.47\}\) & 熊/基准/牛：211.49/\allowbreak359.39/\allowbreak540.47元；概率值：351.24元；空间：-38.6\%。 & 质量基础高；事实输入：官方预告、Q1财务与行情；外部背景：华鑫证券（2026-06-29）。监控：验证：收入/利润兑现、现金流、行业价格和估值分母；反证：利润兑现不及预期、现金流背离或当前估值缺少安全边际 \\
601138 工业富联 & 成长持续期需由盈利兑现与现金流验证 & H1扣非EPS 1.1691元；H2/H1=0.70/1.00/1.30；\(\mathrm{EPS}_s=1.1691\times(1+\{0.70,1.00,1.30\});\ V_s=\mathrm{EPS}_s\times\{18.0,26.0,34.0\}\Rightarrow\{35.77,79.56,91.42\}\) & 熊/基准/牛：35.77/\allowbreak79.56/\allowbreak91.42元；概率值：68.80元；空间：8.1\%。 & 质量基础高；事实输入：官方预告、Q1财务与行情；外部背景：金元证券（2026-06-18）。监控：验证：收入/利润兑现、现金流、行业价格和估值分母；反证：利润兑现不及预期、现金流背离或当前估值缺少安全边际 \\
300475 香农芯创 & 成长持续期需由盈利兑现与现金流验证 & H1扣非EPS 7.7522元；H2/H1=0.70/1.00/1.30；\(\mathrm{EPS}_s=7.7522\times(1+\{0.70,1.00,1.30\});\ V_s=\mathrm{EPS}_s\times\{18.0,26.0,34.0\}\Rightarrow\{165.90,403.12,606.22\}\) & 熊/基准/牛：165.90/\allowbreak403.12/\allowbreak606.22元；概率值：372.57元；空间：68.4\%。 & 质量基础高；事实输入：官方预告、Q1财务与行情；外部背景：华鑫证券（2026-03-22）。监控：验证：收入/利润兑现、现金流、行业价格和估值分母；反证：利润兑现不及预期、现金流背离或当前估值缺少安全边际 \\
002709 天赐材料 & 成长持续期需由盈利兑现与现金流验证 & H1扣非EPS 1.3735元；H2/H1=0.70/1.00/1.30；\(\mathrm{EPS}_s=1.3735\times(1+\{0.70,1.00,1.30\});\ V_s=\mathrm{EPS}_s\times\{16.0,22.0,28.0\}\Rightarrow\{28.16,67.32,88.45\}\) & 熊/基准/牛：28.16/\allowbreak67.32/\allowbreak88.45元；概率值：60.46元；空间：61.0\%。 & 质量基础高；事实输入：官方预告、Q1财务与行情；外部背景：东吴证券（2026-07-10）。监控：验证：出货量、订单、产品价格、产能利用率和营运资金；反证：订单转化慢、价格竞争、利用率不足或库存/应收继续占用现金 \\
002414 高德红外 & 成长持续期需由盈利兑现与现金流验证 & H1扣非EPS 0.3103元；H2/H1=0.70/1.05/1.40；\(\mathrm{EPS}_s=0.3103\times(1+\{0.70,1.05,1.40\});\ V_s=\mathrm{EPS}_s\times\{25.0,35.0,45.0\}\Rightarrow\{9.99,22.26,33.51\}\) & 熊/基准/牛：9.99/\allowbreak22.26/\allowbreak33.51元；概率值：20.83元；空间：56.4\%。 & 质量基础高；事实输入：官方预告、Q1财务与行情；外部背景：太平洋（2025-11-06）。监控：验证：订单交付、合同负债、产品结构和应收账款回款；反证：订单确认延后、交付节奏放缓、应收账款继续扩张或估值透支交付 \\
601869 长飞光纤 & 成长持续期需由盈利兑现与现金流验证 & H1扣非EPS 2.7781元；H2/H1=0.70/1.00/1.30；\(\mathrm{EPS}_s=2.7781\times(1+\{0.70,1.00,1.30\});\ V_s=\mathrm{EPS}_s\times\{18.0,26.0,34.0\}\Rightarrow\{85.01,157.30,217.25\}\) & 熊/基准/牛：85.01/\allowbreak157.30/\allowbreak217.25元；概率值：147.60元；空间：-65.3\%。 & 质量基础高；事实输入：官方预告、Q1财务与行情；外部背景：华鑫证券（2026-04-12）。监控：验证：收入/利润兑现、现金流、行业价格和估值分母；反证：利润兑现不及预期、现金流背离或当前估值缺少安全边际 \\
002432 九安医疗 & 成长持续期需由盈利兑现与现金流验证 & H1扣非EPS 6.7612元；H2/H1=0.70/1.00/1.30；\(\mathrm{EPS}_s=6.7612\times(1+\{0.70,1.00,1.30\});\ V_s=\mathrm{EPS}_s\times\{18.0,25.0,32.0\}\Rightarrow\{50.24,338.06,497.63\}\) & 熊/基准/牛：50.24/\allowbreak338.06/\allowbreak497.63元；概率值：283.63元；空间：323.4\%。 & 质量基础高；事实输入：官方预告、Q1财务与行情；外部背景：东吴证券（2025-05-06）。监控：验证：产品放量、集采/价格、研发进度、许可收入和现金流；反证：产品放量不及预期、价格压力、研发里程碑延后或一次性收益回落 \\
300223 北京君正 & 成长持续期需由盈利兑现与现金流验证 & H1扣非EPS 2.3797元；H2/H1=0.70/1.00/1.30；\(\mathrm{EPS}_s=2.3797\times(1+\{0.70,1.00,1.30\});\ V_s=\mathrm{EPS}_s\times\{18.0,26.0,34.0\}\Rightarrow\{72.82,123.75,186.09\}\) & 熊/基准/牛：72.82/\allowbreak123.75/\allowbreak186.09元；概率值：120.94元；空间：-36.4\%。 & 质量基础高；事实输入：官方预告、Q1财务与行情；外部背景：国信证券（2026-05-06）。监控：验证：收入/利润兑现、现金流、行业价格和估值分母；反证：利润兑现不及预期、现金流背离或当前估值缺少安全边际 \\
002636 金安国纪 & 成长持续期需由盈利兑现与现金流验证 & H1扣非EPS 1.0721元；H2/H1=0.70/1.00/1.30；\(\mathrm{EPS}_s=1.0721\times(1+\{0.70,1.00,1.30\});\ V_s=\mathrm{EPS}_s\times\{18.0,26.0,34.0\}\Rightarrow\{32.81,55.75,83.84\}\) & 熊/基准/牛：32.81/\allowbreak55.75/\allowbreak83.84元；概率值：54.49元；空间：-42.5\%。 & 质量基础高；事实输入：官方预告、Q1财务与行情；外部背景：开源证券（2021-09-02）。监控：验证：收入/利润兑现、现金流、行业价格和估值分母；反证：利润兑现不及预期、现金流背离或当前估值缺少安全边际 \\
600884 杉杉股份 & 成长持续期需由盈利兑现与现金流验证 & H1扣非EPS 0.3445元；H2/H1=0.70/1.00/1.30；\(\mathrm{EPS}_s=0.3445\times(1+\{0.70,1.00,1.30\});\ V_s=\mathrm{EPS}_s\times\{16.0,22.0,28.0\}\Rightarrow\{9.11,15.16,22.19\}\) & 熊/基准/牛：9.11/\allowbreak15.16/\allowbreak22.19元；概率值：14.75元；空间：21.5\%。 & 质量基础高；事实输入：官方预告、Q1财务与行情；外部背景：民生证券（2024-04-28）。监控：验证：出货量、订单、产品价格、产能利用率和营运资金；反证：订单转化慢、价格竞争、利用率不足或库存/应收继续占用现金 \\
002653 海思科 & 成长持续期需由盈利兑现与现金流验证 & H1扣非EPS 0.6518元；H2/H1=0.70/1.00/1.30；\(\mathrm{EPS}_s=0.6518\times(1+\{0.70,1.00,1.30\});\ V_s=\mathrm{EPS}_s\times\{18.0,25.0,32.0\}\Rightarrow\{19.95,32.59,47.97\}\) & 熊/基准/牛：19.95/\allowbreak32.59/\allowbreak47.97元；概率值：31.87元；空间：-57.0\%。 & 质量基础高；事实输入：官方预告、Q1财务与行情；外部背景：太平洋（2026-06-12）。监控：验证：产品放量、集采/价格、研发进度、许可收入和现金流；反证：产品放量不及预期、价格压力、研发里程碑延后或一次性收益回落 \\
603629 利通电子 & 成长持续期需由盈利兑现与现金流验证 & H1扣非EPS 1.6674元；H2/H1=0.70/1.00/1.30；\(\mathrm{EPS}_s=1.6674\times(1+\{0.70,1.00,1.30\});\ V_s=\mathrm{EPS}_s\times\{18.0,26.0,34.0\}\Rightarrow\{51.02,86.71,130.39\}\) & 熊/基准/牛：51.02/\allowbreak86.71/\allowbreak130.39元；概率值：84.74元；空间：-24.4\%。 & 质量基础高；事实输入：官方预告、Q1财务与行情；外部背景：东莞证券（2018-12-04）。监控：验证：收入/利润兑现、现金流、行业价格和估值分母；反证：利润兑现不及预期、现金流背离或当前估值缺少安全边际 \\
600601 方正科技 & 成长持续期需由盈利兑现与现金流验证 & H1扣非EPS 0.1299元；H2/H1=0.70/1.00/1.30；\(\mathrm{EPS}_s=0.1299\times(1+\{0.70,1.00,1.30\});\ V_s=\mathrm{EPS}_s\times\{18.0,26.0,34.0\}\Rightarrow\{3.97,6.75,10.16\}\) & 熊/基准/牛：3.97/\allowbreak6.75/\allowbreak10.16元；概率值：6.60元；空间：-43.6\%。 & 质量基础高；事实输入：官方预告、Q1财务与行情；外部背景：中航证券（2025-07-24）。监控：验证：收入/利润兑现、现金流、行业价格和估值分母；反证：利润兑现不及预期、现金流背离或当前估值缺少安全边际 \\
002475 立讯精密 & 成长持续期需由盈利兑现与现金流验证 & H1扣非EPS 0.8400元；H2/H1=0.70/1.00/1.30；\(\mathrm{EPS}_s=0.8400\times(1+\{0.70,1.00,1.30\});\ V_s=\mathrm{EPS}_s\times\{18.0,26.0,34.0\}\Rightarrow\{25.71,43.68,65.69\}\) & 熊/基准/牛：25.71/\allowbreak43.68/\allowbreak65.69元；概率值：42.69元；空间：-29.4\%。 & 质量基础高；事实输入：官方预告、Q1财务与行情；外部背景：招银国际（2025-11-05）。监控：验证：收入/利润兑现、现金流、行业价格和估值分母；反证：利润兑现不及预期、现金流背离或当前估值缺少安全边际 \\
600487 亨通光电 & 成长持续期需由盈利兑现与现金流验证 & H1扣非EPS 1.3485元；H2/H1=0.70/1.00/1.30；\(\mathrm{EPS}_s=1.3485\times(1+\{0.70,1.00,1.30\});\ V_s=\mathrm{EPS}_s\times\{18.0,26.0,34.0\}\Rightarrow\{41.26,70.12,105.45\}\) & 熊/基准/牛：41.26/\allowbreak70.12/\allowbreak105.45元；概率值：70.18元；空间：-1.2\%。 & 质量基础高；事实输入：官方预告、Q1财务与行情；外部背景：群益证券（2026-07-14）。监控：验证：收入/利润兑现、现金流、行业价格和估值分母；反证：利润兑现不及预期、现金流背离或当前估值缺少安全边际 \\
600183 生益科技 & 成长持续期需由盈利兑现与现金流验证 & H1扣非EPS 1.1604元；H2/H1=0.70/1.00/1.30；\(\mathrm{EPS}_s=1.1604\times(1+\{0.70,1.00,1.30\});\ V_s=\mathrm{EPS}_s\times\{18.0,26.0,34.0\}\Rightarrow\{35.51,60.34,90.74\}\) & 熊/基准/牛：35.51/\allowbreak60.34/\allowbreak90.74元；概率值：63.42元；空间：-58.4\%。 & 质量基础高；事实输入：官方预告、Q1财务与行情；外部背景：西南证券（2026-05-19）。监控：验证：收入/利润兑现、现金流、行业价格和估值分母；反证：利润兑现不及预期、现金流背离或当前估值缺少安全边际 \\
000100 TCL科技 & 成长持续期需由盈利兑现与现金流验证 & H1扣非EPS 0.1457元；H2/H1=0.70/1.00/1.30；\(\mathrm{EPS}_s=0.1457\times(1+\{0.70,1.00,1.30\});\ V_s=\mathrm{EPS}_s\times\{18.0,26.0,34.0\}\Rightarrow\{3.76,9.36,11.39\}\) & 熊/基准/牛：3.76/\allowbreak9.36/\allowbreak11.39元；概率值：8.09元；空间：61.2\%。 & 质量基础高；事实输入：官方预告、Q1财务与行情；外部背景：中邮证券（2026-04-27）。监控：验证：收入/利润兑现、现金流、行业价格和估值分母；反证：利润兑现不及预期、现金流背离或当前估值缺少安全边际 \\
001287 中电港 & 成长持续期需由盈利兑现与现金流验证 & H1扣非EPS 0.6777元；H2/H1=0.70/1.00/1.30；\(\mathrm{EPS}_s=0.6777\times(1+\{0.70,1.00,1.30\});\ V_s=\mathrm{EPS}_s\times\{18.0,26.0,34.0\}\Rightarrow\{19.99,47.84,53.00\}\) & 熊/基准/牛：19.99/\allowbreak47.84/\allowbreak53.00元；概率值：40.52元；空间：52.0\%。 & 质量基础高；事实输入：官方预告、Q1财务与行情；外部背景：爱建证券（2026-07-13）。监控：验证：收入/利润兑现、现金流、行业价格和估值分母；反证：利润兑现不及预期、现金流背离或当前估值缺少安全边际 \\
000725 京东方A & 成长持续期需由盈利兑现与现金流验证 & H1扣非EPS 0.0862元；H2/H1=0.70/1.00/1.30；\(\mathrm{EPS}_s=0.0862\times(1+\{0.70,1.00,1.30\});\ V_s=\mathrm{EPS}_s\times\{18.0,26.0,34.0\}\Rightarrow\{2.64,6.74,6.76\}\) & 熊/基准/牛：2.64/\allowbreak6.74/\allowbreak6.76元；概率值：5.51元；空间：-13.6\%。 & 质量基础高；事实输入：官方预告、Q1财务与行情；外部背景：金元证券（2026-07-10）。监控：验证：收入/利润兑现、现金流、行业价格和估值分母；反证：利润兑现不及预期、现金流背离或当前估值缺少安全边际 \\
605117 德业股份 & 成长持续期需由盈利兑现与现金流验证 & H1扣非EPS 1.9815元；H2/H1=0.70/1.00/1.30；\(\mathrm{EPS}_s=1.9815\times(1+\{0.70,1.00,1.30\});\ V_s=\mathrm{EPS}_s\times\{16.0,22.0,28.0\}\Rightarrow\{53.90,96.25,127.61\}\) & 熊/基准/牛：53.90/\allowbreak96.25/\allowbreak127.61元；概率值：93.84元；空间：9.9\%。 & 质量基础高；事实输入：官方预告、Q1财务与行情；外部背景：群益证券（2026-06-30）。监控：验证：出货量、订单、产品价格、产能利用率和营运资金；反证：订单转化慢、价格竞争、利用率不足或库存/应收继续占用现金 \\
300037 新宙邦 & 成长持续期需由盈利兑现与现金流验证 & H1扣非EPS 1.3066元；H2/H1=0.70/1.00/1.30；\(\mathrm{EPS}_s=1.3066\times(1+\{0.70,1.00,1.30\});\ V_s=\mathrm{EPS}_s\times\{16.0,22.0,28.0\}\Rightarrow\{35.54,70.62,84.14\}\) & 熊/基准/牛：35.54/\allowbreak70.62/\allowbreak84.14元；概率值：68.85元；空间：2.5\%。 & 质量基础高；事实输入：官方预告、Q1财务与行情；外部背景：东吴证券（2026-07-10）。监控：验证：出货量、订单、产品价格、产能利用率和营运资金；反证：订单转化慢、价格竞争、利用率不足或库存/应收继续占用现金 \\
600522 中天科技 & 成长持续期需由盈利兑现与现金流验证 & H1扣非EPS 0.6543元；H2/H1=0.70/1.00/1.30；\(\mathrm{EPS}_s=0.6543\times(1+\{0.70,1.00,1.30\});\ V_s=\mathrm{EPS}_s\times\{18.0,26.0,34.0\}\Rightarrow\{20.02,34.02,51.16\}\) & 熊/基准/牛：20.02/\allowbreak34.02/\allowbreak51.16元；概率值：33.25元；空间：-18.3\%。 & 质量基础高；事实输入：官方预告、Q1财务与行情；外部背景：中邮证券（2026-02-25）。监控：验证：收入/利润兑现、现金流、行业价格和估值分母；反证：利润兑现不及预期、现金流背离或当前估值缺少安全边际 \\
600267 海正药业 & 成长持续期需由盈利兑现与现金流验证 & H1扣非EPS 0.4321元；H2/H1=0.70/1.00/1.30；\(\mathrm{EPS}_s=0.4321\times(1+\{0.70,1.00,1.30\});\ V_s=\mathrm{EPS}_s\times\{18.0,25.0,32.0\}\Rightarrow\{8.98,21.60,31.80\}\) & 熊/基准/牛：8.98/\allowbreak21.60/\allowbreak31.80元；概率值：19.85元；空间：65.7\%。 & 质量基础高；事实输入：官方预告、Q1财务与行情；外部背景：中银证券（2022-08-24）。监控：验证：产品放量、集采/价格、研发进度、许可收入和现金流；反证：产品放量不及预期、价格压力、研发里程碑延后或一次性收益回落 \\
603268 松发股份 & 成长持续期需由盈利兑现与现金流验证 & H1扣非EPS 3.6054元；H2/H1=0.70/1.05/1.40；\(\mathrm{EPS}_s=3.6054\times(1+\{0.70,1.05,1.40\});\ V_s=\mathrm{EPS}_s\times\{25.0,35.0,45.0\}\Rightarrow\{120.34,259.00,389.38\}\) & 熊/基准/牛：120.34/\allowbreak259.00/\allowbreak389.38元；概率值：243.48元；空间：51.7\%。 & 质量基础高；事实输入：官方预告、Q1财务与行情；外部背景：东吴证券（2026-06-02）。监控：验证：订单交付、合同负债、产品结构和应收账款回款；反证：订单确认延后、交付节奏放缓、应收账款继续扩张或估值透支交付 \\
002384 东山精密 & 成长持续期需由盈利兑现与现金流验证 & H1扣非EPS 1.3376元；H2/H1=0.70/1.00/1.30；\(\mathrm{EPS}_s=1.3376\times(1+\{0.70,1.00,1.30\});\ V_s=\mathrm{EPS}_s\times\{18.0,26.0,34.0\}\Rightarrow\{40.93,97.76,104.60\}\) & 熊/基准/牛：40.93/\allowbreak97.76/\allowbreak104.60元；概率值：82.08元；空间：-68.7\%。 & 质量基础高；事实输入：官方预告、Q1财务与行情；外部背景：开源证券（2026-05-05）。监控：验证：收入/利润兑现、现金流、行业价格和估值分母；反证：利润兑现不及预期、现金流背离或当前估值缺少安全边际 \\
000977 浪潮信息 & 成长持续期需由盈利兑现与现金流验证 & H1扣非EPS 1.5697元；H2/H1=0.70/1.05/1.40；\(\mathrm{EPS}_s=1.5697\times(1+\{0.70,1.05,1.40\});\ V_s=\mathrm{EPS}_s\times\{18.0,28.0,38.0\}\Rightarrow\{48.03,108.08,143.15\}\) & 熊/基准/牛：48.03/\allowbreak108.08/\allowbreak143.15元；概率值：96.37元；空间：13.8\%。 & 质量基础高；事实输入：官方预告、Q1财务与行情；外部背景：群益证券（2026-07-08）。监控：验证：合同订单、项目交付、回款、研发投入和利润率；反证：订单确认延后、回款周期拉长、研发费用率上升或估值继续压缩 \\
603127 昭衍新药 & 成长持续期需由盈利兑现与现金流验证 & H1扣非EPS 0.9361元；H2/H1=0.70/1.00/1.30；\(\mathrm{EPS}_s=0.9361\times(1+\{0.70,1.00,1.30\});\ V_s=\mathrm{EPS}_s\times\{18.0,25.0,32.0\}\Rightarrow\{28.64,46.80,68.89\}\) & 熊/基准/牛：28.64/\allowbreak46.80/\allowbreak68.89元；概率值：45.77元；空间：-9.6\%。 & 质量基础高；事实输入：官方预告、Q1财务与行情；外部背景：国信证券（2026-04-20）。监控：验证：产品放量、集采/价格、研发进度、许可收入和现金流；反证：产品放量不及预期、价格压力、研发里程碑延后或一次性收益回落 \\
000938 紫光股份 & 成长持续期需由盈利兑现与现金流验证 & H1扣非EPS 0.6346元；H2/H1=0.70/1.05/1.40；\(\mathrm{EPS}_s=0.6346\times(1+\{0.70,1.05,1.40\});\ V_s=\mathrm{EPS}_s\times\{18.0,28.0,38.0\}\Rightarrow\{19.42,42.28,57.88\}\) & 熊/基准/牛：19.42/\allowbreak42.28/\allowbreak57.88元；概率值：38.54元；空间：8.5\%。 & 质量基础高；事实输入：官方预告、Q1财务与行情；外部背景：开源证券（2026-07-13）。监控：验证：合同订单、项目交付、回款、研发投入和利润率；反证：订单确认延后、回款周期拉长、研发费用率上升或估值继续压缩 \\
002463 沪电股份 & 成长持续期需由盈利兑现与现金流验证 & H1扣非EPS 1.4576元；H2/H1=0.70/1.00/1.30；\(\mathrm{EPS}_s=1.4576\times(1+\{0.70,1.00,1.30\});\ V_s=\mathrm{EPS}_s\times\{18.0,26.0,34.0\}\Rightarrow\{44.60,77.74,113.99\}\) & 熊/基准/牛：44.60/\allowbreak77.74/\allowbreak113.99元；概率值：75.05元；空间：-45.3\%。 & 质量基础高；事实输入：官方预告、Q1财务与行情；外部背景：中邮证券（2026-04-24）。监控：验证：收入/利润兑现、现金流、行业价格和估值分母；反证：利润兑现不及预期、现金流背离或当前估值缺少安全边际 \\
688183 生益电子 & 成长持续期需由盈利兑现与现金流验证 & H1扣非EPS 1.3243元；H2/H1=0.70/1.00/1.30；\(\mathrm{EPS}_s=1.3243\times(1+\{0.70,1.00,1.30\});\ V_s=\mathrm{EPS}_s\times\{18.0,26.0,34.0\}\Rightarrow\{40.52,68.86,103.56\}\) & 熊/基准/牛：40.52/\allowbreak68.86/\allowbreak103.56元；概率值：67.30元；空间：-47.8\%。 & 质量基础高；事实输入：官方预告、Q1财务与行情；外部背景：华安证券（2025-12-31）。监控：验证：收入/利润兑现、现金流、行业价格和估值分母；反证：利润兑现不及预期、现金流背离或当前估值缺少安全边际 \\
300604 长川科技 & 成长持续期需由盈利兑现与现金流验证 & H1扣非EPS 1.4265元；H2/H1=0.70/1.00/1.30；\(\mathrm{EPS}_s=1.4265\times(1+\{0.70,1.00,1.30\});\ V_s=\mathrm{EPS}_s\times\{18.0,26.0,34.0\}\Rightarrow\{43.65,86.06,111.55\}\) & 熊/基准/牛：43.65/\allowbreak86.06/\allowbreak111.55元；概率值：108.60元；空间：-64.3\%。 & 质量基础高；事实输入：官方预告、Q1财务与行情；外部背景：群益证券（2026-06-26）。监控：验证：收入/利润兑现、现金流、行业价格和估值分母；反证：利润兑现不及预期、现金流背离或当前估值缺少安全边际 \\
002916 深南电路 & 成长持续期需由盈利兑现与现金流验证 & H1扣非EPS 3.2004元；H2/H1=0.70/1.00/1.30；\(\mathrm{EPS}_s=3.2004\times(1+\{0.70,1.00,1.30\});\ V_s=\mathrm{EPS}_s\times\{18.0,26.0,34.0\}\Rightarrow\{97.93,196.82,250.27\}\) & 熊/基准/牛：97.93/\allowbreak196.82/\allowbreak250.27元；概率值：177.84元；空间：-52.9\%。 & 质量基础高；事实输入：官方预告、Q1财务与行情；外部背景：招银国际（2026-03-16）。监控：验证：收入/利润兑现、现金流、行业价格和估值分母；反证：利润兑现不及预期、现金流背离或当前估值缺少安全边际 \\
002174 游族网络 & 成长持续期需由盈利兑现与现金流验证 & H1扣非EPS 0.0204元；H2/H1=0.70/1.05/1.40；\(\mathrm{EPS}_s=0.0204\times(1+\{0.70,1.05,1.40\});\ V_s=\mathrm{EPS}_s\times\{10.0,14.0,18.0\}\Rightarrow\{0.35,0.59,0.88\}\) & 熊/基准/牛：0.35/\allowbreak0.59/\allowbreak0.88元；概率值：0.58元；空间：-95.4\%。 & 质量基础高；事实输入：官方预告、Q1财务与行情；外部背景：东吴证券（2020-08-28）。监控：验证：产品上线、流水/付费、内容成本和经营现金流；反证：产品或内容兑现延后、用户/付费不及预期或现金流不能跟随利润 \\
603296 华勤技术 & 成长持续期需由盈利兑现与现金流验证 & H1扣非EPS 1.1097元；H2/H1=0.70/1.00/1.30；\(\mathrm{EPS}_s=1.1097\times(1+\{0.70,1.00,1.30\});\ V_s=\mathrm{EPS}_s\times\{18.0,26.0,34.0\}\Rightarrow\{33.96,86.78,122.72\}\) & 熊/基准/牛：33.96/\allowbreak86.78/\allowbreak122.72元；概率值：78.12元；空间：-0.1\%。 & 质量基础高；事实输入：官方预告、Q1财务与行情；外部背景：太平洋（2026-05-22）。监控：验证：收入/利润兑现、现金流、行业价格和估值分母；反证：利润兑现不及预期、现金流背离或当前估值缺少安全边际 \\
\bottomrule
\end{longtable}
\normalsize

\section{M3｜金融资产：以PB/ROE与盈利混合估值}
适用于证券、保险和类金融资产。PB保留资产负债表与ROE锚，PE检验盈利能力，二者混合可降低单季投资收益或估值波动造成的偏差。

\footnotesize
\setlength{\tabcolsep}{2.5pt}
\begin{longtable}{L{1.55cm}L{2.45cm}L{5.0cm}L{2.75cm}L{4.1cm}}
\toprule
\textbf{代码/标的} & \textbf{适用理由} & \textbf{分母与公式代入} & \textbf{情景与概率值} & \textbf{证据与验证} \\
\midrule
\endfirsthead
\toprule
\textbf{代码/标的} & \textbf{适用理由} & \textbf{分母与公式代入} & \textbf{情景与概率值} & \textbf{证据与验证} \\
\midrule
\endhead
601688 华泰证券 & 资产净值与盈利能力双锚 & Q1 BPS 19.73元；H1扣非EPS 1.2761元；\(V_s=0.65\times(\mathrm{BPS}=19.73\times\{0.75,1.05,1.35\})+0.35\times(\mathrm{EPS}_s=\{2.1694,2.5523,2.8713\}\times\{8.0,11.0,14.0\})\Rightarrow\{15.47,23.29,31.38\}\) & 熊/基准/牛：15.47/\allowbreak23.29/\allowbreak31.38元；概率值：22.56元；空间：9.4\%。 & 质量基础高；事实输入：官方预告、Q1财务与行情；外部背景：华源证券（2026-05-09）。监控：验证：资本市场成交、两融/权益业务、投资收益和净资本约束；反证：市场成交回落、投资收益波动、信用风险上升或估值缺乏安全边际 \\
601336 新华保险 & 资产净值与盈利能力双锚 & Q1 BPS 39.60元；H1扣非EPS 7.1400元；\(V_s=0.65\times(\mathrm{BPS}=39.60\times\{0.75,1.05,1.35\})+0.35\times(\mathrm{EPS}_s=\{12.1380,14.2800,16.0650\}\times\{8.0,11.0,14.0\})\Rightarrow\{49.16,82.01,113.47\}\) & 熊/基准/牛：49.16/\allowbreak82.01/\allowbreak113.47元；概率值：78.45元；空间：19.7\%。 & 质量基础高；事实输入：官方预告、Q1财务与行情；外部背景：东吴证券（2026-07-14）。监控：验证：资本市场成交、两融/权益业务、投资收益和净资本约束；反证：市场成交回落、投资收益波动、信用风险上升或估值缺乏安全边际 \\
000567 海德股份 & 资产净值与盈利能力双锚 & Q1 BPS 3.08元；H1扣非EPS 0.7367元；\(V_s=0.65\times(\mathrm{BPS}=3.08\times\{0.75,1.05,1.35\})+0.35\times(\mathrm{EPS}_s=\{1.2524,1.4734,1.6576\}\times\{8.0,11.0,14.0\})\Rightarrow\{4.63,7.78,10.83\}\) & 熊/基准/牛：4.63/\allowbreak7.78/\allowbreak10.83元；概率值：7.45元；空间：20.8\%。 & 质量基础高；事实输入：官方预告、Q1财务与行情；外部背景：太平洋（2024-06-20）。监控：验证：资本市场成交、两融/权益业务、投资收益和净资本约束；反证：市场成交回落、投资收益波动、信用风险上升或估值缺乏安全边际 \\
600095 湘财股份 & 资产净值与盈利能力双锚 & Q1 BPS 4.30元；H1扣非EPS 0.1801元；\(V_s=0.65\times(\mathrm{BPS}=4.30\times\{0.75,1.05,1.35\})+0.35\times(\mathrm{EPS}_s=\{0.3062,0.3602,0.4053\}\times\{8.0,11.0,14.0\})\Rightarrow\{2.95,4.32,5.76\}\) & 熊/基准/牛：2.95/\allowbreak4.32/\allowbreak5.76元；概率值：4.20元；空间：-46.2\%。 & 质量基础高；事实输入：官方预告、Q1财务与行情；外部背景：国金证券（2025-07-30）。监控：验证：资本市场成交、两融/权益业务、投资收益和净资本约束；反证：市场成交回落、投资收益波动、信用风险上升或估值缺乏安全边际 \\
600918 中泰证券 & 资产净值与盈利能力双锚 & Q1 BPS 5.64元；H1扣非EPS 0.2144元；\(V_s=0.65\times(\mathrm{BPS}=5.64\times\{0.75,1.05,1.35\})+0.35\times(\mathrm{EPS}_s=\{0.3645,0.4289,0.4825\}\times\{8.0,11.0,14.0\})\Rightarrow\{3.77,5.50,7.32\}\) & 熊/基准/牛：3.77/\allowbreak5.50/\allowbreak7.32元；概率值：5.35元；空间：-6.3\%。 & 质量基础高；事实输入：官方预告、Q1财务与行情；外部背景：太平洋（2025-10-31）。监控：验证：资本市场成交、两融/权益业务、投资收益和净资本约束；反证：市场成交回落、投资收益波动、信用风险上升或估值缺乏安全边际 \\
000987 越秀资本 & 资产净值与盈利能力双锚 & Q1 BPS 6.64元；H1扣非EPS 0.4648元；\(V_s=0.65\times(\mathrm{BPS}=6.64\times\{0.75,1.05,1.35\})+0.35\times(\mathrm{EPS}_s=\{0.7902,0.9296,1.0459\}\times\{8.0,11.0,14.0\})\Rightarrow\{5.45,8.11,10.95\}\) & 熊/基准/牛：5.45/\allowbreak8.11/\allowbreak10.95元；概率值：7.88元；空间：1.4\%。 & 质量基础高；事实输入：官方预告、Q1财务与行情；外部背景：太平洋（2019-08-19）。监控：验证：资本市场成交、两融/权益业务、投资收益和净资本约束；反证：市场成交回落、投资收益波动、信用风险上升或估值缺乏安全边际 \\
601377 兴业证券 & 资产净值与盈利能力双锚 & Q1 BPS 6.56元；H1扣非EPS 0.2426元；\(V_s=0.65\times(\mathrm{BPS}=6.56\times\{0.75,1.05,1.35\})+0.35\times(\mathrm{EPS}_s=\{0.4124,0.4852,0.5458\}\times\{8.0,11.0,14.0\})\Rightarrow\{4.35,6.34,8.43\}\) & 熊/基准/牛：4.35/\allowbreak6.34/\allowbreak8.43元；概率值：6.16元；空间：-0.8\%。 & 质量基础高；事实输入：官方预告、Q1财务与行情；外部背景：开源证券（2026-05-05）。监控：验证：资本市场成交、两融/权益业务、投资收益和净资本约束；反证：市场成交回落、投资收益波动、信用风险上升或估值缺乏安全边际 \\
600120 浙江东方 & 资产净值与盈利能力双锚 & Q1 BPS 4.80元；H1扣非EPS 0.2459元；\(V_s=0.65\times(\mathrm{BPS}=4.80\times\{0.75,1.05,1.35\})+0.35\times(\mathrm{EPS}_s=\{0.4181,0.4919,0.5534\}\times\{8.0,11.0,14.0\})\Rightarrow\{3.51,5.17,6.93\}\) & 熊/基准/牛：3.51/\allowbreak5.17/\allowbreak6.93元；概率值：5.02元；空间：2.2\%。 & 质量中低；事实输入：官方预告、Q1财务与行情；外部背景：无当前券商目标锚。监控：验证：资本市场成交、两融/权益业务、投资收益和净资本约束；反证：市场成交回落、投资收益波动、信用风险上升或估值缺乏安全边际 \\
601628 中国人寿 & 资产净值与盈利能力双锚 & Q1 BPS 21.17元；H1扣非EPS 4.7153元；\(V_s=0.65\times(\mathrm{BPS}=21.17\times\{0.75,1.05,1.35\})+0.35\times(\mathrm{EPS}_s=\{8.0160,9.4306,10.6094\}\times\{8.0,11.0,14.0\})\Rightarrow\{30.32,50.75,70.56\}\) & 熊/基准/牛：30.32/\allowbreak50.75/\allowbreak70.56元；概率值：48.58元；空间：20.2\%。 & 质量基础高；事实输入：官方预告、Q1财务与行情；外部背景：国信证券（2026-07-16）。监控：验证：资本市场成交、两融/权益业务、投资收益和净资本约束；反证：市场成交回落、投资收益波动、信用风险上升或估值缺乏安全边际 \\
600999 招商证券 & 资产净值与盈利能力双锚 & Q1 BPS 14.30元；H1扣非EPS 1.2074元；\(V_s=0.65\times(\mathrm{BPS}=14.30\times\{0.75,1.05,1.35\})+0.35\times(\mathrm{EPS}_s=\{2.0525,2.4148,2.7166\}\times\{8.0,11.0,14.0\})\Rightarrow\{12.72,19.05,25.86\}\) & 熊/基准/牛：12.72/\allowbreak19.05/\allowbreak25.86元；概率值：18.51元；空间：-11.4\%。 & 质量基础高；事实输入：官方预告、Q1财务与行情；外部背景：国信证券（2026-07-10）。监控：验证：资本市场成交、两融/权益业务、投资收益和净资本约束；反证：市场成交回落、投资收益波动、信用风险上升或估值缺乏安全边际 \\
000776 广发证券 & 资产净值与盈利能力双锚 & Q1 BPS 17.61元；H1扣非EPS 1.4940元；\(V_s=0.65\times(\mathrm{BPS}=17.61\times\{0.75,1.05,1.35\})+0.35\times(\mathrm{EPS}_s=\{2.5397,2.9879,3.3614\}\times\{8.0,11.0,14.0\})\Rightarrow\{15.69,23.52,31.92\}\) & 熊/基准/牛：15.69/\allowbreak23.52/\allowbreak31.92元；概率值：22.85元；空间：0.3\%。 & 质量基础高；事实输入：官方预告、Q1财务与行情；外部背景：华源证券（2026-04-30）。监控：验证：资本市场成交、两融/权益业务、投资收益和净资本约束；反证：市场成交回落、投资收益波动、信用风险上升或估值缺乏安全边际 \\
600030 中信证券 & 资产净值与盈利能力双锚 & Q1 BPS 19.67元；H1扣非EPS 1.5931元；\(V_s=0.65\times(\mathrm{BPS}=19.67\times\{0.75,1.05,1.35\})+0.35\times(\mathrm{EPS}_s=\{2.7082,3.1861,3.5844\}\times\{8.0,11.0,14.0\})\Rightarrow\{17.17,25.69,34.82\}\) & 熊/基准/牛：17.17/\allowbreak25.69/\allowbreak34.82元；概率值：24.96元；空间：-12.7\%。 & 质量基础高；事实输入：官方预告、Q1财务与行情；外部背景：东吴证券（2026-07-10）。监控：验证：资本市场成交、两融/权益业务、投资收益和净资本约束；反证：市场成交回落、投资收益波动、信用风险上升或估值缺乏安全边际 \\
601995 中金公司 & 资产净值与盈利能力双锚 & Q1 BPS 21.27元；H1扣非EPS 1.6173元；\(V_s=0.65\times(\mathrm{BPS}=21.27\times\{0.75,1.05,1.35\})+0.35\times(\mathrm{EPS}_s=\{2.7494,3.2345,3.6389\}\times\{8.0,11.0,14.0\})\Rightarrow\{18.07,26.97,36.49\}\) & 熊/基准/牛：18.07/\allowbreak26.97/\allowbreak36.49元；概率值：26.20元；空间：-31.7\%。 & 质量基础高；事实输入：官方预告、Q1财务与行情；外部背景：国信证券（2026-05-05）。监控：验证：资本市场成交、两融/权益业务、投资收益和净资本约束；反证：市场成交回落、投资收益波动、信用风险上升或估值缺乏安全边际 \\
601066 中信建投 & 资产净值与盈利能力双锚 & Q1 BPS 11.12元；H1扣非EPS 1.0043元；\(V_s=0.65\times(\mathrm{BPS}=11.12\times\{0.75,1.05,1.35\})+0.35\times(\mathrm{EPS}_s=\{1.7073,2.0086,2.2597\}\times\{8.0,11.0,14.0\})\Rightarrow\{10.20,15.32,20.83\}\) & 熊/基准/牛：10.20/\allowbreak15.32/\allowbreak20.83元；概率值：14.89元；空间：-46.5\%。 & 质量基础高；事实输入：官方预告、Q1财务与行情；外部背景：东吴证券（2026-04-29）。监控：验证：资本市场成交、两融/权益业务、投资收益和净资本约束；反证：市场成交回落、投资收益波动、信用风险上升或估值缺乏安全边际 \\
601211 国泰海通 & 资产净值与盈利能力双锚 & Q1 BPS 18.52元；H1扣非EPS 1.1063元；\(V_s=0.65\times(\mathrm{BPS}=18.52\times\{0.75,1.05,1.35\})+0.35\times(\mathrm{EPS}_s=\{1.8807,2.2126,2.4892\}\times\{8.0,11.0,14.0\})\Rightarrow\{13.93,21.16,28.45\}\) & 熊/基准/牛：13.93/\allowbreak21.16/\allowbreak28.45元；概率值：20.45元；空间：10.1\%。 & 质量基础高；事实输入：官方预告、Q1财务与行情；外部背景：东吴证券（2026-07-05）。监控：验证：资本市场成交、两融/权益业务、投资收益和净资本约束；反证：市场成交回落、投资收益波动、信用风险上升或估值缺乏安全边际 \\
000783 长江证券 & 资产净值与盈利能力双锚 & Q1 BPS 6.66元；H1扣非EPS 0.5758元；\(V_s=0.65\times(\mathrm{BPS}=6.66\times\{0.75,1.05,1.35\})+0.35\times(\mathrm{EPS}_s=\{0.9788,1.1515,1.2955\}\times\{8.0,11.0,14.0\})\Rightarrow\{5.99,8.98,12.19\}\) & 熊/基准/牛：5.99/\allowbreak8.98/\allowbreak12.19元；概率值：8.72元；空间：-3.8\%。 & 质量基础高；事实输入：官方预告、Q1财务与行情；外部背景：中航证券（2021-05-06）。监控：验证：资本市场成交、两融/权益业务、投资收益和净资本约束；反证：市场成交回落、投资收益波动、信用风险上升或估值缺乏安全边际 \\
600909 华安证券 & 资产净值与盈利能力双锚 & Q1 BPS 5.25元；H1扣非EPS 0.4299元；\(V_s=0.65\times(\mathrm{BPS}=5.25\times\{0.75,1.05,1.35\})+0.35\times(\mathrm{EPS}_s=\{0.7309,0.8598,0.9673\}\times\{8.0,11.0,14.0\})\Rightarrow\{4.61,6.89,9.35\}\) & 熊/基准/牛：4.61/\allowbreak6.89/\allowbreak9.35元；概率值：6.70元；空间：-29.5\%。 & 质量基础高；事实输入：官方预告、Q1财务与行情；外部背景：太平洋（2026-04-23）。监控：验证：资本市场成交、两融/权益业务、投资收益和净资本约束；反证：市场成交回落、投资收益波动、信用风险上升或估值缺乏安全边际 \\
600906 财达证券 & 资产净值与盈利能力双锚 & Q1 BPS 3.86元；H1扣非EPS 0.2344元；\(V_s=0.65\times(\mathrm{BPS}=3.86\times\{0.75,1.05,1.35\})+0.35\times(\mathrm{EPS}_s=\{0.3984,0.4687,0.5273\}\times\{8.0,11.0,14.0\})\Rightarrow\{3.00,4.44,5.97\}\) & 熊/基准/牛：3.00/\allowbreak4.44/\allowbreak5.97元；概率值：4.31元；空间：-40.6\%。 & 质量基础高；事实输入：官方预告、Q1财务与行情；外部背景：东兴证券（2022-02-24）。监控：验证：资本市场成交、两融/权益业务、投资收益和净资本约束；反证：市场成交回落、投资收益波动、信用风险上升或估值缺乏安全边际 \\
\bottomrule
\end{longtable}
\normalsize

\section{M4｜常规盈利：以扣非EPS情景估值}
适用于正扣非分母、商业模式相对稳定且未触发低基数/结算型阻断的公司。H2转化率明确写为情景假设，不将半年度数据伪装为全年共识。

\footnotesize
\setlength{\tabcolsep}{2.5pt}
\begin{longtable}{L{1.55cm}L{2.45cm}L{5.0cm}L{2.75cm}L{4.1cm}}
\toprule
\textbf{代码/标的} & \textbf{适用理由} & \textbf{分母与公式代入} & \textbf{情景与概率值} & \textbf{证据与验证} \\
\midrule
\endfirsthead
\toprule
\textbf{代码/标的} & \textbf{适用理由} & \textbf{分母与公式代入} & \textbf{情景与概率值} & \textbf{证据与验证} \\
\midrule
\endhead
000858 五粮液 & 正扣非分母，采用H2情景而非全年点估 & H1扣非EPS 2.2400元；H2/H1=0.65/0.90/1.20；\(\mathrm{EPS}_s=2.2400\times(1+\{0.65,0.90,1.20\});\ V_s=\mathrm{EPS}_s\times\{12.0,18.0,24.0\}\Rightarrow\{44.35,76.61,118.27\}\) & 熊/基准/牛：44.35/\allowbreak76.61/\allowbreak118.27元；概率值：75.26元；空间：-1.5\%。 & 质量基础高；事实输入：官方预告、Q1财务与行情；外部背景：中邮证券（2026-05-06）。监控：验证：销量、渠道库存、批价、终端动销和现金回款；反证：终端动销弱于预期、渠道库存上升、批价承压或现金回款恶化 \\
601061 中信金属 & 正扣非分母，采用H2情景而非全年点估 & H1扣非EPS 0.5398元；H2/H1=0.65/0.90/1.20；\(\mathrm{EPS}_s=0.5398\times(1+\{0.65,0.90,1.20\});\ V_s=\mathrm{EPS}_s\times\{12.0,18.0,24.0\}\Rightarrow\{8.12,18.46,28.50\}\) & 熊/基准/牛：8.12/\allowbreak18.46/\allowbreak28.50元；概率值：17.37元；空间：60.5\%。 & 质量基础高；事实输入：官方预告、Q1财务与行情；外部背景：西南证券（2025-09-05）。监控：验证：商品价格、贸易量、库存周转、毛利和经营现金流；反证：贸易量或价差下行、库存周转恶化、毛利承压或现金流弱于利润 \\
600292 电投水电 & 正扣非分母，采用H2情景而非全年点估 & H1扣非EPS 0.2701元；H2/H1=0.65/0.90/1.20；\(\mathrm{EPS}_s=0.2701\times(1+\{0.65,0.90,1.20\});\ V_s=\mathrm{EPS}_s\times\{9.0,13.0,17.0\}\Rightarrow\{4.01,6.67,10.10\}\) & 熊/基准/牛：4.01/\allowbreak6.67/\allowbreak10.10元；概率值：6.56元；空间：-47.5\%。 & 质量基础高；事实输入：官方预告、Q1财务与行情；外部背景：西南证券（2025-09-12）。监控：验证：利用小时、上网电价、燃料/资源成本、资本开支和现金流；反证：利用率或电价下行、成本抬升、资本开支超预期或现金流不达标 \\
000155 川能动力 & 正扣非分母，采用H2情景而非全年点估 & H1扣非EPS 0.3656元；H2/H1=0.65/0.90/1.20；\(\mathrm{EPS}_s=0.3656\times(1+\{0.65,0.90,1.20\});\ V_s=\mathrm{EPS}_s\times\{9.0,13.0,17.0\}\Rightarrow\{5.43,9.03,13.67\}\) & 熊/基准/牛：5.43/\allowbreak9.03/\allowbreak13.67元；概率值：10.04元；空间：-9.5\%。 & 质量基础高；事实输入：官方预告、Q1财务与行情；外部背景：国信证券（2026-04-20）。监控：验证：利用小时、上网电价、燃料/资源成本、资本开支和现金流；反证：利用率或电价下行、成本抬升、资本开支超预期或现金流不达标 \\
301377 鼎泰高科 & 正扣非分母，采用H2情景而非全年点估 & H1扣非EPS 1.5564元；H2/H1=0.65/0.90/1.20；\(\mathrm{EPS}_s=1.5564\times(1+\{0.65,0.90,1.20\});\ V_s=\mathrm{EPS}_s\times\{14.0,20.0,28.0\}\Rightarrow\{35.95,61.20,95.88\}\) & 熊/基准/牛：35.95/\allowbreak61.20/\allowbreak95.88元；概率值：60.56元；空间：-87.1\%。 & 质量基础高；事实输入：官方预告、Q1财务与行情；外部背景：东吴证券（2026-07-15）。监控：验证：收入/利润兑现、现金流、行业价格和估值分母；反证：利润兑现不及预期、现金流背离或当前估值缺少安全边际 \\
688257 新锐股份 & 正扣非分母，采用H2情景而非全年点估 & H1扣非EPS 1.6608元；H2/H1=0.65/0.90/1.20；\(\mathrm{EPS}_s=1.6608\times(1+\{0.65,0.90,1.20\});\ V_s=\mathrm{EPS}_s\times\{14.0,20.0,28.0\}\Rightarrow\{38.36,63.11,102.30\}\) & 熊/基准/牛：38.36/\allowbreak63.11/\allowbreak102.30元；概率值：63.52元；空间：-20.7\%。 & 质量基础高；事实输入：官方预告、Q1财务与行情；外部背景：东吴证券（2026-07-15）。监控：验证：收入/利润兑现、现金流、行业价格和估值分母；反证：利润兑现不及预期、现金流背离或当前估值缺少安全边际 \\
600176 中国巨石 & 正扣非分母，采用H2情景而非全年点估 & H1扣非EPS 0.7437元；H2/H1=0.65/0.90/1.20；\(\mathrm{EPS}_s=0.7437\times(1+\{0.65,0.90,1.20\});\ V_s=\mathrm{EPS}_s\times\{12.0,18.0,24.0\}\Rightarrow\{14.73,26.98,39.27\}\) & 熊/基准/牛：14.73/\allowbreak26.98/\allowbreak39.27元；概率值：28.18元；空间：-46.8\%。 & 质量基础高；事实输入：官方预告、Q1财务与行情；外部背景：群益证券（2026-05-26）。监控：验证：收入/利润兑现、现金流、行业价格和估值分母；反证：利润兑现不及预期、现金流背离或当前估值缺少安全边际 \\
002266 浙富控股 & 正扣非分母，采用H2情景而非全年点估 & H1扣非EPS 0.2395元；H2/H1=0.65/0.90/1.20；\(\mathrm{EPS}_s=0.2395\times(1+\{0.65,0.90,1.20\});\ V_s=\mathrm{EPS}_s\times\{12.0,18.0,24.0\}\Rightarrow\{3.82,8.19,12.65\}\) & 熊/基准/牛：3.82/\allowbreak8.19/\allowbreak12.65元；概率值：7.77元；空间：52.6\%。 & 质量基础高；事实输入：官方预告、Q1财务与行情；外部背景：中邮证券（2026-06-29）。监控：验证：收入/利润兑现、现金流、行业价格和估值分母；反证：利润兑现不及预期、现金流背离或当前估值缺少安全边际 \\
002008 大族激光 & 正扣非分母，采用H2情景而非全年点估 & H1扣非EPS 1.3598元；H2/H1=0.65/0.90/1.20；\(\mathrm{EPS}_s=1.3598\times(1+\{0.65,0.90,1.20\});\ V_s=\mathrm{EPS}_s\times\{14.0,20.0,28.0\}\Rightarrow\{31.41,52.40,83.76\}\) & 熊/基准/牛：31.41/\allowbreak52.40/\allowbreak83.76元；概率值：52.38元；空间：-54.2\%。 & 质量基础高；事实输入：官方预告、Q1财务与行情；外部背景：东吴证券（2026-07-10）。监控：验证：收入/利润兑现、现金流、行业价格和估值分母；反证：利润兑现不及预期、现金流背离或当前估值缺少安全边际 \\
000685 中山公用 & 正扣非分母，采用H2情景而非全年点估 & H1扣非EPS 0.9724元；H2/H1=0.65/0.90/1.20；\(\mathrm{EPS}_s=0.9724\times(1+\{0.65,0.90,1.20\});\ V_s=\mathrm{EPS}_s\times\{12.0,18.0,24.0\}\Rightarrow\{8.39,33.26,51.34\}\) & 熊/基准/牛：8.39/\allowbreak33.26/\allowbreak51.34元；概率值：29.41元；空间：162.8\%。 & 质量基础高；事实输入：官方预告、Q1财务与行情；外部背景：国信证券（2026-05-11）。监控：验证：收入/利润兑现、现金流、行业价格和估值分母；反证：利润兑现不及预期、现金流背离或当前估值缺少安全边际 \\
000833 粤桂股份 & 正扣非分母，采用H2情景而非全年点估 & H1扣非EPS 0.8141元；H2/H1=0.65/0.90/1.20；\(\mathrm{EPS}_s=0.8141\times(1+\{0.65,0.90,1.20\});\ V_s=\mathrm{EPS}_s\times\{12.0,18.0,24.0\}\Rightarrow\{15.05,27.84,42.99\}\) & 熊/基准/牛：15.05/\allowbreak27.84/\allowbreak42.99元；概率值：27.03元；空间：34.7\%。 & 质量基础高；事实输入：官方预告、Q1财务与行情；外部背景：东兴证券（2018-08-31）。监控：验证：收入/利润兑现、现金流、行业价格和估值分母；反证：利润兑现不及预期、现金流背离或当前估值缺少安全边际 \\
002458 益生股份 & 正扣非分母，采用H2情景而非全年点估 & H1扣非EPS 0.1977元；H2/H1=0.65/0.90/1.20；\(\mathrm{EPS}_s=0.1977\times(1+\{0.65,0.90,1.20\});\ V_s=\mathrm{EPS}_s\times\{12.0,18.0,24.0\}\Rightarrow\{3.92,8.82,10.44\}\) & 熊/基准/牛：3.92/\allowbreak8.82/\allowbreak10.44元；概率值：7.67元；空间：-15.0\%。 & 质量基础高；事实输入：官方预告、Q1财务与行情；外部背景：太平洋（2026-03-30）。监控：验证：收入/利润兑现、现金流、行业价格和估值分母；反证：利润兑现不及预期、现金流背离或当前估值缺少安全边际 \\
301200 大族数控 & 正扣非分母，采用H2情景而非全年点估 & H1扣非EPS 1.9427元；H2/H1=0.65/0.90/1.20；\(\mathrm{EPS}_s=1.9427\times(1+\{0.65,0.90,1.20\});\ V_s=\mathrm{EPS}_s\times\{14.0,20.0,28.0\}\Rightarrow\{44.88,85.00,119.67\}\) & 熊/基准/牛：44.88/\allowbreak85.00/\allowbreak119.67元；概率值：79.90元；空间：-74.4\%。 & 质量基础高；事实输入：官方预告、Q1财务与行情；外部背景：东吴证券（2026-07-10）。监控：验证：收入/利润兑现、现金流、行业价格和估值分母；反证：利润兑现不及预期、现金流背离或当前估值缺少安全边际 \\
001248 华润新能源 & 正扣非分母，采用H2情景而非全年点估 & H1扣非EPS 0.2722元；H2/H1=0.65/0.90/1.20；\(\mathrm{EPS}_s=0.2722\times(1+\{0.65,0.90,1.20\});\ V_s=\mathrm{EPS}_s\times\{9.0,13.0,17.0\}\Rightarrow\{4.04,6.72,10.18\}\) & 熊/基准/牛：4.04/\allowbreak6.72/\allowbreak10.18元；概率值：6.61元；空间：-50.7\%。 & 质量基础高；事实输入：官方预告、Q1财务与行情；外部背景：华金证券（2026-06-21）。监控：验证：利用小时、上网电价、燃料/资源成本、资本开支和现金流；反证：利用率或电价下行、成本抬升、资本开支超预期或现金流不达标 \\
601633 长城汽车 & 正扣非分母，采用H2情景而非全年点估 & H1扣非EPS 0.1900元；H2/H1=0.65/0.90/1.20；\(\mathrm{EPS}_s=0.1900\times(1+\{0.65,0.90,1.20\});\ V_s=\mathrm{EPS}_s\times\{12.0,16.0,20.0\}\Rightarrow\{3.76,8.36,23.65\}\) & 熊/基准/牛：3.76/\allowbreak8.36/\allowbreak23.65元；概率值：11.25元；空间：-27.8\%。 & 质量基础高；事实输入：官方预告、Q1财务与行情；外部背景：国金证券（2026-06-03）。监控：验证：销量、单车盈利、产品结构、价格竞争和应收回款；反证：销量或单车盈利下滑、价格战加剧、产品结构恶化或现金流转弱 \\
002611 东方精工 & 正扣非分母，采用H2情景而非全年点估 & H1扣非EPS 0.0739元；H2/H1=0.65/0.90/1.20；\(\mathrm{EPS}_s=0.0739\times(1+\{0.65,0.90,1.20\});\ V_s=\mathrm{EPS}_s\times\{14.0,20.0,28.0\}\Rightarrow\{1.71,2.81,4.55\}\) & 熊/基准/牛：1.71/\allowbreak2.81/\allowbreak4.55元；概率值：2.83元；空间：-81.2\%。 & 质量基础高；事实输入：官方预告、Q1财务与行情；外部背景：爱建证券（2025-12-25）。监控：验证：收入/利润兑现、现金流、行业价格和估值分母；反证：利润兑现不及预期、现金流背离或当前估值缺少安全边际 \\
\bottomrule
\end{longtable}
\normalsize

\section{R1｜恢复估值：以正常化EPS、PB或PS重建分母}
适用于低基数扭亏、项目结算、近零EPS或利润周期明显扭曲的公司。先切换至正常化EPS、PB或收入PS，再决定是否使用外部锚。

\footnotesize
\setlength{\tabcolsep}{2.5pt}
\begin{longtable}{L{1.55cm}L{2.45cm}L{5.0cm}L{2.75cm}L{4.1cm}}
\toprule
\textbf{代码/标的} & \textbf{适用理由} & \textbf{分母与公式代入} & \textbf{情景与概率值} & \textbf{证据与验证} \\
\midrule
\endfirsthead
\toprule
\textbf{代码/标的} & \textbf{适用理由} & \textbf{分母与公式代入} & \textbf{情景与概率值} & \textbf{证据与验证} \\
\midrule
\endhead
000042 中洲控股 & 低基数/结算扭曲，先切换正常化分母 & 项目结算、税费与利息正常化EPS；\(V_s=\mathrm{EPS}_{norm}\times\mathrm{PE}_s=\{0.80,1.10,1.50\}\times\{6.0,8.0,9.0\}\Rightarrow\{4.80,8.80,13.50\}\) & 熊/基准/牛：4.80/\allowbreak8.80/\allowbreak13.50元；概率值：8.54元；空间：6.5\%。 & 质量基础高；事实输入：官方预告、Q1财务与行情；外部背景：无当前券商目标锚。监控：验证：2026年EPS不低于1.10元，项目结算及税费/利息计提得到确认，债务与现金流压力改善且股价守住8元；反证：结算利润回转、土地增值税或利息计提上升、销售回款走弱，或净负债/流动性恶化 \\
002460 赣锋锂业 & 低基数/结算扭曲，先切换正常化分母 & 2026E周期调整EPS；\(V_s=\mathrm{EPS}_{norm}\times\mathrm{PE}_s=\{4.00,4.45,5.35\}\times\{10.0,12.0,14.0\}\Rightarrow\{40.00,53.40,74.90\}\) & 熊/基准/牛：40.00/\allowbreak53.40/\allowbreak74.90元；概率值：53.68元；空间：4.2\%。 & 质量基础高；事实输入：官方预告、Q1财务与行情；外部背景：东吴证券（2026-01-28）。监控：验证：锂价、产量、资源自给率和H2经营现金流共同验证2026年EPS 4.45元；反证：锂价回落、资源项目延期、库存减值，或现金流弱于利润 \\
600037 歌华有线 & 低基数/结算扭曲，先切换正常化分母 & 2026E持续经营EPS；\(V_s=\mathrm{EPS}_{norm}\times\mathrm{PE}_s=\{0.02,0.03,0.05\}\times\{120.0,180.0,220.0\}\Rightarrow\{2.40,5.40,11.00\}\) & 熊/基准/牛：2.40/\allowbreak5.40/\allowbreak11.00元；概率值：5.62元；空间：-17.3\%。 & 质量基础高；事实输入：官方预告、Q1财务与行情；外部背景：中邮证券（2026-03-23）。监控：验证：通信业务收入、经营现金流和全年EPS达到公开预测，并连续两个季度得到验证；反证：扭亏由一次性收益驱动、通信业务不及预期，或经营现金流转弱 \\
002048 宁波华翔 & 低基数/结算扭曲，先切换正常化分母 & 汽车零部件与机器人订单正常化EPS；\(V_s=\mathrm{EPS}_{norm}\times\mathrm{PE}_s=\{1.60,1.98,2.20\}\times\{12.0,15.0,18.0\}\Rightarrow\{19.20,29.64,39.60\}\) & 熊/基准/牛：19.20/\allowbreak29.64/\allowbreak39.60元；概率值：28.50元；空间：33.0\%。 & 质量基础高；事实输入：官方预告、Q1财务与行情；外部背景：国金证券（2025-10-30）。监控：验证：汽车订单、机器人批量供货和H2扣非利润共同验证2026年EPS 1.98元；反证：主机厂降价、机器人订单延迟，或扣非利润低于预期 \\
600688 上海石化 & 低基数/结算扭曲，先切换正常化分母 & Q1 BPS与炼化周期调整资产价值；\(V_s=\mathrm{BPS}_{norm}\times\mathrm{PB}_s=\{2.23,2.23,2.23\}\times\{0.9,1.2,1.5\}\Rightarrow\{2.01,2.68,3.35\}\) & 熊/基准/牛：2.01/\allowbreak2.68/\allowbreak3.35元；概率值：2.61元；空间：-8.1\%。 & 质量基础高；事实输入：官方预告、Q1财务与行情；外部背景：西南证券（2023-03-30）。监控：验证：炼化价差、经营现金流和资产负债表改善，同时新一期机构预测恢复；反证：炼化价差继续收窄、现金流恶化，或资产减值上升 \\
601360 三六零 & 低基数/结算扭曲，先切换正常化分母 & 2026E收入95.91亿元与AI/安全变现；\(V_s=(R_{2026E}=95.91/\mathrm{Shares}=69.9956)\times\{7.5,10.0,11.0\}\Rightarrow\{10.28,13.70,15.07\}\) & 熊/基准/牛：10.28/\allowbreak13.70/\allowbreak15.07元；概率值：13.10元；空间：45.6\%。 & 质量基础高；事实输入：官方预告、Q1财务与行情；外部背景：信达证券（2023-04-25）。监控：验证：2026年收入接近95.91亿元，AI/安全形成可重复变现，EPS共识不低于0.06元且经营现金流转正；反证：AI产品变现不及预期、2026年EPS低于0.03元、经营现金流仍为负，或PS低于可比区间 \\
001309 德明利 & 低基数/结算扭曲，先切换正常化分母 & 存储周期调整前瞻EPS；\(V_s=\mathrm{EPS}_{norm}\times\mathrm{PE}_s=\{30.00,39.16,45.00\}\times\{11.0,13.1,14.0\}\Rightarrow\{330.00,513.00,630.00\}\) & 熊/基准/牛：330.00/\allowbreak513.00/\allowbreak630.00元；概率值：481.50元；空间：-27.3\%。 & 质量基础高；事实输入：官方预告、Q1财务与行情；外部背景：国信证券（2026-04-21）。监控：验证：存储价格、企业级产品出货和H2经营现金流共同验证2026年EPS 39.16元；反证：存储价格反转、库存减值、客户集中风险暴露，或经营现金流持续为负 \\
300390 天华新能 & 低基数/结算扭曲，先切换正常化分母 & 锂价敏感的2026E周期调整EPS；\(V_s=\mathrm{EPS}_{norm}\times\mathrm{PE}_s=\{4.70,5.70,6.40\}\times\{12.0,14.0,16.0\}\Rightarrow\{56.40,79.80,102.40\}\) & 熊/基准/牛：56.40/\allowbreak79.80/\allowbreak102.40元；概率值：77.30元；空间：28.8\%。 & 质量基础高；事实输入：官方预告、Q1财务与行情；外部背景：华源证券（2026-06-06）。监控：验证：氢氧化锂价格、客户采购、产量和H2经营现金流共同验证2026年EPS 5.70元；反证：锂价回落、客户集中风险暴露，或经营现金流继续弱于利润 \\
002812 恩捷股份 & 低基数/结算扭曲，先切换正常化分母 & 隔膜价格/利用率恢复后的2027E EPS；\(V_s=\mathrm{EPS}_{norm}\times\mathrm{PE}_s=\{1.80,5.16,6.31\}\times\{12.0,20.0,20.0\}\Rightarrow\{21.60,103.20,126.20\}\) & 熊/基准/牛：21.60/\allowbreak103.20/\allowbreak126.20元；概率值：85.29元；空间：55.2\%。 & 质量基础高；事实输入：官方预告、Q1财务与行情；外部背景：东吴证券（2026-07-10）。监控：验证：隔膜涨价、产能利用率、海外/3C订单和H2现金流共同验证2027年EPS 5.16元；反证：供需修复不及预期、价格回落、客户集中，或产能利用率不足 \\
002240 盛新锂能 & 低基数/结算扭曲，先切换正常化分母 & 锂价、自有矿与现金流桥接EPS；\(V_s=\mathrm{EPS}_{norm}\times\mathrm{PE}_s=\{1.80,2.45,3.28\}\times\{12.0,16.0,18.0\}\Rightarrow\{21.60,39.20,59.04\}\) & 熊/基准/牛：21.60/\allowbreak39.20/\allowbreak59.04元；概率值：37.89元；空间：21.8\%。 & 质量基础高；事实输入：官方预告、Q1财务与行情；外部背景：东吴证券（2026-07-09）。监控：验证：锂价、出货、自有矿贡献和H2经营现金流共同验证2026年EPS 2.45元；反证：锂价回落、现金流持续为负，或资源项目贡献不及预期 \\
600736 苏州高新 & 低基数/结算扭曲，先切换正常化分母 & 园区/结算/投资收益正常化EPS；\(V_s=\mathrm{EPS}_{norm}\times\mathrm{PE}_s=\{0.12,0.14,0.16\}\times\{35.0,55.0,65.0\}\Rightarrow\{4.20,7.70,10.40\}\) & 熊/基准/牛：4.20/\allowbreak7.70/\allowbreak10.40元；概率值：7.19元；空间：4.7\%。 & 质量基础高；事实输入：官方预告、Q1财务与行情；外部背景：东吴证券（2026-04-28）。监控：验证：产业园利润、地产结算、投资收益和经营现金流共同验证2026年EPS 0.14元；反证：地产结算延期、投资收益回落，或经营现金流继续为负 \\
000415 渤海租赁 & 低基数/结算扭曲，先切换正常化分母 & 租赁收入、资产质量与汇率正常化EPS；\(V_s=\mathrm{EPS}_{norm}\times\mathrm{PE}_s=\{0.60,0.77,0.85\}\times\{5.0,6.0,7.0\}\Rightarrow\{3.00,4.62,5.95\}\) & 熊/基准/牛：3.00/\allowbreak4.62/\allowbreak5.95元；概率值：4.40元；空间：-1.1\%。 & 质量基础高；事实输入：官方预告、Q1财务与行情；外部背景：国金证券（2026-03-17）。监控：验证：租赁收入、资产质量、汇率和H2现金流共同验证2026年EPS 0.77元；反证：飞机减值、汇率损失、融资成本上升，或现金流风险上行 \\
\bottomrule
\end{longtable}
\normalsize

\section{B1｜上市边界：不在无二级市场价格时制造目标价}
尚无有效二级市场价格时，任何当前价目标和空间均不成立；保留发行价与上市状态，待交易形成后再建立估值模型。

\footnotesize
\setlength{\tabcolsep}{2.5pt}
\begin{longtable}{L{1.55cm}L{2.45cm}L{5.0cm}L{2.75cm}L{4.1cm}}
\toprule
\textbf{代码/标的} & \textbf{适用理由} & \textbf{分母与公式代入} & \textbf{情景与概率值} & \textbf{证据与验证} \\
\midrule
\endfirsthead
\toprule
\textbf{代码/标的} & \textbf{适用理由} & \textbf{分母与公式代入} & \textbf{情景与概率值} & \textbf{证据与验证} \\
\midrule
\endhead
688825 长鑫科技 & 未上市，无法建立当前价模型 & 发行价8.66元；尚无二级市场当前价；不适用：尚未上市，待形成可验证二级市场价格后重建模型。 & 待上市交易；当前价、目标价与空间均不适用。 & 质量低；事实输入：官方预告、Q1财务与行情；外部背景：无当前券商目标锚。监控：形成有效二级市场交易价格后重新评估 \\
\bottomrule
\end{longtable}
\normalsize

\begin{sourcequalitybox}[附录阅读说明]
表中的倍数、H2转化率和概率权重是本机构情景假设；官方预告、Q1财务、行情和原始研报是事实输入。外部目标价只有在来源可审计且时点可比时才获得正权重。
\end{sourcequalitybox}

\chapter{优先池的证据升级路径}

39只优先池的价值在于安排下一轮研究资源，而不在于将模型空间直接转化为配置结论。
本附录第一张表只覆盖39只优先池；100\%以上模型空间的全量样本另列在第二张表，以解释为什么有些高空间票没有进入优先池或正式池。

\footnotesize
\begin{longtable}{L{1.0cm}L{1.6cm}L{2.5cm}R{1.1cm}R{1.0cm}L{7.0cm}}
\toprule
\textbf{代码} & \textbf{标的} & \textbf{状态} & \textbf{概率值} & \textbf{空间} & \textbf{下一轮必须补齐的证据} \\
\midrule
\endfirsthead
\toprule
\textbf{代码} & \textbf{标的} & \textbf{状态} & \textbf{概率值} & \textbf{空间} & \textbf{下一轮必须补齐的证据} \\
\midrule
\endhead
000623 & 吉林敖东 & 条件性高空间验证 & 77.82 & 317.0\% & 缺少当前可正权使用的外部估值锚；最新研报缺失或早于2026-04-01；原始研报已归档但未披露目标价/合理价值区间 \\
600739 & 辽宁成大 & 条件性高空间验证 & 43.80 & 300.0\% & 缺少当前可正权使用的外部估值锚；最新研报缺失或早于2026-04-01 \\
600150 & 中国船舶 & 条件性高空间验证 & 83.33 & 142.9\% & 缺少当前可正权使用的外部估值锚；原始研报已归档但未披露目标价/合理价值区间 \\
000703 & 恒逸石化 & 可复核当前价模型 & 23.42 & 61.7\% & 正式模型：当前价格、情景估值、外部锚和原始研报证据均可复核 \\
601061 & 中信金属 & 条件性高空间验证 & 17.37 & 60.5\% & 缺少当前可正权使用的外部估值锚；最新研报缺失或早于2026-04-01 \\
601600 & 中国铝业 & 条件性高空间验证 & 14.07 & 56.7\% & 基础证据链不足以支持正式模型；缺少当前可正权使用的外部估值锚；最新研报缺失或早于2026-04-01；公开源检索未取得可审计外部目标价/合理价值区间；公开源检索未取得本地原始研报证据路径 \\
601360 & 三六零 & 条件性高空间验证 & 13.10 & 45.6\% & 缺少当前可正权使用的外部估值锚；最新研报缺失或早于2026-04-01 \\
002379 & 宏桥控股 & 可复核当前价模型 & 25.80 & 38.3\% & 正式模型：当前价格、情景估值、外部锚和原始研报证据均可复核 \\
002048 & 宁波华翔 & 条件性高空间验证 & 28.50 & 33.0\% & 缺少当前可正权使用的外部估值锚；最新研报缺失或早于2026-04-01；原始研报已归档但未披露目标价/合理价值区间 \\
000737 & 北方铜业 & 条件性高空间验证 & 15.40 & 20.8\% & 基础证据链不足以支持正式模型；缺少当前可正权使用的外部估值锚；最新研报缺失或早于2026-04-01；公开源检索未取得可审计外部目标价/合理价值区间；公开源检索未取得本地原始研报证据路径 \\
000567 & 海德股份 & 条件性高空间验证 & 7.45 & 20.8\% & 缺少当前可正权使用的外部估值锚；最新研报缺失或早于2026-04-01；原始研报已归档但未披露目标价/合理价值区间 \\
601336 & 新华保险 & 条件性高空间验证 & 78.45 & 19.7\% & 缺少当前可正权使用的外部估值锚；原始研报已归档但未披露目标价/合理价值区间 \\
002602 & 世纪华通 & 条件性高空间验证 & 16.51 & 15.8\% & 缺少当前可正权使用的外部估值锚；原始研报已归档但未披露目标价/合理价值区间 \\
300014 & 亿纬锂能 & 可复核当前价模型 & 60.43 & 15.6\% & 正式模型：当前价格、情景估值、外部锚和原始研报证据均可复核 \\
002558 & 巨人网络 & 条件性安全边际观察 & 33.99 & 15.0\% & 缺少当前可正权使用的外部估值锚；原始研报已归档但未披露目标价/合理价值区间 \\
600026 & 中远海能 & 条件性安全边际观察 & 17.45 & 10.2\% & 缺少当前可正权使用的外部估值锚；最新研报缺失或早于2026-04-01 \\
601688 & 华泰证券 & 条件性安全边际观察 & 22.56 & 9.4\% & 缺少当前可正权使用的外部估值锚；原始研报已归档但未披露目标价/合理价值区间 \\
000042 & 中洲控股 & 条件性安全边际观察 & 8.54 & 6.5\% & 缺少当前可正权使用的外部估值锚；最新研报缺失或早于2026-04-01 \\
600489 & 中金黄金 & 条件性安全边际观察 & 21.16 & 5.2\% & 缺少当前可正权使用的外部估值锚；原始研报已归档但未披露目标价/合理价值区间 \\
000975 & 山金国际 & 条件性安全边际观察 & 20.31 & 5.0\% & 缺少当前可正权使用的外部估值锚；原始研报已归档但未披露目标价/合理价值区间 \\
002460 & 赣锋锂业 & 条件性安全边际观察 & 53.68 & 4.2\% & 缺少当前可正权使用的外部估值锚；最新研报缺失或早于2026-04-01；原始研报已归档但未披露目标价/合理价值区间 \\
600120 & 浙江东方 & 条件性安全边际观察 & 5.02 & 2.2\% & 基础证据链不足以支持正式模型；缺少当前可正权使用的外部估值锚；最新研报缺失或早于2026-04-01；公开源检索未取得可审计外部目标价/合理价值区间；公开源检索未取得本地原始研报证据路径 \\
000987 & 越秀资本 & 条件性安全边际观察 & 7.88 & 1.4\% & 缺少当前可正权使用的外部估值锚；最新研报缺失或早于2026-04-01 \\
601377 & 兴业证券 & 估值纪律/风险观察 & 6.16 & -0.8\% & 缺少当前可正权使用的外部估值锚；原始研报已归档但未披露目标价/合理价值区间；概率目标低于当前价，仅保留为下行风险对照 \\
000858 & 五粮液 & 估值纪律/风险观察 & 75.26 & -1.5\% & 缺少当前可正权使用的外部估值锚；原始研报已归档但未披露目标价/合理价值区间；概率目标低于当前价，仅保留为下行风险对照 \\
000301 & 东方盛虹 & 估值纪律/风险观察 & 11.18 & -3.7\% & 缺少当前可正权使用的外部估值锚；原始研报已归档但未披露目标价/合理价值区间；概率目标低于当前价，仅保留为下行风险对照 \\
002756 & 永兴材料 & 估值纪律/风险观察 & 42.10 & -6.2\% & 缺少当前可正权使用的外部估值锚；原始研报已归档但未披露目标价/合理价值区间；概率目标低于当前价，仅保留为下行风险对照 \\
600918 & 中泰证券 & 估值纪律/风险观察 & 5.35 & -6.3\% & 缺少当前可正权使用的外部估值锚；最新研报缺失或早于2026-04-01；原始研报已归档但未披露目标价/合理价值区间；概率目标低于当前价，仅保留为下行风险对照 \\
002466 & 天齐锂业 & 估值纪律/风险观察 & 44.40 & -6.4\% & 缺少当前可正权使用的外部估值锚；原始研报已归档但未披露目标价/合理价值区间；概率目标低于当前价，仅保留为下行风险对照 \\
600346 & 恒力石化 & 估值纪律/风险观察 & 14.21 & -6.8\% & 缺少当前可正权使用的外部估值锚；原始研报已归档但未披露目标价/合理价值区间；概率目标低于当前价，仅保留为下行风险对照 \\
600688 & 上海石化 & 估值纪律/风险观察 & 2.61 & -8.1\% & 缺少当前可正权使用的外部估值锚；最新研报缺失或早于2026-04-01；概率目标低于当前价，仅保留为下行风险对照 \\
000155 & 川能动力 & 可复核当前价模型 & 10.04 & -9.5\% & 正式下行风险对照：证据链完整，但概率目标低于当前价 \\
600037 & 歌华有线 & 估值纪律/风险观察 & 5.62 & -17.3\% & 缺少当前可正权使用的外部估值锚；最新研报缺失或早于2026-04-01；原始研报已归档但未披露目标价/合理价值区间；概率目标低于当前价，仅保留为下行风险对照 \\
002738 & 中矿资源 & 估值纪律/风险观察 & 39.26 & -20.4\% & 缺少当前可正权使用的外部估值锚；原始研报已归档但未披露目标价/合理价值区间；概率目标低于当前价，仅保留为下行风险对照 \\
000426 & 兴业银锡 & 估值纪律/风险观察 & 23.56 & -22.9\% & 缺少当前可正权使用的外部估值锚；原始研报已归档但未披露目标价/合理价值区间；概率目标低于当前价，仅保留为下行风险对照 \\
001339 & 智微智能 & 估值纪律/风险观察 & 65.43 & -26.7\% & 缺少当前可正权使用的外部估值锚；原始研报已归档但未披露目标价/合理价值区间；概率目标低于当前价，仅保留为下行风险对照 \\
601069 & 西部黄金 & 估值纪律/风险观察 & 12.74 & -44.5\% & 缺少当前可正权使用的外部估值锚；最新研报缺失或早于2026-04-01；原始研报已归档但未披露目标价/合理价值区间；概率目标低于当前价，仅保留为下行风险对照 \\
600095 & 湘财股份 & 估值纪律/风险观察 & 4.20 & -46.2\% & 缺少当前可正权使用的外部估值锚；最新研报缺失或早于2026-04-01；概率目标低于当前价，仅保留为下行风险对照 \\
600292 & 电投水电 & 估值纪律/风险观察 & 6.56 & -47.5\% & 缺少当前可正权使用的外部估值锚；最新研报缺失或早于2026-04-01；原始研报已归档但未披露目标价/合理价值区间；概率目标低于当前价，仅保留为下行风险对照 \\
\bottomrule
\end{longtable}
\normalsize

\section{原中等证据标的补证后的候选准入结论}

本节只覆盖补证前旧口径为中等证据的标的。补证后，基础证据质量已经不再是阻断项；是否进入候选取决于空间、优先池状态、现金流、外部锚时效和价格纪律。

\footnotesize
\begin{longtable}{L{1.0cm}L{1.55cm}L{2.1cm}R{1.0cm}L{2.4cm}L{5.3cm}}
\toprule
\textbf{代码} & \textbf{标的} & \textbf{准入结论} & \textbf{空间} & \textbf{外部锚状态} & \textbf{使用方式/阻断项} \\
\midrule
\endfirsthead
\toprule
\textbf{代码} & \textbf{标的} & \textbf{准入结论} & \textbf{空间} & \textbf{外部锚状态} & \textbf{使用方式/阻断项} \\
\midrule
\endhead
002432 & 九安医疗 & 扩展事件验证候选 & 323.4\% & 原始PDF无目标价 & 可进入扩展观察候选池；缺少当前可正权使用的外部估值锚；经营现金流仍需验证 \\
000623 & 吉林敖东 & 优先验证候选 & 317.0\% & 原始PDF无目标价 & 进入优先验证候选池；缺少当前可正权使用的外部估值锚；最新研报缺失或早于2026-04-01；原始研报已归档但未披露目标价/合理价值区间 \\
600739 & 辽宁成大 & 优先验证候选 & 300.0\% & 陈旧/低权重外部锚 & 进入优先验证候选池；缺少当前可正权使用的外部估值锚；最新研报缺失或早于2026-04-01；经营现金流仍需验证 \\
000612 & 焦作万方 & 扩展事件验证候选 & 90.6\% & 原始PDF无目标价 & 可进入扩展观察候选池；缺少当前可正权使用的外部估值锚 \\
000933 & 神火股份 & 扩展事件验证候选 & 73.3\% & 原始PDF无目标价 & 可进入扩展观察候选池；缺少当前可正权使用的外部估值锚 \\
300475 & 香农芯创 & 扩展事件验证候选 & 68.4\% & 原始PDF无目标价 & 可进入扩展观察候选池；缺少当前可正权使用的外部估值锚 \\
600267 & 海正药业 & 扩展事件验证候选 & 65.7\% & 原始PDF无目标价 & 可进入扩展观察候选池；缺少当前可正权使用的外部估值锚 \\
601061 & 中信金属 & 优先验证候选 & 60.5\% & 陈旧/低权重外部锚 & 进入优先验证候选池；缺少当前可正权使用的外部估值锚；最新研报缺失或早于2026-04-01 \\
002414 & 高德红外 & 扩展事件验证候选 & 56.4\% & 陈旧/低权重外部锚 & 可进入扩展观察候选池；缺少当前可正权使用的外部估值锚；经营现金流仍需验证 \\
601360 & 三六零 & 优先验证候选 & 45.6\% & 公开共识锚 & 进入优先验证候选池；缺少当前可正权使用的外部估值锚；最新研报缺失或早于2026-04-01；经营现金流仍需验证 \\
000833 & 粤桂股份 & 扩展事件验证候选 & 34.7\% & 原始PDF无目标价 & 可进入扩展观察候选池；缺少当前可正权使用的外部估值锚；经营现金流仍需验证 \\
002048 & 宁波华翔 & 优先验证候选 & 33.0\% & 无正权外部锚 & 进入优先验证候选池；缺少当前可正权使用的外部估值锚；最新研报缺失或早于2026-04-01；原始研报已归档但未披露目标价/合理价值区间 \\
002192 & 融捷股份 & 扩展事件验证候选 & 26.2\% & 原始PDF无目标价 & 可进入扩展观察候选池；缺少当前可正权使用的外部估值锚；经营现金流仍需验证 \\
600884 & 杉杉股份 & 扩展事件验证候选 & 21.5\% & 原始PDF无目标价 & 可进入扩展观察候选池；缺少当前可正权使用的外部估值锚 \\
000567 & 海德股份 & 优先验证候选 & 20.8\% & 原始PDF无目标价 & 进入优先验证候选池；缺少当前可正权使用的外部估值锚；最新研报缺失或早于2026-04-01；原始研报已归档但未披露目标价/合理价值区间 \\
600362 & 江西铜业 & 扩展事件验证候选 & 20.1\% & 陈旧/低权重外部锚 & 可进入扩展观察候选池；缺少当前可正权使用的外部估值锚 \\
000792 & 盐湖股份 & 扩展事件验证候选 & 16.5\% & 原始PDF无目标价 & 可进入扩展观察候选池；缺少当前可正权使用的外部估值锚 \\
600026 & 中远海能 & 安全边际观察 & 10.2\% & 陈旧/低权重外部锚 & 进入观察候选池；缺少当前可正权使用的外部估值锚；最新研报缺失或早于2026-04-01 \\
002460 & 赣锋锂业 & 安全边际观察 & 4.2\% & 无正权外部锚 & 进入观察候选池；缺少当前可正权使用的外部估值锚；最新研报缺失或早于2026-04-01；原始研报已归档但未披露目标价/合理价值区间 \\
000987 & 越秀资本 & 安全边际观察 & 1.4\% & 陈旧/低权重外部锚 & 进入观察候选池；缺少当前可正权使用的外部估值锚；最新研报缺失或早于2026-04-01 \\
000415 & 渤海租赁 & 暂不纳入 & -1.1\% & 无正权外部锚 & 暂不进入新增候选；缺少当前可正权使用的外部估值锚；概率目标低于现价 \\
002842 & 翔鹭钨业 & 暂不纳入 & -1.3\% & 原始PDF无目标价 & 暂不进入新增候选；缺少当前可正权使用的外部估值锚；经营现金流仍需验证；概率目标低于现价 \\
002468 & 申通快递 & 暂不纳入 & -2.2\% & 陈旧/低权重外部锚 & 暂不进入新增候选；缺少当前可正权使用的外部估值锚；概率目标低于现价 \\
000783 & 长江证券 & 暂不纳入 & -3.8\% & 原始PDF无目标价 & 暂不进入新增候选；缺少当前可正权使用的外部估值锚；概率目标低于现价 \\
600918 & 中泰证券 & 估值风险观察 & -6.3\% & 原始PDF无目标价 & 保留监控，不新增为候选；缺少当前可正权使用的外部估值锚；最新研报缺失或早于2026-04-01；原始研报已归档但未披露目标价/合理价值区间；概率目标低于当前价，仅保留为下行风险对照；概率目标低于现价 \\
600688 & 上海石化 & 估值风险观察 & -8.1\% & 无正权外部锚 & 保留监控，不新增为候选；缺少当前可正权使用的外部估值锚；最新研报缺失或早于2026-04-01；概率目标低于当前价，仅保留为下行风险对照；经营现金流仍需验证；概率目标低于现价 \\
002440 & 闰土股份 & 暂不纳入 & -12.1\% & 原始PDF无目标价 & 暂不进入新增候选；缺少当前可正权使用的外部估值锚；概率目标低于现价 \\
002458 & 益生股份 & 暂不纳入 & -15.0\% & 原始PDF无目标价 & 暂不进入新增候选；缺少当前可正权使用的外部估值锚；概率目标低于现价 \\
600037 & 歌华有线 & 估值风险观察 & -17.3\% & 无正权外部锚 & 保留监控，不新增为候选；缺少当前可正权使用的外部估值锚；最新研报缺失或早于2026-04-01；原始研报已归档但未披露目标价/合理价值区间；概率目标低于当前价，仅保留为下行风险对照；概率目标低于现价 \\
600522 & 中天科技 & 暂不纳入 & -18.3\% & 原始PDF无目标价 & 暂不进入新增候选；缺少当前可正权使用的外部估值锚；经营现金流仍需验证；概率目标低于现价 \\
603629 & 利通电子 & 暂不纳入 & -24.4\% & 原始PDF无目标价 & 暂不进入新增候选；缺少当前可正权使用的外部估值锚；概率目标低于现价 \\
002475 & 立讯精密 & 暂不纳入 & -29.4\% & 陈旧/低权重外部锚 & 暂不进入新增候选；缺少当前可正权使用的外部估值锚；经营现金流仍需验证；概率目标低于现价 \\
600906 & 财达证券 & 暂不纳入 & -40.6\% & 原始PDF无目标价 & 暂不进入新增候选；缺少当前可正权使用的外部估值锚；概率目标低于现价 \\
002636 & 金安国纪 & 暂不纳入 & -42.5\% & 原始PDF无目标价 & 暂不进入新增候选；缺少当前可正权使用的外部估值锚；经营现金流仍需验证；概率目标低于现价 \\
600601 & 方正科技 & 暂不纳入 & -43.6\% & 原始PDF无目标价 & 暂不进入新增候选；缺少当前可正权使用的外部估值锚；概率目标低于现价 \\
601069 & 西部黄金 & 估值风险观察 & -44.5\% & 原始PDF无目标价 & 保留监控，不新增为候选；缺少当前可正权使用的外部估值锚；最新研报缺失或早于2026-04-01；原始研报已归档但未披露目标价/合理价值区间；概率目标低于当前价，仅保留为下行风险对照；概率目标低于现价 \\
600095 & 湘财股份 & 估值风险观察 & -46.2\% & 陈旧/低权重外部锚 & 保留监控，不新增为候选；缺少当前可正权使用的外部估值锚；最新研报缺失或早于2026-04-01；概率目标低于当前价，仅保留为下行风险对照；概率目标低于现价 \\
600292 & 电投水电 & 估值风险观察 & -47.5\% & 原始PDF无目标价 & 保留监控，不新增为候选；缺少当前可正权使用的外部估值锚；最新研报缺失或早于2026-04-01；原始研报已归档但未披露目标价/合理价值区间；概率目标低于当前价，仅保留为下行风险对照；概率目标低于现价 \\
688183 & 生益电子 & 暂不纳入 & -47.8\% & 原始PDF无目标价 & 暂不进入新增候选；缺少当前可正权使用的外部估值锚；概率目标低于现价 \\
601088 & 中国神华 & 暂不纳入 & -50.4\% & 原始PDF无目标价 & 暂不进入新增候选；缺少当前可正权使用的外部估值锚；概率目标低于现价 \\
002378 & 章源钨业 & 暂不纳入 & -52.4\% & 原始PDF无目标价 & 暂不进入新增候选；缺少当前可正权使用的外部估值锚；经营现金流仍需验证；概率目标低于现价 \\
002916 & 深南电路 & 暂不纳入 & -52.9\% & 陈旧/低权重外部锚 & 暂不进入新增候选；缺少当前可正权使用的外部估值锚；概率目标低于现价 \\
600392 & 盛和资源 & 暂不纳入 & -55.6\% & 原始PDF无目标价 & 暂不进入新增候选；缺少当前可正权使用的外部估值锚；经营现金流仍需验证；概率目标低于现价 \\
601869 & 长飞光纤 & 暂不纳入 & -65.3\% & 原始PDF无目标价 & 暂不进入新增候选；缺少当前可正权使用的外部估值锚；概率目标低于现价 \\
002407 & 多氟多 & 暂不纳入 & -69.3\% & 原始PDF无目标价 & 暂不进入新增候选；缺少当前可正权使用的外部估值锚；经营现金流仍需验证；概率目标低于现价 \\
002611 & 东方精工 & 暂不纳入 & -81.2\% & 原始PDF无目标价 & 暂不进入新增候选；缺少当前可正权使用的外部估值锚；经营现金流仍需验证；概率目标低于现价 \\
002221 & 东华能源 & 暂不纳入 & -81.4\% & 原始PDF无目标价 & 暂不进入新增候选；缺少当前可正权使用的外部估值锚；经营现金流仍需验证；概率目标低于现价 \\
002174 & 游族网络 & 暂不纳入 & -95.4\% & 原始PDF无目标价 & 暂不进入新增候选；缺少当前可正权使用的外部估值锚；经营现金流仍需验证；概率目标低于现价 \\
\bottomrule
\end{longtable}
\normalsize

\begin{sourcequalitybox}[准入结论]
原中等证据标的中，没有一只直接升级为正式模型；可进入优先验证或扩展观察的标的必须继续等待当前外部锚、H2扣非利润和现金流验证。估值风险观察和概率目标低于现价的标的不新增为候选。
\end{sourcequalitybox}

\section{100\%以上空间样本的全量拦截审计}

下表覆盖全142只候选中模型空间超过100\%的全部样本。原始API、原始PDF、目标价字段和Q1现金流均已逐票核验；表中不再保留“待补目标价”占位，而是公开核验结果、准入结论和仍需等待的未来事件。

\footnotesize
\begin{longtable}{L{1.0cm}L{1.45cm}R{1.0cm}L{4.8cm}L{3.3cm}L{3.3cm}}
\toprule
\textbf{代码} & \textbf{标的} & \textbf{空间} & \textbf{来源覆盖/外部锚} & \textbf{最终准入} & \textbf{剩余事件验证} \\
\midrule
\endfirsthead
\toprule
\textbf{代码} & \textbf{标的} & \textbf{空间} & \textbf{来源覆盖/外部锚} & \textbf{最终准入} & \textbf{剩余事件验证} \\
\midrule
\endhead
002432 & 九安医疗 & 323.4\% & API/PDF 2/2；当前原始目标0；2024-12-23 70.28元，权重0\%；Q1现金流-0.92亿元 & 保留候选观察，维持未进入优先池及非正式模型 & 重新满足优先池的价格位置与行业阶段规则；H2扣非利润持续性与经营现金流转正；2026-04-01之后原始券商目标价或合理价值区间 \\
000623 & 吉林敖东 & 317.0\% & API/PDF 2/2；当前原始目标0；2024-03-27 18.02元，权重0\%；Q1现金流0.11亿元 & 保留优先池验证，维持非正式模型 & 2026-04-01之后原始券商目标价或合理价值区间；H2扣非利润持续性与投资收益质量；经营现金流持续为正 \\
600739 & 辽宁成大 & 300.0\% & API/PDF 4/4；当前原始目标0；2017-10-30 22.79元，权重0\%；Q1现金流-1.76亿元 & 保留优先池验证，维持非正式模型 & 2026-04-01之后原始券商目标价或合理价值区间；H2扣非利润持续性与广发证券投资收益质量；经营现金流转正 \\
000685 & 中山公用 & 162.8\% & API/PDF 6/6；当前原始目标0；2025-03-04 10.50--11.11元，权重0\%；Q1现金流-0.27亿元 & 保留候选观察，维持未进入优先池及非正式模型 & 重新满足优先池的价格位置与行业阶段规则；H2扣非利润与投资收益可持续性；经营现金流转正；当前原始券商目标价或合理价值区间 \\
600150 & 中国船舶 & 142.9\% & API/PDF 56/8；当前原始目标0；2023-03-27 31.36元，权重0\%；Q1现金流81.52亿元 & 保留优先池验证，维持非正式模型 & 当前原始券商目标价或合理价值区间；H2订单交付、船价、利润率与应收回款；集团整合后的持续盈利与现金流 \\
301308 & 江波龙 & 119.2\% & API/PDF 12/8；当前原始目标0；原始语料未找到可接受目标价；Q1现金流-28.75亿元 & 保留候选观察，维持未进入优先池及非正式模型 & 重新满足优先池的价格位置与行业阶段规则；当前原始券商目标价或合理价值区间；H2存储价格、出货、毛利率与库存周期；经营现金流显著改善 \\
\bottomrule
\end{longtable}
\normalsize

\end{document}
